% Les médias et leur impact sur l’opinion publique — ECM 1re Tech — Cours signé OnBuch+
% Programme officiel MINESEC. Compiler avec : tectonic N03-medias-leur-impact-opinion-publique.tex
\newcommand{\DOCMATIERE}{ECM}
\newcommand{\DOCNIVEAU}{1re Tech}
\newcommand{\DOCMODULE}{ECM – classe de Première technique}
\newcommand{\DOCLECON}{N03}
\newcommand{\DOCTITRE}{Les médias et leur impact sur l’opinion publique}
% OnBuch+ — préambule « cours signés » ENRICHI (compilé avec tectonic / XeLaTeX)
% Système visuel « Pop » d'OnBuch+ : fond crème (jamais blanc), bordures encre
% épaisses, ombres « dures » décalées, titres Archivo Black en capitales.
% Version MATHÉMATIQUES : graphes (pgfplots/tikz), boîtes pédagogiques
% (définition, propriété, méthode, attention, le savais-tu, exercice résolu,
% exercice, TP, bilan), page de garde et signature OnBuch+ ancrée en bas.
%
% Macros attendues dans main.tex AVANT \input{preamble} :
%   \DOCMATIERE  \DOCNIVEAU  \DOCMODULE  \DOCLECON (numéro)  \DOCTITRE
\documentclass[11pt]{article}

\usepackage{fontspec}
\usepackage[french]{babel}
\usepackage[a4paper,margin=2cm,top=3.4cm,bottom=2.8cm,headheight=1.4cm,footskip=1.3cm]{geometry}
\usepackage{amsmath,amssymb,amsfonts}
\usepackage{mathtools} % \xmapsto, \coloneqq... souvent utilisés par les modèles
\usepackage{cancel} % simplifications barrées (bilans de forces)
% \abs{z} (module, valeur absolue) et \norm{u} (norme), utilisés par les modèles
\providecommand{\abs}{}\let\abs\relax\DeclarePairedDelimiter{\abs}{\lvert}{\rvert}
\providecommand{\norm}{}\let\norm\relax\DeclarePairedDelimiter{\norm}{\lVert}{\rVert}
\usepackage[table]{xcolor} % [table] : fournit aussi \rowcolors (alternance de lignes)
\usepackage{pagecolor}
\usepackage{tikz}
\usetikzlibrary{arrows.meta,positioning,calc,shapes.geometric,shapes.misc,shapes.arrows,decorations.pathmorphing,decorations.pathreplacing,decorations.markings,patterns,shadows,mindmap,trees,fit,backgrounds,matrix,intersections,angles,quotes}
\usepackage{pgfplots}
\usepackage[version=4]{mhchem} % équations nucléaires \ce{^{235}_{92}U}
\usepackage[european,siunitx]{circuitikz} % schémas électriques
\pgfplotsset{compat=1.18}
\usepgfplotslibrary{fillbetween,groupplots} % aires entre courbes (calcul intégral)
\usepackage{enumitem}
\usepackage{fancyhdr}
% [most] charge toutes les bibliothèques tcolorbox (listings, minted, raster,
% documentation...) et gonfle inutilement la mémoire TeX sur un document long
% (~300 boîtes) : on ne charge que ce qui est réellement utilisé.
\usepackage[skins,breakable]{tcolorbox}
\usepackage{graphicx}
\usepackage{colortbl}
\usepackage{booktabs}
\usepackage{tabularx}
\usepackage{multirow}
\usepackage{diagbox} % case d'en-tête diagonale des tableaux à double entrée
% Colonnes extensibles centrée / gauche / droite (C, L, R), souvent utilisées par les modèles
\newcolumntype{C}{>{\centering\arraybackslash}X}
\newcolumntype{L}{>{\raggedright\arraybackslash}X}
\newcolumntype{R}{>{\raggedleft\arraybackslash}X}
\usepackage{array}
\usepackage{multicol}
\usepackage{siunitx}
% parse-numbers=false : accepte aussi \SI{2,5 \times 10^{-3}}{...}
\sisetup{locale=FR,output-decimal-marker={,},group-minimum-digits=4,parse-numbers=false}
\DeclareSIUnit\ohm{\ensuremath{\symup{\Omega}}} % U+2126 absent de la police
\DeclareSIUnit\division{div} % réglages d'oscilloscope (ms/div, V/div)
\DeclareSIUnit\tr{tr} % tours (vitesse de rotation, ex. \SI{1500}{\tr\per\minute})
\DeclareSIUnit\molar{M} % concentration molaire (ex. \SI{3}{\molar}, \SI{1}{\micro\molar})
\DeclareSIUnit\an{an} % année (le modèle confond parfois avec \year, un compteur TeX) — ex. \SI{5}{\kilo\meter\per\an}
\DeclareSIUnit\FCFA{FCFA} % franc CFA (ex. \SI{500}{\FCFA\per\kilo\gram})
\DeclareSIUnit\degreeBrix{\textdegree Brix} % teneur en sucre d'un sirop/jus (réfractométrie)
\DeclareSIUnit\month{mois} % le modèle utilise parfois \month, un compteur TeX, comme unité de durée
\DeclareSIUnit\microliter{\micro\liter} % le modèle écrit parfois \microliter en un seul token
\DeclareSIUnit\million{millions} % dénombrement (ex. \SI{15}{\million\per\mL} spermatozoïdes)
\usepackage{qrcode}
% \degree : commande fréquemment utilisée par les modèles pour un angle ou une
% température en dehors de \SI{}{} (non fournie par les paquets chargés).
\providecommand{\degree}{^{\circ}}
% \qed (fin de démonstration), utilisé par les modèles sans amsthm
\providecommand{\qed}{\hfill$\square$}
% \barre : double barre des valeurs interdites dans un tableau de variations (tkz-tab)
\providecommand{\barre}{\ensuremath{\|}}
% environnement proof (amsthm non chargé) utilisé par les modèles
\ifdefined\proof\else\newenvironment{proof}[1][Démonstration]{\par\noindent\textit{#1.}\ }{\qed\par}\fi
\usepackage{lastpage}
\usepackage{hyperref}
\hypersetup{hidelinks}

% ============================================================
% Palette « Pop » OnBuch+ — mêmes hex que la classe P (plus_ui.dart)
% ============================================================
\definecolor{poporange}{HTML}{F59321}
\definecolor{popdark}{HTML}{E07A0C}
\definecolor{popcream}{HTML}{FFF6E8}
\definecolor{popink}{HTML}{1C1714}
\definecolor{poppurple}{HTML}{7A5AE0}
\definecolor{popgreen}{HTML}{2E9E6A}
\definecolor{poppink}{HTML}{E0457B}
\definecolor{popblue}{HTML}{2F7DE1}
\definecolor{popgold}{HTML}{C98A12}
\definecolor{popmuted}{HTML}{8A7F76}
\definecolor{popline}{HTML}{EADFD2}
\definecolor{popgray}{HTML}{8A7F76}
\colorlet{popgrey}{popgray}
% Alias tolérants (le modèle nomme parfois les couleurs différemment)
\colorlet{popwhite}{white}
\colorlet{popblack}{popink}
\colorlet{popred}{poppink}
\colorlet{popyellow}{popgold}
% Teintes claires (fonds de tableaux/encadrés) — le modèle nomme parfois
% les couleurs avec un suffixe L (ex. popblueL) sans les avoir définies.
\colorlet{poporangeL}{poporange!20}
\colorlet{popdarkL}{popdark!20}
\colorlet{poppurpleL}{poppurple!20}
\colorlet{popgreenL}{popgreen!20}
\colorlet{poppinkL}{poppink!20}
\colorlet{popblueL}{popblue!20}
\colorlet{popgoldL}{popgold!20}
\colorlet{popredL}{popred!20}
\colorlet{popgrayL}{popgray!20}
% Teintes claires pour fonds de tableaux / zones
\colorlet{popblueL}{popblue!10!white}
\colorlet{poporangeL}{poporange!14!white}
\colorlet{popgreenL}{popgreen!12!white}
\colorlet{poppurpleL}{poppurple!12!white}
\colorlet{poppinkL}{poppink!12!white}
\colorlet{popgoldL}{popgold!14!white}

\pagecolor{popcream}

% ============================================================
% Typographie
% ============================================================
\defaultfontfeatures{Ligatures=TeX}
\setmainfont{PlusJakartaSans-Medium.ttf}[Path=../../../fonts/,BoldFont=PlusJakartaSans-Bold.ttf]
% Maths en sans-serif (Fira Math) avec chiffres et lettres droites de Plus
% Jakarta, pour que les formules se fondent dans le texte.
\usepackage[math-style=ISO,bold-style=ISO]{unicode-math}
\setmathfont{FiraMath-Regular.otf}[Path=../../../fonts/]
\setmathfont{PlusJakartaSans-Medium.ttf}[Path=../../../fonts/,range={up/num,up/{Latin,latin},\mathup}]
\usepackage{icomma} % virgule décimale sans espace en mode math
% Caractères mathématiques Unicode absents des polices de texte (Plus Jakarta,
% Archivo Black) : rendus par leur équivalent mathématique, en texte comme en
% formule — sinon ils s'affichent en carré vide (ex. « Arithmétique dans ℤ »).
\usepackage{newunicodechar}
\newunicodechar{ℝ}{\ensuremath{\mathbb{R}}}
\newunicodechar{ℤ}{\ensuremath{\mathbb{Z}}}
\newunicodechar{ℂ}{\ensuremath{\mathbb{C}}}
\newunicodechar{ℕ}{\ensuremath{\mathbb{N}}}
\newunicodechar{ℚ}{\ensuremath{\mathbb{Q}}}
\newunicodechar{𝔻}{\ensuremath{\mathbb{D}}}
\newunicodechar{α}{\ensuremath{\alpha}}
\newunicodechar{β}{\ensuremath{\beta}}
\newunicodechar{γ}{\ensuremath{\gamma}}
\newunicodechar{δ}{\ensuremath{\delta}}
\newunicodechar{ε}{\ensuremath{\varepsilon}}
\newunicodechar{θ}{\ensuremath{\theta}}
\newunicodechar{λ}{\ensuremath{\lambda}}
\newunicodechar{μ}{\ensuremath{\mu}}
\newunicodechar{σ}{\ensuremath{\sigma}}
\newunicodechar{τ}{\ensuremath{\tau}}
\newunicodechar{φ}{\ensuremath{\varphi}}
\newunicodechar{ψ}{\ensuremath{\psi}}
\newunicodechar{ω}{\ensuremath{\omega}}
\newunicodechar{Σ}{\ensuremath{\Sigma}}
% majuscules produites par \MakeUppercase dans les titres (α-aminés -> Α)
\newunicodechar{Α}{\ensuremath{\alpha}}
\newunicodechar{Β}{\ensuremath{\beta}}
\newunicodechar{Γ}{\ensuremath{\Gamma}}
\newunicodechar{Δ}{\ensuremath{\Delta}}
\newunicodechar{Ε}{\ensuremath{\varepsilon}}
\newunicodechar{Θ}{\ensuremath{\Theta}}
\newunicodechar{Λ}{\ensuremath{\Lambda}}
\newunicodechar{Μ}{\ensuremath{\mu}}
\newunicodechar{Τ}{\ensuremath{\tau}}
\newunicodechar{Φ}{\ensuremath{\Phi}}
\newunicodechar{Ψ}{\ensuremath{\Psi}}
\newunicodechar{Ω}{\ensuremath{\Omega}}
\newunicodechar{↦}{\ensuremath{\mapsto}}
\newunicodechar{∈}{\ensuremath{\in}}
\newunicodechar{∉}{\ensuremath{\notin}}
\newunicodechar{∧}{\ensuremath{\wedge}}
\newunicodechar{⇔}{\ensuremath{\Leftrightarrow}}
\newunicodechar{⟺}{\ensuremath{\Longleftrightarrow}}
\newunicodechar{⇒}{\ensuremath{\Rightarrow}}
\newunicodechar{⟹}{\ensuremath{\Longrightarrow}}
\newunicodechar{∘}{\ensuremath{\circ}}
\newunicodechar{∪}{\ensuremath{\cup}}
\newunicodechar{∩}{\ensuremath{\cap}}
\newunicodechar{≡}{\ensuremath{\equiv}}
\newunicodechar{⊂}{\ensuremath{\subset}}
\newunicodechar{⊥}{\ensuremath{\perp}}
\newunicodechar{∃}{\ensuremath{\exists}}
\newunicodechar{∀}{\ensuremath{\forall}}
\newunicodechar{∛}{\ensuremath{\sqrt[3]{\,}}}
\newunicodechar{∜}{\ensuremath{\sqrt[4]{\,}}}
\newunicodechar{⃗}{\ensuremath{{}^{\to}}}
\newunicodechar{ⁿ}{\ifmmode^{n}\else\textsuperscript{\ensuremath{n}}\fi}
\newunicodechar{ˣ}{\ifmmode^{x}\else\textsuperscript{\ensuremath{x}}\fi}
\newunicodechar{ᵇ}{\ifmmode^{b}\else\textsuperscript{\ensuremath{b}}\fi}
\newunicodechar{ʳ}{\ifmmode^{r}\else\textsuperscript{\ensuremath{r}}\fi}
\newunicodechar{ᵘ}{\ifmmode^{u}\else\textsuperscript{\ensuremath{u}}\fi}
\newunicodechar{ʸ}{\ifmmode^{y}\else\textsuperscript{\ensuremath{y}}\fi}
\newunicodechar{ᵃ}{\ifmmode^{a}\else\textsuperscript{\ensuremath{a}}\fi}
\newunicodechar{ᵉ}{\ifmmode^{e}\else\textsuperscript{\ensuremath{e}}\fi}
\newunicodechar{ᵏ}{\ifmmode^{k}\else\textsuperscript{\ensuremath{k}}\fi}
\newunicodechar{ᵖ}{\ifmmode^{p}\else\textsuperscript{\ensuremath{p}}\fi}
\newunicodechar{ᴺ}{\ifmmode^{N}\else\textsuperscript{\ensuremath{N}}\fi}
\newunicodechar{ᶜ}{\ifmmode^{c}\else\textsuperscript{\ensuremath{c}}\fi}
\newunicodechar{ʰ}{\ifmmode^{h}\else\textsuperscript{\ensuremath{h}}\fi}
\newunicodechar{ᵠ}{\ifmmode^{q}\else\textsuperscript{\ensuremath{q}}\fi}
\newunicodechar{ₙ}{\ifmmode_{n}\else\textsubscript{\ensuremath{n}}\fi}
\newunicodechar{ₐ}{\ifmmode_{a}\else\textsubscript{\ensuremath{a}}\fi}
\newunicodechar{ᵢ}{\ifmmode_{i}\else\textsubscript{\ensuremath{i}}\fi}
\newunicodechar{ᵣ}{\ifmmode_{r}\else\textsubscript{\ensuremath{r}}\fi}
\newunicodechar{ₖ}{\ifmmode_{k}\else\textsubscript{\ensuremath{k}}\fi}
\newunicodechar{ᵦ}{\ifmmode_{\beta}\else\textsubscript{\ensuremath{\beta}}\fi}
\newunicodechar{ₓ}{\ifmmode_{x}\else\textsubscript{\ensuremath{x}}\fi}
\newfontfamily\popdisplay{ArchivoBlack-Regular.ttf}[Path=../../../fonts/]
\newfontfamily\poplogo{SpaceGrotesk-Bold.ttf}[Path=../../../fonts/]
\newfontfamily\popsemi{PlusJakartaSans-SemiBold.ttf}[Path=../../../fonts/]
\newfontfamily\popxbold{PlusJakartaSans-ExtraBold.ttf}[Path=../../../fonts/]
\newfontfamily\popxboldls{PlusJakartaSans-ExtraBold.ttf}[Path=../../../fonts/,LetterSpace=3.0]

\setlist[enumerate]{leftmargin=1.7em,itemsep=5pt,topsep=4pt}
\setlist[itemize]{leftmargin=1.7em,itemsep=4pt,topsep=4pt,label={\color{poporange}\textbullet}}
\setlength{\parindent}{0pt}
\setlength{\parskip}{5pt}
\renewcommand{\arraystretch}{1.25}

% pgfplots : style Pop par défaut
\pgfplotsset{
  every axis/.append style={axis line style={popink,line width=1pt},tick style={popink},
    grid style={popline,line width=0.5pt},label style={font=\small},tick label style={font=\footnotesize}},
  popcurve/.style={poporange,line width=1.8pt,no markers,smooth},
}
% Même style disponible dans un \draw TikZ simple (hors pgfplots)
\tikzset{popcurve/.style={poporange,line width=1.8pt}}
\tikzset{popgrid/.style={popline,very thin}}
\colorlet{popcurve}{poporange} % le nom du style est parfois utilisé comme couleur

% ============================================================
% Pill / squiggle
% ============================================================
\newcommand{\poppill}[3]{%
  \tikz[baseline=-0.55ex]\node[draw=popink,line width=1.2pt,rounded corners=2.6pt,%
    fill=#2,inner xsep=6.5pt,inner ysep=3pt]{%
    \popxboldls\fontsize{7}{7}\selectfont\color{#3}\MakeUppercase{#1}};%
}
\newcommand{\popsquiggle}[1]{%
  \tikz[baseline=-0.4ex]{
    \draw[#1,line width=2.6pt,line cap=round]
      plot [smooth] coordinates {(0,0.09) (0.34,0.01) (0.68,0.21) (1.02,0.01) (1.36,0.10)};
  }%
}
% Mot-clé surligné (marqueur orange sous le texte)
\newcommand{\cle}[1]{{\popxbold\color{popdark}#1}}

% ============================================================
% En-tête / pied
% ============================================================
\pagestyle{fancy}
\fancyhf{}
\renewcommand{\headrulewidth}{0pt}
\renewcommand{\footrulewidth}{0pt}
\lhead{\raisebox{-0.15ex}{\poplogo\fontsize{13}{13}\selectfont\color{popink}OnBuch\color{poporange}+}}
\rhead{\poppill{\DOCMATIERE\ \textbullet\ \DOCNIVEAU\ \textbullet\ Leçon \DOCLECON}{white}{popink}}
\lfoot{\popxbold\fontsize{7.6}{7.6}\selectfont\color{popmuted}\MakeUppercase{Cours signé OnBuch+}\ \textbullet\ {\popsemi\DOCTITRE}}
\rfoot{\tikz[baseline=-0.6ex]\node[rounded corners=8.5pt,fill=popink,minimum height=17pt,minimum width=17pt,inner xsep=5pt,inner sep=0]{\popxbold\fontsize{8}{8}\selectfont\color{popcream}\thepage\,/\,\pageref*{LastPage}};}

% ============================================================
% Titres
% ============================================================
\newcommand{\chaptitre}[2]{%
  \begin{tcolorbox}[enhanced,colback=poporange,colframe=popink,boxrule=2.6pt,arc=11pt,
      drop shadow=popink,left=15pt,right=15pt,top=12pt,bottom=11pt,before skip=3pt,after skip=18pt]
    \poppill{#2}{popink}{popcream}\par\vspace{9pt}
    {\raggedright\hyphenpenalty=10000\exhyphenpenalty=10000\popdisplay\fontsize{22}{24}\selectfont\color{popcream}\MakeUppercase{#1}\par}
  \end{tcolorbox}
}
% Section numérotée : pastille orange avec le numéro + titre Archivo
\newcounter{popsec}
\newcommand{\coursec}[1]{%
  \stepcounter{popsec}\setcounter{popsub}{0}%
  \par\vspace{16pt}\noindent
  \begin{minipage}{\linewidth}
  \tikz[baseline=-0.7ex]\node[draw=popink,line width=1.6pt,fill=poporange,rounded corners=4pt,minimum size=22pt,inner sep=0]{\popdisplay\fontsize{12}{12}\selectfont\color{popcream}\thepopsec};%
  \hspace{8pt}{\popdisplay\fontsize{15}{17}\selectfont\color{popink}\MakeUppercase{#1}}\par
  \vspace{4pt}\noindent{\color{poporange}\rule{36pt}{3.6pt}}
  \end{minipage}\par\nopagebreak\vspace{8pt}
}
\newcounter{popsub}
\newcommand{\courssub}[1]{%
  \stepcounter{popsub}%
  \par\vspace{9pt}\noindent{\popxbold\fontsize{11.5}{13}\selectfont\color{popink}{\color{poporange}\thepopsec.\thepopsub}\hspace{6pt}#1}\par\nopagebreak\vspace{4pt}
}

% ============================================================
% Boîtes pédagogiques — même recette : fond blanc, bordure épaisse colorée,
% ombre dure encre, badge plein. Titre optionnel [#1].
% ============================================================
\tcbset{popbase/.style={breakable,enhanced,colback=white,boxrule=2.3pt,arc=9pt,
  shadow={3pt}{-3pt}{0pt}{popink},left=14pt,right=14pt,top=11pt,bottom=11pt,before skip=11pt,after skip=15pt,
  fontupper=\popsemi\fontsize{10.6}{14.5}\selectfont\color{popink},
  fontlower=\popsemi\fontsize{10.6}{14.5}\selectfont\color{popink}}}
% Coche dessinée (le glyphe ✓ n'existe pas dans les polices du document)
\newcommand{\popcheck}[1]{\tikz[baseline=-0.5ex]\draw[#1,line width=1.6pt,line cap=round,line join=round](-0.13,0.0)--(-0.04,-0.09)--(0.13,0.1);}
\providecommand{\coche}{\popcheck{popgreen}}
\providecommand{\resultat}[1]{\par\smallskip\noindent\fcolorbox{popgreen}{popgreenL}{\parbox{\dimexpr\linewidth-2\fboxsep-2\fboxrule\relax}{\textbf{Résultat.} #1}}\par\smallskip}
\newcommand{\popboxhead}[3]{\poppill{#1}{#2}{white}\ifx\relax#3\relax\else\hspace{6pt}{\popxbold\fontsize{10.6}{12}\selectfont\color{popink}#3}\fi\par\vspace{7pt}}

\newtcolorbox{aretenir}[1][]{popbase,colframe=popgreen,before upper={\popboxhead{À retenir}{popgreen}{#1}}}
\newtcolorbox{exemplebox}[1][]{popbase,colframe=poporange,before upper={\popboxhead{Exemple}{poporange}{#1}}}
\newtcolorbox{definition}[1][]{popbase,colframe=popblue,before upper={\popboxhead{Définition}{popblue}{#1}}}
\newtcolorbox{propriete}[1][]{popbase,colframe=poppurple,before upper={\popboxhead{Propriété}{poppurple}{#1}}}
\newtcolorbox{methode}[1][]{popbase,colframe=popgold,before upper={\popboxhead{Méthode}{popgold}{#1}}}
\newtcolorbox{attention}[1][]{popbase,colframe=poppink,before upper={\popboxhead{Attention}{poppink}{#1}}}
\newtcolorbox{experience}[1][]{popbase,colframe=popblue,colback=popblueL,before upper={\popboxhead{Expérience}{popblue}{#1}}}
% « Le savais-tu ? » : carte violette pleine (contexte, culture, Cameroun)
\newtcolorbox{savaistu}[1][]{popbase,colback=poppurple,colframe=popink,
  fontupper=\popsemi\fontsize{10.4}{14.2}\selectfont\color{white},
  before upper={\poppill{Le savais-tu\,?}{white}{poppurple}\ifx\relax#1\relax\else\hspace{6pt}{\popxbold\color{white}#1}\fi\par\vspace{7pt}}}
% Exercice résolu : énoncé en haut, \tcblower puis la solution sur fond crème
\newcounter{popexo}\newcounter{popexr}
\newtcolorbox{exoresolu}[1][]{popbase,colframe=poporange,colbacklower=popcream,
  segmentation style={popink,line width=1.2pt,dashed},
  before upper={\stepcounter{popexr}\popboxhead{Exercice résolu \thepopexr}{poporange}{#1}},
  before lower={\poppill{Solution}{popink}{popcream}\par\vspace{6pt}}}
% Exercice (non résolu en place) — la correction est dans la partie Corrigés
\newtcolorbox{exercice}[1][]{popbase,colframe=popink,
  before upper={\stepcounter{popexo}\popboxhead{Exercice \thepopexo}{popink}{#1}}}
% Corrigé
\newtcolorbox{corrige}[1][]{popbase,colframe=popgreen,colback=popgreenL,
  before upper={\popboxhead{Corrigé}{popgreen}{#1}}}
% Objectifs / prérequis / bilan
\newtcolorbox{objectifs}{popbase,colframe=popink,colback=poporangeL,
  before upper={\popboxhead{Objectifs de la leçon}{popink}{}}}
\newtcolorbox{prerequis}{popbase,colframe=popink,colback=popblueL,
  before upper={\popboxhead{Prérequis}{popblue}{}}}
\newtcolorbox{bilan}{popbase,colframe=popink,colback=white,boxrule=2.8pt,
  before upper={\popboxhead{Fiche bilan}{poporange}{L'essentiel à connaître pour l'examen}}}
% Figure encadrée avec légende
\newtcolorbox{popfigure}[1][]{enhanced,breakable=false,colback=white,colframe=popline,boxrule=1.4pt,arc=8pt,
  left=10pt,right=10pt,top=10pt,bottom=8pt,before skip=10pt,after skip=12pt,halign=center,
  lower separated=false,halign lower=center,
  fontlower=\popsemi\fontsize{9}{11.5}\selectfont\color{popmuted},
  #1}
\newcommand{\legende}[1]{\par\vspace{6pt}{\popsemi\fontsize{9}{11.5}\selectfont\color{popmuted}#1}}

% ============================================================
% Signature OnBuch+
% ============================================================
\newcommand{\poplogoinline}[1]{{\poplogo\fontsize{#1}{#1}\selectfont\color{popink}OnBuch\color{poporange}+}}
\newcommand{\signatureonbuch}{%
  \begin{tcolorbox}[enhanced,colback=popink,colframe=popink,boxrule=2.2pt,arc=10pt,
      drop shadow=poporange,left=14pt,right=14pt,top=10pt,bottom=10pt,before skip=0pt,after skip=0pt]
    \tikz[baseline=-0.7ex]\node[circle,fill=popgreen,draw=popcream,line width=1.3pt,minimum size=22pt,inner sep=0]{\popcheck{white}};%
    \hspace{9pt}\begin{minipage}[c]{0.62\linewidth}
      {\popxbold\fontsize{12}{13}\selectfont\color{popcream}Signé\ \textbullet\ {\poplogo OnBuch\color{poporange}+}}\par\vspace{2pt}
      {\raggedright\popsemi\fontsize{8}{10.5}\selectfont\color{popline}\MakeUppercase{Équipe pédagogique OnBuch+}\\\MakeUppercase{Conforme au programme officiel MINESEC}\par}
    \end{minipage}\hfill
    {\poplogo\fontsize{22}{22}\selectfont\color{popcream}OnBuch\color{poporange}+}
  \end{tcolorbox}%
}

% ============================================================
% Page de garde
% #1 titre · #2 matière · #3 niveau · #4 module · #5 n° leçon · #6 durée ·
% #7 description / objectifs (texte ou itemize)
% ============================================================
\newcommand{\pagegarde}[7]{%
  \thispagestyle{empty}
  \newgeometry{a4paper,margin=1.7cm,top=1.6cm,bottom=1.4cm}%
  \begin{tikzpicture}[remember picture,overlay]
    \fill[poporange!18] ([xshift=-3.2cm,yshift=-3.6cm]current page.north east) circle (4.6cm);
    \fill[poppurple!14] ([xshift=2.4cm,yshift=7.6cm]current page.south west) circle (3.8cm);
  \end{tikzpicture}%
  {\poplogo\fontsize{30}{30}\selectfont\color{popink}OnBuch\color{poporange}+}\hfill
  \raisebox{0.6ex}{\poppill{Cours signé \textbullet\ Premium}{popink}{popcream}}\par
  \vspace{1pt}\popsquiggle{poppurple}\par
  \vspace{8pt}
  \begin{tcolorbox}[enhanced,colback=poporange,colframe=popink,boxrule=2.8pt,arc=13pt,
      drop shadow=popink,left=18pt,right=18pt,top=12pt,bottom=13pt,after skip=10pt]
    \poppill{#2 \textbullet\ #3}{popink}{popcream}\hspace{5pt}\poppill{#4}{white}{popink}\par\vspace{8pt}
    {\popdisplay\fontsize{14}{14}\selectfont\color{popink}LEÇON #5}\par\vspace{4pt}
    {\raggedright\hyphenpenalty=10000\exhyphenpenalty=10000\popdisplay\fontsize{25}{27}\selectfont\color{popcream}\MakeUppercase{#1}\par}
    \vspace{8pt}
    {\popxbold\fontsize{9.5}{9.5}\selectfont\color{popink}\MakeUppercase{Durée conseillée : #6}}
  \end{tcolorbox}
  \begin{tcolorbox}[enhanced,colback=white,colframe=popink,boxrule=2.2pt,arc=10pt,
      drop shadow=popink,left=16pt,right=16pt,top=10pt,bottom=10pt,after skip=8pt]
    \poppill{Dans cette leçon}{poppurple}{white}\par\vspace{5pt}
    \popsemi\fontsize{9.3}{12.6}\selectfont\color{popink}#7
  \end{tcolorbox}
  \noindent\begin{minipage}[t]{0.487\linewidth}
    \begin{tcolorbox}[enhanced,colback=popgreen,colframe=popink,boxrule=2.2pt,arc=10pt,
        drop shadow=popink,left=11pt,right=11pt,top=10pt,bottom=10pt,height=3.1cm,valign=center]
      \tcbox[colback=white,colframe=popink,boxrule=1.3pt,arc=5pt,boxsep=0pt,nobeforeafter,
        left=3pt,right=3pt,top=3pt,bottom=3pt,valign=center]{\qrcode[height=1.9cm]{https://play.google.com/store/apps/details?id=com.luvvix.onbuch&hl=fr}}%
      \hspace{8pt}\begin{minipage}[c]{0.5\linewidth}
        {\popdisplay\fontsize{11}{12.5}\selectfont\color{white}\MakeUppercase{Télécharge OnBuch}}\par\vspace{4pt}
        {\raggedright\popsemi\fontsize{8.4}{11}\selectfont\color{white}Gratuit \textbullet\ Google Play\\Tuteur IA, annales, concours, résultats\par}
      \end{minipage}
    \end{tcolorbox}
  \end{minipage}\hfill
  \begin{minipage}[t]{0.487\linewidth}
    \begin{tcolorbox}[enhanced,colback=poppurple,colframe=popink,boxrule=2.2pt,arc=10pt,
        drop shadow=popink,left=11pt,right=11pt,top=10pt,bottom=10pt,height=3.1cm,valign=center]
      \tikz[baseline=-0.7ex]\node[circle,fill=white,draw=popink,line width=1.3pt,minimum size=1.6cm,inner sep=0]{\popdisplay\fontsize{24}{24}\selectfont\color{poppurple}+};%
      \hspace{8pt}\begin{minipage}[c]{0.6\linewidth}
        {\popdisplay\fontsize{11}{12.5}\selectfont\color{white}\MakeUppercase{Passe à OnBuch+}}\par\vspace{4pt}
        {\raggedright\popsemi\fontsize{8.4}{11}\selectfont\color{white}Tous les cours signés, Tuteur IA illimité, fiches PDF — onglet OnBuch+ de l'app\par}
      \end{minipage}
    \end{tcolorbox}
  \end{minipage}
  \vspace{8pt}
  \popcontenu
  \vfill
  \signatureonbuch
  \restoregeometry
  \newpage
}

% Rangée « Ce cours contient » (page de garde)
\newcommand{\poptile}[3]{% #1 couleur #2 gros texte #3 légende
  \begin{tcolorbox}[enhanced,colback=#1,colframe=popink,boxrule=2pt,arc=9pt,shadow={2.5pt}{-2.5pt}{0pt}{popink},
      left=6pt,right=6pt,top=9pt,bottom=9pt,halign=center,valign=center,height=1.75cm,width=\linewidth,nobeforeafter]
    {\popdisplay\fontsize{12.5}{13}\selectfont\color{white}\MakeUppercase{#2}}\par\vspace{3pt}
    {\centering\popsemi\fontsize{7.8}{9.5}\selectfont\color{white}#3\par}
  \end{tcolorbox}}
\newcommand{\popcontenu}{%
  \par\noindent{\popxboldls\fontsize{7.5}{7.5}\selectfont\color{popmuted}CE COURS SIGNÉ CONTIENT}\par\vspace{7pt}
  \noindent\begin{minipage}[t]{0.235\linewidth}\poptile{popblue}{Cours}{complet, illustré, pas à pas}\end{minipage}\hfill
  \begin{minipage}[t]{0.235\linewidth}\poptile{poporange}{Méthodes}{et exercices résolus}\end{minipage}\hfill
  \begin{minipage}[t]{0.235\linewidth}\poptile{poppink}{Exercices}{type examen + corrigés}\end{minipage}\hfill
  \begin{minipage}[t]{0.235\linewidth}\poptile{popgreen}{Fiche bilan}{l'essentiel à retenir}\end{minipage}\par}

% Sommaire
\newtcolorbox{sommaire}{enhanced,colback=white,colframe=popink,boxrule=2.2pt,arc=10pt,
  drop shadow=popink,left=18pt,right=18pt,top=13pt,bottom=13pt,before skip=6pt,after skip=20pt,
  before upper={\poppill{Sommaire}{poppurple}{white}\par\vspace{9pt}\popsemi\fontsize{10.6}{14.5}\selectfont\color{popink}}}

% Signature de fin de cours — poussée tout en bas de la dernière page
\newcommand{\findecours}{%
  \par\vfill\nopagebreak
  \begin{center}\popsquiggle{poporange}\end{center}\vspace{4pt}
  \signatureonbuch
}
\AtBeginDocument{\def\llbracket{\mathopen{[\mkern-3.5mu[}}\def\rrbracket{\mathclose{]\mkern-3.5mu]}}} % intervalles d'entiers (glyphe ⟦ absent de la police)
% Glyphes absents de Fira Math : redéfinis à partir de glyphes disponibles
\AtBeginDocument{%
  \def\setminus{\mathbin{\backslash}}\let\smallsetminus\setminus
  \def\vdots{\mathord{\vbox{\baselineskip3.2pt\lineskiplimit0pt\kern4pt\hbox{.}\hbox{.}\hbox{.}}}}%
  \def\ddots{\mathinner{\mkern1mu\raise6.5pt\hbox{.}\mkern2mu\raise3.6pt\hbox{.}\mkern2mu\raise0.7pt\hbox{.}\mkern1mu}}%
  \def\triangle{\mathord{\scalebox{0.9}{$\symup{\Delta}$}}}%
  \def\nabla{\mathord{\raisebox{\depth}{\scalebox{1}[-1]{$\symup{\Delta}$}}}}%
  \def\checkmark{\popcheck{popdark}}%
  \def\star{\mathbin{\ast}}\def\bigstar{\mathord{\ast}}%
  \def\diamond{\mathbin{\scalebox{0.6}{$\Diamond$}}}%
  \def\bigcup{\mathop{\mathchoice{\raisebox{-0.3em}{\scalebox{1.7}{$\cup$}}}{\scalebox{1.25}{$\cup$}}{\cup}{\cup}}\displaylimits}%
  \def\bigcap{\mathop{\mathchoice{\raisebox{-0.3em}{\scalebox{1.7}{$\cap$}}}{\scalebox{1.25}{$\cap$}}{\cap}{\cap}}\displaylimits}%
  \def\lesseqqgtr{\mathrel{\lessgtr}}\def\bigoplus{\mathop{\oplus}\displaylimits}%
}
\providecommand{\ssi}{si et seulement si} % abréviation fréquente
\AtBeginDocument{\providecommand{\wideparen}{\overparen}} % arc de cercle

%% FIGFIX-BEGIN v5 — correctifs globaux des figures (docs/onbuch-generation/tools/figfix)
\usepackage{adjustbox}\usepackage{etoolbox}\tcbuselibrary{raster}
% toute figure TikZ/pgfplots plus large que la ligne est réduite à la largeur disponible
\BeforeBeginEnvironment{tikzpicture}{\begin{adjustbox}{max width=\linewidth}}
\AfterEndEnvironment{tikzpicture}{\end{adjustbox}}
% légendes pgfplots lisibles par-dessus les courbes
\pgfplotsset{every axis/.append style={legend style={fill=white,fill opacity=0.92,text opacity=1,draw=black!15,inner sep=3pt}}}
% étiquettes d'arêtes (syntaxe edge["..."]) avec fond blanc
\tikzset{every edge quotes/.append style={fill=white,inner sep=1.2pt}}
%% FIGFIX-END
\begin{document}

\pagegarde{Les médias et leur impact sur l’opinion publique}{ECM}{1re Tech}{ECM – classe de Première technique}{N03}{3 séances (fiche de progression harmonisée)}{Cette leçon étudie les médias, leurs différentes formes et leur influence sur la formation de l'opinion publique au Cameroun. Elle présente les mécanismes de protection de l'opinion publique et analyse, à travers un travail dirigé, le rôle de la publicité dans les médias.
\vspace{4pt}
\begin{multicols}{2}\small
\begin{itemize}
\item À la fin de cette leçon, je sais définir les notions de média, d'opinion publique et de manipulation de l'information
\item Identifier et classer les différents types de médias présents au Cameroun (presse écrite, radio, télévision, médias numériques)
\item Expliquer à partir d'exemples camerounais comment les médias influencent l'opinion publique
\item Distinguer information, désinformation, propagande et publicité dans un document médiatique
\item Décrire les institutions et les textes qui protègent l'opinion publique au Cameroun
\item Analyser un message publicitaire diffusé dans les médias et en dégager les techniques de persuasion
\end{itemize}\end{multicols}}

\begin{sommaire}
\begin{multicols}{2}
\begin{enumerate}
\item Les médias : définitions et typologie
\item L'opinion publique : formation et caractéristiques
\item L'impact des médias sur l'opinion publique
\item La protection de l'opinion publique au Cameroun
\item TD3 : La publicité dans les médias
\item Activité d'intégration
\item Méthodes et exercices résolus
\item Exercices
\item Corrigés des exercices
\item Fiche bilan
\end{enumerate}
\end{multicols}
\end{sommaire}

\begin{objectifs}
\begin{itemize}
\item À la fin de cette leçon, je sais définir les notions de média, d'opinion publique et de manipulation de l'information
\item Identifier et classer les différents types de médias présents au Cameroun (presse écrite, radio, télévision, médias numériques)
\item Expliquer à partir d'exemples camerounais comment les médias influencent l'opinion publique
\item Distinguer information, désinformation, propagande et publicité dans un document médiatique
\item Décrire les institutions et les textes qui protègent l'opinion publique au Cameroun
\item Analyser un message publicitaire diffusé dans les médias et en dégager les techniques de persuasion
\item Appliquer les notions de l'unité à des situations concrètes de la vie courante (rumeurs, campagnes médiatiques, réseaux sociaux)
\item Rédiger et justifier une démarche, une réponse ou une conclusion à partir de documents fournis
\end{itemize}
\end{objectifs}

\begin{prerequis}
\begin{itemize}
\item Notions de citoyenneté et de vie en société (classes de 5e-4e)
\item Les droits et devoirs du citoyen (ECM, Seconde)
\item La liberté d'expression et ses limites (ECM, Seconde)
\item Notions de base sur l'État et ses institutions (ECM, début de Première)
\end{itemize}
\end{prerequis}

\newpage

\coursec{Les médias : définitions et typologie}

\textbf{Situation d'accroche.} En avril 2024, une rumeur circulant sur WhatsApp affirme qu'une épidémie frappe les marchés de Yaoundé : en quelques heures, des familles retirent leurs enfants de l'école et des commerçants ferment leurs boutiques. La CRTV et le MINCOM démentent, mais la panique a déjà fait des dégâts. Comment une simple information peut-elle bouleverser le comportement de milliers de Camerounais ? Qui contrôle ce que nous lisons, écoutons et regardons ? Et comment l'État protège-t-il les citoyens contre les dérives des médias ? Cette leçon nous aide à comprendre le pouvoir des médias et les moyens de s'en protéger.

\textbf{Problématique.} Qu'est-ce qu'un média ? Quels sont les grands types de médias au Cameroun et quelles fonctions remplissent-ils dans la vie de la cité ?

\courssub{Qu'est-ce qu'un média ?}

Commençons par une observation simple. Le matin, dans un quartier de Douala, un apprenti mécanicien écoute la radio dans l'atelier, son patron lit \emph{Mutations} en buvant son café, et sa sœur consulte Facebook sur son téléphone. Tous trois reçoivent de l'information, mais par des \cle{canaux} différents. Ces canaux, ce sont les médias.

\begin{definition}[Média et média de masse]
Un \cle{média} (du latin \emph{medium}, « milieu, intermédiaire ») est un \textbf{moyen de diffusion de l'information vers un large public}. On parle de \cle{média de masse} (ou \emph{mass media}) lorsque ce moyen de diffusion atteint simultanément un très grand nombre de personnes, souvent anonymes et dispersées sur un vaste territoire. La presse écrite, la radio, la télévision et Internet sont les grands médias de masse.
\end{definition}

Retenons trois caractéristiques essentielles d'un média de masse :
\begin{itemize}
\item il s'adresse à un \textbf{public nombreux} (des milliers, voire des millions de personnes) ;
\item il diffuse l'information \textbf{rapidement et simultanément} ;
\item il repose sur une \textbf{technique de reproduction et de transmission} (imprimerie, ondes hertziennes, réseaux numériques).
\end{itemize}

\begin{attention}[Média ou message ?]
Il ne faut pas confondre le \cle{média}, qui est le \textbf{canal de diffusion} (le journal, le poste de radio, l'application WhatsApp), et le \cle{message}, qui est le \textbf{contenu diffusé} (l'article, le reportage, la rumeur). Un même message peut passer par plusieurs médias, et un même média peut transmettre des messages vrais ou faux. Critiquer un message ne signifie donc pas condamner le média qui l'a porté.
\end{attention}

\courssub{La typologie des médias au Cameroun}

On classe traditionnellement les médias en deux grandes familles : les \cle{médias traditionnels} (presse écrite, radio, télévision), apparus avant l'ère d'Internet, et les \cle{médias numériques}, nés des nouvelles technologies de l'information et de la communication.

\textbf{1. La presse écrite.} C'est le plus ancien média de masse. Au Cameroun, on distingue :
\begin{itemize}
\item \emph{Cameroon Tribune}, quotidien national édité par la SOPECAM, société d'État ;
\item des journaux privés comme \emph{Le Messager}, \emph{Mutations} ou \emph{L'Œil du Sahel}, journal paraissant dans le Grand Nord.
\end{itemize}
La presse écrite offre des analyses approfondies, mais elle suppose de savoir lire et d'acheter le journal : son public est surtout urbain et scolarisé.

\textbf{2. La radio.} C'est le média le plus accessible : un poste coûte peu cher, fonctionne à piles, et les émissions sont en langues nationales comme en langues locales. On trouve :
\begin{itemize}
\item \emph{CRTV Radio}, la radio publique nationale, avec ses stations régionales ;
\item des radios privées comme \emph{Equinoxe Radio} ou \emph{Sweet FM} ;
\item des \cle{radios communautaires et rurales}, proches des populations, qui traitent des préoccupations locales (agriculture, santé, marchés).
\end{itemize}

\textbf{3. La télévision.} Elle combine image et son, ce qui lui donne une forte force de conviction. Les principales chaînes sont la \emph{CRTV} (chaîne publique), \emph{Canal 2 International}, \emph{STV} et \emph{Equinoxe TV}. Son coût (téléviseur, électricité, antenne) limite cependant sa diffusion en milieu rural.

\textbf{4. Les médias numériques.} Ce sont les sites d'information en ligne, les blogs et surtout les \cle{réseaux sociaux} (Facebook, WhatsApp, TikTok), très populaires chez les jeunes. Leur originalité : chacun peut y devenir émetteur d'information, sans formation ni contrôle préalable. C'est à la fois une richesse (liberté d'expression) et un danger (rumeurs, fausses nouvelles).

\begin{popfigure}
\begin{center}
\begin{tikzpicture}[font=\small,
  fam/.style={draw=popdark, fill=popblueL, rounded corners, thick, minimum width=3.4cm, minimum height=0.9cm, align=center},
  med/.style={draw=popblue, fill=popcream, rounded corners, thick, minimum width=3.1cm, minimum height=0.8cm, align=center},
  ex/.style={align=center, text=popmuted, font=\footnotesize}]
\node[draw=popdark, fill=poporangeL, rounded corners, thick, minimum width=5cm, minimum height=1cm, font=\bfseries] (root) at (0,0) {LES MÉDIAS};
\node[fam] (trad) at (-4.6,-2) {Médias traditionnels};
\node[fam] (num) at (4.6,-2) {Médias numériques};
\node[med] (presse) at (-7.2,-4) {Presse écrite};
\node[med] (radio) at (-4.6,-4) {Radio};
\node[med] (tv) at (-2,-4) {Télévision};
\node[med] (sites) at (2.4,-4) {Sites d'info, blogs};
\node[med] (rs) at (6.2,-4) {Réseaux sociaux};
\node[ex, align=center] at (-7.2,-4.9){Cameroon Tribune,\\ Mutations, L'Œil du Sahel};
\node[ex, align=center] at (-4.6,-4.9){CRTV Radio, Equinoxe,\\ radios rurales};
\node[ex, align=center] at (-2,-4.9){CRTV, Canal 2,\\ STV, Equinoxe TV};
\node[ex, align=center] at (2.4,-4.9){CamerounWeb,\\ blogs d'opinion};
\node[ex, align=center] at (6.2,-4.9){Facebook, WhatsApp,\\ TikTok};
\draw[thick, popdark] (root) -- (trad);
\draw[thick, popdark] (root) -- (num);
\draw[thick, popblue] (trad) -- (presse);
\draw[thick, popblue] (trad) -- (radio);
\draw[thick, popblue] (trad) -- (tv);
\draw[thick, popblue] (num) -- (sites);
\draw[thick, popblue] (num) -- (rs);
\end{tikzpicture}
\end{center}
\legende{Classification des médias en deux grandes familles, avec des exemples camerounais pour chaque type.}
\end{popfigure}

\courssub{Les fonctions des médias}

Pourquoi les médias occupent-ils une place si importante dans nos sociétés ? Parce qu'ils remplissent cinq grandes fonctions :
\begin{enumerate}
\item \textbf{Informer} : rapporter les faits (politique, économie, sport, faits divers) de manière rapide et véridique ;
\item \textbf{Éduquer} : diffuser des connaissances (campagnes de vaccination, émissions agricoles, programmes scolaires) ;
\item \textbf{Divertir} : offrir des loisirs (musique, films, jeux, séries) ;
\item \textbf{Former l'opinion} : proposer des analyses et des débats qui aident les citoyens à se forger un jugement ;
\item \textbf{Contrôler les pouvoirs} : dénoncer les abus, enquêter sur la mauvaise gestion, interpeller les gouvernants.
\end{enumerate}

\begin{definition}[Les médias, « quatrième pouvoir »]
On appelle \cle{contre-pouvoir} toute instance capable de surveiller, de critiquer et de limiter l'action des pouvoirs publics. Parce qu'ils enquêtent sur les gouvernants et dénoncent les abus, les médias sont surnommés le \textbf{quatrième pouvoir}, à côté des trois pouvoirs classiques de l'État : le pouvoir exécutif, le pouvoir législatif et le pouvoir judiciaire. Cette fonction de contre-pouvoir est essentielle en démocratie, mais elle exige des journalistes rigueur, honnêteté et vérification des faits.
\end{definition}

\begin{exemplebox}[Le contre-pouvoir en action]
Lorsqu'un journal enquête sur le détournement de fonds destinés à la construction d'un lycée technique et publie ses preuves, l'opinion s'émeut et l'État ouvre une enquête judiciaire. Le média a joué son rôle de contre-pouvoir : il a alerté les citoyens et poussé les autorités à agir. À l'inverse, lorsqu'une page Facebook accuse sans preuve un commerçant de vendre de la viande avariée, elle détruit sa réputation injustement : le pouvoir des médias devient alors dangereux.
\end{exemplebox}

\courssub{Les médias au Cameroun en chiffres}

Pour mesurer le poids réel de chaque média, on utilise les \cle{taux de pénétration}, c'est-à-dire la part de la population qui a accès à un média donné. Au Cameroun, les données disponibles (enquêtes nationales, rapports internationaux) montrent une hiérarchie nette : la radio reste le média le plus accessible (environ 80\,\% des ménages), devant la télévision (environ 60\,\% des ménages, surtout en ville), tandis qu'Internet progresse rapidement (plus de 40\,\% de la population en 2024, contre moins de 15\,\% en 2015). La presse écrite reste confidentielle (moins de 10\,\% de lecteurs réguliers).

\begin{popfigure}
\begin{center}
\begin{tikzpicture}
\begin{axis}[
  ybar, bar width=16pt,
  width=0.92\linewidth, height=6.2cm,
  ymin=0, ymax=10,
  ylabel={Usagers réguliers (en millions, données simulées)},
  symbolic x coords={Radio, Télévision, Internet, Presse écrite},
  xtick=data,
  nodes near coords,
  every node near coord/.append style={font=\footnotesize},
  ymajorgrids, grid style={popline!40},
  axis line style={popdark},
]
\addplot[fill=popblue!70, draw=popblue] coordinates {(Radio,8.5) (Télévision,6.2) (Internet,5.8) (Presse écrite,1.2)};
\end{axis}
\end{tikzpicture}
\end{center}
\legende{Nombre estimé d'auditeurs, téléspectateurs ou usagers réguliers par type de média au Cameroun (données simulées réalistes, population d'environ 28 millions d'habitants).}
\end{popfigure}

\begin{exemplebox}[L'essor spectaculaire d'Internet]
En 2015, le Cameroun comptait environ 3 millions d'internautes ; en 2024, ils sont plus de 11 millions, soit une multiplication par près de 3,7 en moins de dix ans. Cette croissance s'explique par la baisse du prix des smartphones, l'extension des réseaux 3G/4G et la jeunesse de la population. Conséquence directe : les réseaux sociaux deviennent, surtout chez les jeunes, la première source d'information — devant parfois la radio et la télévision.
\end{exemplebox}

\begin{savaistu}[La radio, reine des campagnes]
Dans les zones rurales du Cameroun, la radio reste le média le plus écouté : plus de 8 ménages ruraux sur 10 y ont accès, contre moins de 4 sur 10 pour la télévision. Les raisons sont simples : le poste à piles ne dépend pas du réseau électrique, les radios communautaires émettent en langues locales (fufuldé, ewondo, douala, gbaya\ldots) et les émissions traitent des préoccupations quotidiennes des villageois : prix des denrées, techniques agricoles, santé, annonces du chef de village.
\end{savaistu}

\begin{exoresolu}[Identifier média, message et fonction]
Énoncé. Une radio communautaire de l'Extrême-Nord diffuse une émission en fufuldé expliquant aux éleveurs comment vacciner leurs bœufs contre la péripneumonie. Pour cette situation, identifie : a) le média ; b) le message ; c) la fonction principale remplie.
\tcblower
Solution détaillée.
a) Le \textbf{média} est le canal de diffusion : la radio communautaire (média traditionnel, de proximité).
b) Le \textbf{message} est le contenu diffusé : les conseils pratiques de vaccination des bovins.
c) La \textbf{fonction principale} est la fonction d'\textbf{éducation} : la radio transmet des connaissances utiles qui améliorent la vie des éleveurs. On peut ajouter qu'elle remplit aussi une fonction d'information, mais l'objectif premier est ici pédagogique.
\end{exoresolu}

\begin{aretenir}[L'essentiel de la section]
\begin{itemize}
\item Un \cle{média} est un moyen de diffusion de l'information vers un large public ; un \cle{média de masse} touche simultanément des millions de personnes.
\item On distingue les \cle{médias traditionnels} (presse écrite, radio, télévision) et les \cle{médias numériques} (sites, blogs, réseaux sociaux).
\item Les médias remplissent cinq fonctions : informer, éduquer, divertir, former l'opinion et contrôler les pouvoirs (\cle{contre-pouvoir} ou « quatrième pouvoir »).
\item Au Cameroun, la radio domine (surtout en milieu rural), la télévision est urbaine, la presse écrite est faible, et Internet progresse très vite.
\item Il faut toujours distinguer le \cle{média} (le canal) du \cle{message} (le contenu).
\end{itemize}
\end{aretenir}

\begin{exercice}[À toi de jouer]
1. Définis en deux phrases le mot « média » et donne deux exemples de médias présents dans ton quartier ou ton village.
2. Classe les médias suivants dans le tableau traditionnels / numériques : CRTV, WhatsApp, \emph{Mutations}, TikTok, Equinoxe Radio, un blog sportif.
3. Pourquoi dit-on que les médias sont un « quatrième pouvoir » ? Illustre ta réponse par un exemple.
4. Explique pourquoi la radio reste le média le plus accessible en milieu rural camerounais.
\end{exercice}

\begin{corrige}[Exercice n°1]
1. Un média est un moyen technique de diffuser une information vers un large public. Exemples attendus : une radio locale, un journal, la télévision, Facebook, etc.
2. Traditionnels : CRTV, \emph{Mutations}, Equinoxe Radio. Numériques : WhatsApp, TikTok, blog sportif.
3. Parce qu'ils surveillent et critiquent les trois pouvoirs de l'État (exécutif, législatif, judiciaire) et dénoncent leurs abus ; exemple : une enquête journalistique sur un détournement de fonds publics.
4. Le poste de radio est bon marché, fonctionne à piles sans électricité, et les radios rurales émettent en langues locales sur des sujets proches des villageois.
\end{corrige}

\coursec{L'opinion publique : formation et caractéristiques}

Dans la section précédente, nous avons défini les médias et dressé leur typologie (presse écrite, radio, télévision, médias numériques). Nous avons vu qu'ils sont des intermédiaires entre l'information et le public. Mais pourquoi les médias occupent-ils une place si importante dans nos sociétés ? Parce qu'ils agissent sur quelque chose de précieux et de fragile : \cle{l'opinion publique}. C'est cette notion que nous allons maintenant étudier : qu'est-ce que l'opinion publique, comment se forme-t-elle, qui la façonne et comment la mesure-t-on ?

\courssub{Définition de l'opinion publique}

Commençons par une situation simple. Dans ton quartier, la mairie annonce la fermeture du marché pour travaux. Très vite, les commerçants protestent, les clients s'inquiètent, les discussions s'animent au carrefour et sur les groupes WhatsApp du quartier. Certains sont favorables aux travaux, d'autres y sont hostiles. Peu à peu, un sentiment dominant se dégage : << les travaux sont nécessaires, mais il faut un site de repli pour les vendeuses >>. Ce sentiment partagé par une grande partie de la population du quartier, c'est déjà une forme d'opinion publique.

\begin{definition}[Opinion publique]
L'\cle{opinion publique} est l'ensemble des idées, des attitudes et des jugements partagés par une population (ou une fraction importante de celle-ci) sur une question d'intérêt général, à un moment donné. Elle porte sur des sujets qui concernent la vie collective : la politique, l'économie, la sécurité, l'éducation, la santé, les mœurs.
\end{definition}

Trois éléments de cette définition méritent d'être soulignés :

\begin{itemize}
\item \textbf{Un objet d'intérêt général} : l'opinion publique ne porte pas sur les goûts personnels (préférer le ndolé au eru), mais sur des questions qui engagent la collectivité (le prix du carburant, la construction d'une route, une élection).
\item \textbf{Un partage} : il ne s'agit pas de l'avis d'une seule personne, mais d'idées \emph{communes} à un groupe large. Une opinion devient << publique >> lorsqu'elle est exprimée et connue de tous.
\item \textbf{Un moment donné} : l'opinion publique est mouvante. Elle peut changer en quelques semaines sous l'effet d'un événement, d'une révélation médiatique ou d'une crise.
\end{itemize}

\begin{attention}[Une opinion publique n'est ni unanime ni forcément rationnelle]
Deux pièges à éviter :
\begin{itemize}
\item L'opinion publique n'est \textbf{jamais unanime}. Sur toute question importante, il existe des courants divergents : majorité, minorités, indécis. Dire << l'opinion publique pense que... >> est toujours un raccourci qui désigne la tendance dominante.
\item L'opinion publique n'est pas toujours \textbf{rationnelle}. Elle peut être influencée par les émotions, les rumeurs, les préjugés ou les fausses informations. Une opinion largement partagée n'est pas pour autant une opinion vraie ou juste.
\end{itemize}
\end{attention}

\courssub{Les acteurs de la formation de l'opinion publique}

L'opinion publique ne naît pas spontanément : elle est \emph{fabriquée} par de nombreux acteurs qui interagissent. Observons comment un jeune apprenti en mécanique se forge un avis sur les élections : il entend les informations à la radio de l'atelier, il discute le soir avec sa famille, il écoute son chef d'atelier qui est très respecté, il suit un jeune influenceur sur TikTok, et il retient ce qu'il a appris en classe d'ECM. Chacun de ces canaux a contribué à former son opinion.

\begin{itemize}
\item \textbf{Les médias} : ils sélectionnent les sujets dont on parle (on dit qu'ils << font l'agenda >>), hiérarchisent les informations et donnent la parole à certains acteurs plutôt qu'à d'autres. Leur rôle est central, comme nous le verrons dans la section suivante.
\item \textbf{La famille} : c'est le premier cadre de socialisation. Les convictions politiques, religieuses et morales des parents influencent durablement celles des enfants.
\item \textbf{L'école} : par l'éducation à la citoyenneté, elle transmet des valeurs (justice, tolérance, esprit critique) et apprend à débattre de manière argumentée.
\item \textbf{Les leaders d'opinion} et \textbf{les influenceurs} : ce sont des personnes dont la parole pèse davantage que celle des autres.
\end{itemize}

\begin{definition}[Leader d'opinion et influenceur]
Un \cle{leader d'opinion} est une personne qui, par son statut, sa compétence ou son prestige (chef traditionnel, religieux, professeur, syndicaliste, artiste, chef d'entreprise), exerce une influence sur les idées et les comportements de ceux qui l'écoutent.

Un \cle{influenceur} est une personne qui, grâce à sa notoriété sur les réseaux sociaux (YouTube, TikTok, Facebook, Instagram), oriente les opinions, les goûts et les choix de sa communauté d'abonnés, souvent dans un cadre commercial ou politique.
\end{definition}

\begin{popfigure}
\begin{center}
\begin{tikzpicture}[scale=1]
\node[circle, draw=popblue, fill=popblueL, thick, minimum size=2.6cm, align=center, font=\small\bfseries] (op) at (0,0) {Opinion\\publique};
\node[circle, draw=poporange, fill=poporangeL, thick, minimum size=1.7cm, align=center, font=\footnotesize] (med) at (90:3.4) {Médias};
\node[circle, draw=popgreen, fill=popgreenL, thick, minimum size=1.7cm, align=center, font=\footnotesize] (fam) at (162:3.4) {Famille};
\node[circle, draw=poppurple, fill=poppurpleL, thick, minimum size=1.7cm, align=center, font=\footnotesize] (eco) at (234:3.4) {École};
\node[circle, draw=popgold, fill=popgoldL, thick, minimum size=1.7cm, align=center, font=\footnotesize] (lea) at (306:3.4) {Leaders\\d'opinion};
\node[circle, draw=poppink, fill=poppinkL, thick, minimum size=1.7cm, align=center, font=\footnotesize] (inf) at (18:3.4) {Influenceurs};
\draw[-{Stealth}, thick, popmuted] (med) -- (op);
\draw[-{Stealth}, thick, popmuted] (fam) -- (op);
\draw[-{Stealth}, thick, popmuted] (eco) -- (op);
\draw[-{Stealth}, thick, popmuted] (lea) -- (op);
\draw[-{Stealth}, thick, popmuted] (inf) -- (op);
\end{tikzpicture}
\end{center}
\legende{Les cinq grandes catégories d'acteurs qui influencent la formation de l'opinion publique : médias, famille, école, leaders d'opinion et influenceurs. Ces influences se combinent et parfois se contredisent.}
\end{popfigure}

\courssub{Les étapes de formation d'une opinion publique}

Une opinion publique ne surgit pas d'un coup : elle suit un processus que l'on peut décrire en quatre étapes. Prenons l'exemple d'une rumeur de fermeture d'un lycée technique dans une ville.

\begin{enumerate}
\item \textbf{L'information} : un événement survient et est diffusé (par les médias, les réseaux sociaux, le bouche-à-oreille). Sans information, pas d'opinion : on ne peut pas avoir d'avis sur ce qu'on ignore.
\item \textbf{La discussion} : les gens échangent, comparent leurs points de vue, débattent en famille, au travail, dans les buvettes, sur les groupes WhatsApp. Les arguments s'affrontent.
\item \textbf{La cristallisation} : peu à peu, les positions se stabilisent. Des courants d'opinion se forment : les favorables, les opposants, les indécis. L'opinion devient identifiable.
\item \textbf{L'expression} : l'opinion se manifeste publiquement : pétitions, manifestations, votes, articles de presse, messages sur les réseaux sociaux, débats à la radio.
\end{enumerate}

\begin{popfigure}
\begin{center}
\begin{tikzpicture}[node distance=0.55cm,
etape/.style={rectangle, rounded corners=4pt, draw=popblue, fill=popblueL, thick, minimum width=2.6cm, minimum height=1.1cm, align=center, font=\small}]
\node[etape, align=center] (e1){1. Information\\(l'événement est diffusé)};
\node[etape, right=of e1, align=center] (e2){2. Discussion\\(échanges et débats)};
\node[etape, right=of e2, align=center] (e3){3. Cristallisation\\(positions stabilisées)};
\node[etape, right=of e3, align=center] (e4){4. Expression\\(opinion manifestée)};
\draw[-{Stealth}, very thick, poporange] (e1) -- (e2);
\draw[-{Stealth}, very thick, poporange] (e2) -- (e3);
\draw[-{Stealth}, very thick, poporange] (e3) -- (e4);
\end{tikzpicture}
\end{center}
\legende{Les quatre étapes de formation d'une opinion publique, de la diffusion de l'information jusqu'à son expression publique.}
\end{popfigure}

\begin{exemplebox}[Un cas camerounais : les grands débats nationaux]
Au Cameroun, plusieurs sujets ont donné lieu à de vastes débats publics relayés par les médias : l'organisation de la Coupe d'Afrique des Nations (CAN) et la construction des stades, le coût de la vie et les prix des denrées, la réforme du système éducatif, la lutte contre la corruption ou encore la gestion des ressources naturelles. À chaque fois, on observe le même schéma : un fait est révélé (information), les émissions de débat radiophoniques et les plateaux télévisés s'en emparent (discussion), des camps se dessinent (cristallisation), puis des citoyens, des associations ou des partis prennent publiquement position (expression).
\end{exemplebox}

\courssub{Les instruments de mesure de l'opinion publique}

Comment savoir ce que << l'opinion >> pense réellement ? On ne peut pas interroger environ 28 millions de Camerounais un par un. On utilise donc des instruments de mesure :

\begin{itemize}
\item \textbf{Les sondages} : on interroge un échantillon représentatif de la population (par exemple 1\,000 personnes choisies selon des critères d'âge, de sexe, de région) et on généralise les résultats avec une marge d'erreur.
\item \textbf{Les enquêtes} : plus approfondies que les sondages, elles combinent questionnaires et entretiens pour comprendre les motivations des opinions.
\item \textbf{Les débats publics} : émissions de radio et de télévision, forums, réunions communautaires, où l'on entend directement les arguments des citoyens.
\item \textbf{Les réactions sur les réseaux sociaux} : commentaires, partages, mots-clés tendances donnent une indication rapide, mais biaisée, car tous les citoyens n'y sont pas présents.
\end{itemize}

\begin{methode}[Analyser un sondage en trois questions]
Face à un sondage publié dans la presse ou sur les réseaux sociaux, pose-toi toujours trois questions :
\begin{enumerate}
\item \textbf{Qui a été interrogé ?} Vérifie l'\cle{échantillon} : combien de personnes, comment ont-elles été choisies, sont-elles représentatives de la population (âge, sexe, région, milieu social) ? Un sondage réalisé uniquement auprès des abonnés d'une page Facebook ne vaut rien pour toute la population.
\item \textbf{Quelle question a été posée, et par qui ?} La formulation peut orienter la réponse (<< Êtes-vous d'accord avec cette excellente réforme ? >>). Regarde aussi qui a commandité le sondage : un parti politique, une entreprise, un institut indépendant ?
\item \textbf{Quelle est la marge d'erreur ?} Un sondage donne une estimation, pas une certitude. Pour un échantillon d'environ 1\,000 personnes, la marge d'erreur est d'environ $\pm 3$ points. Si un candidat est crédité de 48\,\% et l'autre de 47\,\%, l'écart est trop faible pour conclure.
\end{enumerate}
\end{methode}

\begin{exoresolu}[Lire un sondage avec esprit critique]
\textbf{Énoncé.} Un site web publie : << Sondage exclusif : 85\,\% des Camerounais soutiennent le projet X ! >> En lisant l'article, tu apprends que 200 personnes ont répondu à un questionnaire diffusé sur la page Facebook du promoteur du projet X. Analyse ce sondage à l'aide de la méthode vue en cours.
\tcblower
\textbf{Solution détaillée.}
\begin{enumerate}
\item \emph{L'échantillon} : 200 personnes, c'est peu. Surtout, elles ont été recrutées parmi les visiteurs de la page du promoteur du projet : ce sont des personnes déjà intéressées, voire déjà favorables au projet. L'échantillon n'est donc \textbf{pas représentatif} de la population camerounaise.
\item \emph{La question et le commanditaire} : le sondage est réalisé par le promoteur lui-même, qui a intérêt à un résultat favorable. La formulation de la question n'est pas indiquée : elle a pu être orientée.
\item \emph{La marge d'erreur} : avec 200 personnes, la marge d'erreur est d'environ $\pm 7$ points, et surtout le biais de sélection rend le chiffre de 85\,\% non généralisable.
\end{enumerate}
\textbf{Conclusion} : ce sondage ne permet pas d'affirmer que 85\,\% des Camerounais soutiennent le projet X. Il mesure seulement l'avis des visiteurs d'une page favorable au projet. Il faut donc accueillir ce chiffre avec une grande prudence.
\end{exoresolu}

\begin{savaistu}[Les sondages sont encadrés au Cameroun]
Au Cameroun, la publication des sondages d'opinion est encadrée par la loi, en particulier en période électorale. La loi n° 90/052 du 19 décembre 1990 sur la communication sociale, modifiée depuis, et la réglementation électorale imposent des règles de transparence : tout sondage publié doit indiquer l'institut qui l'a réalisé, le commanditaire, la taille de l'échantillon, les dates du terrain et la formulation des questions. En période de campagne électorale, la publication de sondages sur les intentions de vote est strictement encadrée afin de ne pas manipuler les électeurs. Objectif : protéger l'opinion publique contre les faux sondages fabriqués pour influencer le vote.
\end{savaistu}

\courssub{Opinion publique et démocratie}

Dans une démocratie, l'opinion publique n'est pas un simple décor : elle est un acteur politique à part entière.

\begin{itemize}
\item \textbf{Dans les élections} : c'est l'opinion publique qui, en définitive, fait et défait les majorités. Les candidats cherchent à la convaincre ; les électeurs expriment leur opinion par le vote.
\item \textbf{Dans les politiques publiques} : les gouvernants tiennent compte de l'opinion pour lancer, ajuster ou retirer des réformes. Une mesure très impopulaire peut être abandonnée sous la pression de l'opinion.
\item \textbf{Dans la dénonciation des abus} : l'opinion publique joue un rôle de contre-pouvoir. La dénonciation de la corruption, des détournements de fonds ou des violences, relayée par les médias et les réseaux sociaux, peut conduire à des enquêtes, des sanctions ou des démissions.
\end{itemize}

\begin{aretenir}[L'essentiel de la section]
\begin{itemize}
\item L'opinion publique est l'ensemble des idées, attitudes et jugements partagés par une population sur une question d'intérêt général, à un moment donné.
\item Elle est façonnée par cinq grandes catégories d'acteurs : les médias, la famille, l'école, les leaders d'opinion et les influenceurs.
\item Sa formation suit quatre étapes : information, discussion, cristallisation, expression.
\item On la mesure par les sondages, les enquêtes, les débats publics et les réactions sur les réseaux sociaux — avec prudence et esprit critique.
\item Elle n'est jamais unanime ni forcément rationnelle, mais elle joue un rôle essentiel en démocratie : élections, politiques publiques, dénonciation des abus.
\end{itemize}
\end{aretenir}

\begin{exercice}[À toi de jouer]
\begin{enumerate}
\item Définis l'opinion publique et donne deux exemples de questions d'intérêt général débattues actuellement dans ta localité.
\item Classe les acteurs suivants selon qu'il s'agit de la famille, de l'école, des médias, d'un leader d'opinion ou d'un influenceur : un chef de village ; une émission de débat à la télévision ; un professeur d'ECM ; un chanteur suivi par 2 millions d'abonnés ; une discussion au repas familial.
\item Un sondage sur les intentions de vote est réalisé auprès de 50 étudiants d'un seul campus universitaire. Peut-on le généraliser à toute la jeunesse camerounaise ? Justifie ta réponse en rédigeant un paragraphe argumenté.
\end{enumerate}
\end{exercice}

\begin{corrige}[Exercice : À toi de jouer]
\begin{enumerate}
\item L'opinion publique est l'ensemble des idées, attitudes et jugements partagés par une population sur une question d'intérêt général, à un moment donné. Exemples possibles : l'état des routes, le prix du transport urbain, la sécurité dans le quartier, la construction d'un marché.
\item Chef de village : leader d'opinion ; émission de débat télévisée : médias ; professeur d'ECM : école ; chanteur suivi sur les réseaux : influenceur ; discussion au repas : famille.
\item Non, on ne peut pas généraliser ce sondage. L'échantillon est trop petit (50 personnes) et surtout non représentatif : il ne comprend que des étudiants d'un seul campus, alors que la jeunesse camerounaise comprend aussi des élèves, des apprentis, des jeunes travailleurs et des jeunes ruraux. Les résultats mesurent l'opinion de ce groupe précis, pas celle de toute la jeunesse. Pour conclure, il faudrait un échantillon plus large et diversifié, avec une marge d'erreur connue.
\end{enumerate}
\end{corrige}

Maintenant que nous savons ce qu'est l'opinion publique, comment elle se forme et comment on la mesure, nous pouvons aborder la question centrale de la leçon : quel est l'impact réel des médias sur cette opinion ? C'est l'objet de la section suivante.

\coursec{Les médias et leur impact sur l'opinion publique}

\coursec{Les médias : définitions et typologie}

\courssub{Définition et fonctions sociales des médias}

\begin{definition}[Média]
Un \cle{média} (pluriel : médias) est un moyen de communication de masse permettant de diffuser une information, un message ou un contenu culturel auprès d'un large public, de manière simultanée et rapide. Le terme vient du latin \emph{medium} (intermédiaire).
\end{definition}

\begin{propriete}[Fonctions fondamentales des médias (Laswell, 1948 ; McQuail)]
\begin{enumerate}
    \item \textbf{Information (surveillance de l'environnement) :} Signaler les événements, menaces, opportunités (météo, actualité politique, alertes sanitaires).
    \item \textbf{Corrélation (interprétation) :} Expliquer, commenter, hiérarchiser les faits pour donner du sens (éditoriaux, analyses, débats).
    \item \textbf{Socialisation (transmission culturelle) :} Transmettre les valeurs, normes, histoire, langue d'une société d'une génération à l'autre (émissions éducatives, fictions, documentaires).
    \item \textbf{Divertissement (détente) :} Offrir une évasion, un loisir (séries, jeux, musique, sport).
    \item \textbf{Mobilisation :} Inciter à l'action collective (campagnes électorales, appels au don, sensibilisation civique).
\end{enumerate}
\end{propriete}

\courssub{Typologie des médias : classification technique et juridique}

\begin{center}
\begin{tabularx}{\linewidth}{|l|X|X|}\hline
\rowcolor{popblueL}
\textbf{Critère} & \textbf{Catégories} & \textbf{Exemples au Cameroun} \\ \hline
\textbf{Nature du support} & \textbf{Presse écrite} (périodiques, quotidiens) & \emph{Cameroon Tribune}, \emph{Le Jour}, \emph{Mutations}, \emph{L'Œil du Sahel} \\ \cline{2-3}
& \textbf{Radio} (hertzienne, FM, web) & CRTV Radio, Radio Balafon, RFI, BBC, radios communautaires (Radio Medumba, etc.) \\ \cline{2-3}
& \textbf{Télévision} (hertzienne, satellite, TNT, web) & CRTV, Canal 2 International, Equinoxe TV, STV, Vision 4 \\ \cline{2-3}
& \textbf{Médias numériques / En ligne} & Sites d'info (\emph{Actu Cameroun}, \emph{Data Cameroon}), réseaux sociaux (Facebook, WhatsApp, X, TikTok), blogs, podcasts \\ \hline
\textbf{Statut juridique} & \textbf{Service public} (mission de service universel, financés par redevance/subventions) & CRTV (Radio + TV), \emph{Cameroon Tribune} \\ \cline{2-3}
& \textbf{Privés commerciaux} (but lucratif, financement pub/abonnements) & Groupe \emph{L'Anecdote} (Equinoxe, \emph{Le Jour}), Groupe \emph{Tele-Sud} (Canal 2, \emph{Mutations}) \\ \cline{2-3}
& \textbf{Associatifs / Communautaires} (non lucratif, proximité) & Radios communautaires (ex. Radio Yet à Yagoua), presse d'opinion \\ \hline
\textbf{Périmètre de diffusion} & \textbf{Nationaux} & CRTV, \emph{Cameroon Tribune}, principaux quotidiens privés \\ \cline{2-3}
& \textbf{Régionaux / Locaux} & Radios communautaires, correspondants locaux, journaux de province \\ \cline{2-3}
& \textbf{Internationaux} & RFI, France 24, TV5Monde, BBC, VOA, Deutsche Welle, Jeune Afrique \\ \hline
\end{tabularx}
\end{center}

\begin{aretenir}[Distinction clé : Médias classiques vs Réseaux sociaux]
\begin{itemize}
    \item \textbf{Médias classiques (Mainstream) :} Structure éditoriale (rédacteur en chef, journalistes professionnels, déontologie, responsabilité légale), ligne éditoriale connue, vérification des faits (\emph{fact-checking}) avant publication, droit de réponse encadré.
    \item \textbf{Réseaux sociaux / Médias citoyens :} Publication instantanée sans filtre éditorial a priori, algorithmes de viralité, anonymat possible, responsabilité juridique diffuse (plateforme vs auteur), bulles de filtre (\emph{filter bubbles}).
\end{itemize}
\end{aretenir}

\courssub{L'écosystème médiatique camerounais : repères chiffrés (ordre de grandeur)}

\begin{center}
\begin{tabularx}{\linewidth}{|l|X|}\hline
\rowcolor{popblueL}
\textbf{Indicateur} & \textbf{Données récentes (estimations 2023-2024)} \\ \hline
Taux de pénétration d'internet & $\sim$ 45\% -- 50\% de la population (source : MINPOSTEL, UIT, DataReportal) \\ \hline
Utilisateurs actifs réseaux sociaux & $\sim$ 6 -- 7 millions de comptes (Facebook dominant, WhatsApp messagerie n°1) \\ \hline
Téléphonie mobile & $\sim$ 85\% -- 90\% de pénétration (multi-SIM fréquent) \\ \hline
Titres de presse écrite déclarés & $>$ 600 titres (majorité paraissent irrégulièrement) \\ \hline
Radios / TV autorisées (CNC) & $\sim$ 600 radios (dont $\sim$ 200 communautaires), $\sim$ 50 chaînes TV \\ \hline
\end{tabularx}
\end{center}

\begin{savaistu}[La fracture numérique au Cameroun]
Malgré la progression du mobile, l'accès à une information de qualité reste inégal : zones rurales mal couvertes (3G/4G instable), coût du data élevé par rapport au SMIG, illettrisme numérique chez les aînés, barrière linguistique (contenu majoritairement en français/anglais, peu en langues nationales). Les radios communautaires en langues locales restent le média de confiance n°1 en zone rurale.
\end{savaistu}

\coursec{L'opinion publique : formation et caractéristiques}

\courssub{Définition et conditions d'existence}

\begin{definition}[Opinion publique]
L'\cle{opinion publique} est l'ensemble des attitudes, croyances, jugements et volontés partagés par une majorité de citoyens sur une question d'intérêt général, à un moment donné, et qui s'exprime publiquement. Elle n'est pas la somme des opinions individuelles, mais une \textbf{construction sociale} issue de l'interaction.
\end{definition}

\begin{propriete}[Conditions de formation de l'opinion publique (Crespi, 1997)]
Pour qu'une opinion publique se constitue sur un sujet, trois conditions sont nécessaires :
\begin{enumerate}
    \item \textbf{Un objet commun :} Un thème d'intérêt général (prix du carburant, réforme éducative, sécurité, élection).
    \item \textbf{Une communication publique :} Échange d'arguments via des espaces publics (médias, places de marché, réseaux sociaux, parlement, associations).
    \item \textbf{Une prise de position collective :} Émergence de tendances majoritaires/minoritaires, cristallisation autour de leaders d'opinion (chefs traditionnels, influenceurs, politiques, experts).
\end{enumerate}
\end{propriete}

\courssub{Caractéristiques principales}

\begin{propriete}[Caractères de l'opinion publique]
\begin{itemize}
    \item \textbf{Volatilité / Instabilité :} Elle peut basculer rapidement sous l'effet d'un événement marquant (choc émotionnel, scandale, performance sportive).
    \item \textbf{Irrationalité partielle :} Elle repose souvent sur des affects, des symboles, des rumeurs, des préjugés (\emph{stéréotypes}), plus que sur une analyse rationnelle coûts/bénéfices.
    \item \textbf{Conformisme / Pression sociale :} Peur de l'isolement (\emph{spirale du silence}, Noelle-Neumann) : la minorité se tait, la majorité perçue s'amplifie.
    \item \textbf{Compétence limitée :} Le citoyen moyen n'a ni le temps, ni l'expertise pour tout comprendre (complexité technique des lois, économie, diplomatie) $\rightarrow$ recours aux \cle{raccourcis cognitifs} (heuristiques : confiance dans un leader, parti, média de référence).
    \item \textbf{Influence des leaders d'opinion :} \emph{Two-step flow} (Lazarsfeld) : Médias $\rightarrow$ Leaders d'opinion $\rightarrow$ Public. Au Cameroun : chefs traditionnels, imams/pasteurs, enseignants, influenceurs web, artistes.
\end{itemize}
\end{propriete}

\courssub{Opinion publique, opinion publiée et sondages}

\begin{attention}[Ne pas confondre]
\begin{itemize}
    \item \textbf{Opinion publique :} Réelle, diffuse, difficile à mesurer exactement.
    \item \textbf{Opinion publiée :} Ce qui s'exprime dans les médias, les tribunes, les réseaux sociaux. Elle surreprésente les plus bruyants, les connectés, les organisés.
    \item \textbf{Sondage d'opinion :} Outil statistique (échantillon représentatif, marge d'erreur, formulation neutre des questions). Au Cameroun, les sondages politiques sont rares, souvent commandités, et leur publication en période électorale est encadrée (Code électoral, CNC).
\end{itemize}
\end{attention}

\coursec{L'impact des médias sur l'opinion publique}

\courssub{L'impact positif des médias : informer, éduquer, sensibiliser}

Les médias remplissent des fonctions vitales dans une société démocratique. Au-delà du divertissement, ils sont le principal vecteur de la circulation de l'information.

\begin{definition}[Fonctions positives des médias au service de l'opinion publique]
\begin{enumerate}
    \item \textbf{Information rapide et massive :} Diffusion quasi instantanée des faits (actualité politique, économique, sociale, sportive) à l'échelle nationale et internationale.
    \item \textbf{Éducation des masses :} Par des émissions éducatives, des documentaires, des chroniques spécialisées (santé, droit, environnement), ils élèvent le niveau de connaissances des citoyens.
    \item \textbf{Sensibilisation et mobilisation :} Relais des campagnes d'intérêt public (vaccination, inscription sur les listes électorales, sécurité routière, lutte contre le choléra, prévention du VIH/SIDA, éducation à la paix).
\end{enumerate}
\end{definition}

\begin{exemplebox}[Campagnes de sensibilisation au Cameroun : l'exemple de la riposte au COVID-19]
Lors de la pandémie de COVID-19 (2020-2021), la CRTV (radio, télévision, web) et le réseau des radios communautaires (plus de 200 stations) ont diffusé des spots en langues nationales (fulfulde, bassa, ewondo, bamileke, pidgin, arabe choa) pour atteindre les populations rurales et peu scolarisées. Les médias ont ainsi contribué à lever les réticences culturelles, à expliquer les gestes barrières et à améliorer l'adhésion aux campagnes de vaccination, en relais du MINSANTE et de l'OMS.
\end{exemplebox}

\begin{aretenir}[Rôle démocratique : la fonction de « chien de garde » (\emph{watchdog})]
En assurant la transparence de la vie publique (débats parlementaires diffusés, comptes rendus du Conseil des ministres, enquêtes journalistiques d'investigation sur la gestion des deniers publics), les médias permettent aux citoyens de contrôler l'action gouvernementale et de voter en connaissance de cause. C'est la condition de la \cle{reddition de comptes} (\emph{accountability}).
\end{aretenir}

\courssub{L'impact négatif : manipulation, propagande, désinformation et \cle{fake news}}

Si les médias peuvent éclairer, ils peuvent aussi obscurcir. La rapidité de diffusion, l'absence de filtre éditorial sur les réseaux sociaux et les intérêts politiques ou commerciaux créent un terreau fertile pour la désinformation.

\begin{definition}[Désinformation, propagande et \cle{fake news}]
\begin{itemize}
    \item \textbf{Désinformation :} Diffusion \textbf{délibérée} d'informations fausses ou trompeuses dans le but de manipuler l'opinion publique ou de nuire à une personne, un groupe ou une institution. Diffère de l'erreur involontaire (\emph{misinformation}).
    \item \textbf{Propagande :} Communication systématique, souvent étatique ou partisane, visant à imposer une idéologie, à glorifier un pouvoir ou à diaboliser un adversaire, en utilisant l'émotion plus que la raison (symboles, slogans, répétition).
    \item \textbf{\cle{Fake news} (infox) :} Information fausse, fabriquée de toutes pièces ou détournée de son contexte, présentée sous l'apparence d'un article de presse légitime, dans le but de tromper, de générer du trafic (clics, monétisation) ou de nuire.
\end{itemize}
\end{definition}

\begin{definition}[Rumeur]
Information non vérifiée, transmise de bouche à oreille ou via les réseaux sociaux, qui se propage rapidement dans l'opinion publique, souvent en période d'incertitude ou de crise. Elle comble un \textbf{vide informationnel} (anxiété + ambiguïté = rumeur, loi d'Allport \& Postman).
\end{definition}

\begin{exemplebox}[Rumeurs virales au Cameroun : conséquences réelles]
\begin{itemize}
    \item \textbf{2020 (COVID-19) :} Rumeur WhatsApp : « L'eau de javel ou l'alcool à 90° ingérés guérissent du COVID-19 ». Résultat : dizaines de cas d'intoxication grave enregistrés aux urgences de l'Hôpital Central et de l'Hôpital Général de Yaoundé, et à Laquintinie (Douala).
    \item \textbf{2023 (Éducation) :} Fausse capture d'écran aux couleurs du MINESEC annonçant la « suppression du Probatoire » partagée massivement fin mai. Panique chez les élèves/parents. Démenti officiel du Ministre (Communiqué n°0000/MINESEC/SG/DREX) 48h plus tard.
    \item \textbf{2024 (Sécurité) :} Vidéo d'une explosion en Turquie (2022) repartagée avec la légende « Attentat à Douala ce matin ». Partages > 50k en 2h. Panique dans les marchés. Démenti Police/Minat.
\end{itemize}
\end{exemplebox}

\begin{exoresolu}[Analyse d'une \cle{fake news} circulant sur les réseaux sociaux]
\textbf{Énoncé :} Une publication Facebook, partagée 12 000 fois en 48h, montre une photo d'un camion militaire brûlé avec la légende : « Affrontements sanglants à Buea ce matin : 15 soldats tués par les séparatistes. L'armée se replie. » La publication est datée d'aujourd'hui. Vérifiez la véracité de cette information.

\textbf{Solution détaillée :}
\begin{enumerate}
    \item \textbf{Observation de la source (S) :} Le compte qui poste est anonyme, créé il y a 3 jours, sans mention légale, sans photo de profil identifiable. Ce n'est pas un média reconnu (CRTV, \emph{Cameroon Tribune}, \emph{Jeune Afrique}, AFP, Reuters, Data Cameroon).
    \item \textbf{Analyse de l'image (D - Date/Contexte) :} Une recherche d'image inversée (Google Images, TinEye, Yandex Images) révèle que la photo date de \textbf{mai 2018} et a été prise à \textbf{Fotokol (Extrême-Nord)} lors d'une attaque de Boko Haram, pas à Buea (Nord-Ouest).
    \item \textbf{Recoupement (R) :} Aucun média crédible ne relaie l'information ce jour. Le communiqué officiel du MINDEF (Ministère de la Défense) de ce jour ne mentionne aucun incident à Buea. Le compte officiel du MINDEF sur X (Twitter) n'a pas publié d'alerte.
    \item \textbf{Conclusion :} C'est une \cle{fake news} par \textbf{détournement de contexte} (\emph{deep context} / \emph{misattribution}). L'image est vraie (réelle, non truquée) mais la légende (lieu, date, auteurs, bilan) est fausse. L'objectif est de créer la panique, de discréditer l'armée et d'alimenter le narratif séparatiste.
\end{enumerate}
\end{exoresolu}

\begin{attention}[Piège du « viral = vrai »]
\emph{Erreur fréquente :} Croire qu'une information est vraie parce qu'elle est « virale » (partagée des milliers de fois) ou parce qu'elle provient d'un « grand frère » ou d'un « contact de confiance ». La viralité mesure l'émotion (colère, peur, espoir, indignation), pas la vérité. Les algorithmes des réseaux sociaux (Facebook, X, TikTok) favorisent les contenus polémiques et clivants (engagement), pas les contenus vérifiés.
\end{attention}

\courssub{Les techniques d'influence médiatique : comment le message façonne la perception}

L'influence ne passe pas seulement par le mensonge, mais par la \textbf{construction du message} (\emph{framing}). Les professionnels de la communication (journalistes, communicants politiques, publicitaires) maîtrisent des techniques pour orienter la lecture du réel.

\begin{propriete}[Principales techniques de cadrage et d'influence médiatique]
\begin{itemize}
    \item \textbf{Répétition du message (Effet de vérité illusoire / \emph{Illusory Truth Effect}) :} Une affirmation répétée souvent (slogans politiques, \emph{talking points} repris en boucle par des porte-parole) finit par être perçue comme vraie, même si elle est fausse.
    \item \textbf{Choix des images et du son (Sémiologie visuelle/sonore) :} Un plan serré sur un visage en pleurs suscite l'empathie ; une foule filmée en contre-plongée semble puissante ; une musique angoissante transforme un reportage neutre en alerte ; ralenti = dramatisation.
    \item \textbf{Titres sensationnalistes (\emph{Clickbait} / \emph{Pute à clics}) :} « Choc : le ministre détourné des milliards ! » (alors qu'il s'agit d'une enquête préliminaire sans mise en examen). Le titre conditionne la lecture de l'article (biais d'ancrage).
    \item \textbf{Cadrage de l'information (\emph{Framing}) :} Présenter un même fait sous un angle différent pour orienter l'interprétation.
    \begin{center}
    \begin{tabularx}{\linewidth}{|X|X|}\hline
    \rowcolor{poporangeL}
    \textbf{Cadre A (Victimaire)} & \textbf{Cadre B (Responsabilité étatique)} \\ \hline
    « Grève des enseignants : les élèves otages » & « Enseignants en grève : le gouvernement refuse le dialogue depuis 6 mois » \\ \hline
    \end{tabularx}
    \end{center}
    \item \textbf{Fixation de l'agenda (\emph{Agenda-Setting}) :} Voir encadré théorique ci-dessous.
    \item \textbf{Amorçage (\emph{Priming}) :} En insistant sur certains thèmes (insécurité, corruption, prix du pain), les médias préparent l'opinion à juger l'action gouvernementale \textbf{principalement sur ces critères}, au détriment d'autres (éducation, santé, infrastructure).
\end{itemize}
\end{propriete}

\begin{definition}[Théorie de l'\emph{Agenda-Setting} (McCombs \& Shaw, 1972 -- Étude élection US 1968)]
Les médias ne disent pas aux gens \emph{comment} penser (direction du vote), mais \emph{à quoi} penser (salience des problèmes). En sélectionnant et en hiérarchisant les sujets (Une, temps d'antenne, fréquence de citation), ils imposent l'agenda des préoccupations publiques. \textbf{Corollaire :} Les sujets absents des médias deviennent inexistants dans l'opinion publique.
\end{definition}

\begin{savaistu}[L'\emph{Agenda-Setting} au Cameroun : repères pour le Bac]
Au Cameroun, l'agenda médiatique dominant est souvent structuré par :
\begin{enumerate}
    \item La communication présidentielle (discours du 31 déc, 10 mai, 20 mai, sorties exceptionnelles).
    \item L'actualité sécuritaire (crise NOSO, Extrême-Nord / Boko Haram).
    \item Les grandes infrastructures (stades CAN, barrages de Nachtigal/Memve'ele, routes, échangeurs).
    \item Les faits divers sensationnels (crimes, accidents, scandales people).
\end{enumerate}
Les sujets structurels (qualité de l'éducation, système de santé, emploi des jeunes, corruption systémique, changement climatique) peinent à s'imposer durablement à l'agenda médiatique dominant, sauf crise aiguë (grève, épidémie).
\end{savaistu}

\begin{popfigure}
\begin{tikzpicture}[
    node distance=1.5cm and 2cm,
    box/.style={rectangle, draw=popdark, thick, fill=popcream, rounded corners, minimum width=2.8cm, minimum height=1.2cm, align=center, font=\small},
    arrow/.style={-Stealth, thick, popdark},
    redarrow/.style={-Stealth, very thick, poporange},
    every node/.style={font=\small}
]
% Nœuds
\node[box, align=center] (source){Source malveillante\\(individu, bot, groupe politique)};
\node[box, right=of source, align=center] (reseaux){Réseaux sociaux\\(WhatsApp, Facebook, X/TikTok)\\\emph{Partage viral, algorithmique}};
\node[box, below=of reseaux, align=center] (medias){Médias classiques\\(Radio, TV, Presse en ligne)\\\emph{Reprise sans vérification (citizen journalism)}};
\node[box, right=of medias, align=center] (opinion){Opinion publique\\Citoyens, électeurs\\\emph{Perception déformée, peur, colère}};
\node[box, below=of opinion, align=center] (conseq){Conséquences sociales\\Panique, stigmatisation,\\violences, abstention électorale};
% Flèches principales
\draw[redarrow] (source) -- (reseaux) node[midway, above, fill=popcream, rounded corners, inner sep=2pt] {Création / Injection};
\draw[redarrow] (reseaux) -- (medias) node[midway, right, fill=popcream, rounded corners, inner sep=2pt] {Contamination};
\draw[redarrow] (medias) -- (opinion) node[midway, above, fill=popcream, rounded corners, inner sep=2pt] {Légitimation};
\draw[redarrow] (opinion) -- (conseq) node[midway, right, fill=popcream, rounded corners, inner sep=2pt] {Impact réel};
% Boucle de rétroaction
\draw[arrow, dashed, poppurple] (opinion) to[out=45, in=-45, looseness=1.5] node[midway, right, fill=popcream, rounded corners, inner sep=2pt, poppurple] {Retour : commentaires, nouveaux partages} (reseaux);
% Légendes étapes
\node[align=left, font=\tiny, popmuted] at (source.west |- conseq.south) {Étape 1 : Fabrication};
\node[align=left, font=\tiny, popmuted] at (reseaux.east |- conseq.south) {Étape 2 : Amplification exponentielle};
\node[align=left, font=\tiny, popmuted] at (medias.east |- conseq.south) {Étape 3 : Validation apparente};
\node[align=left, font=\tiny, popmuted] at (opinion.east |- conseq.south) {Étape 4 : Intériorisation};
\end{tikzpicture}
\legende{Schéma du circuit de propagation d'une \cle{fake news} : de la source à l'opinion publique via les réseaux sociaux et la reprise par les médias classiques. La boucle de rétroaction (pointillés) montre l'auto-alimentation du cycle.}
\end{popfigure}

\courssub{Vitesse de diffusion : rumeur vs information vérifiée}

La structure des réseaux sociaux (partage en un clic, algorithmes favorisant l'engagement émotionnel fort) donne un avantage structurel à la rumeur sur l'information vérifiée, qui demande du temps (enquête, recoupement, validation hiérarchique, droit de réponse).

\begin{popfigure}
\begin{tikzpicture}
\begin{axis}[
    width=\linewidth,
    height=8cm,
    xlabel={Temps (heures)},
    ylabel={Nombre de personnes atteintes (échelle logarithmique)},
    xmin=0, xmax=48,
    ymin=1, ymax=1000000,
    ymode=log,
    grid=major,
    grid style={popline},
    legend style={at={(0.02,0.98)}, anchor=north west, fill=popcream, draw=popdark, rounded corners, font=\small},
    tick label style={font=\small},
    label style={font=\small},
    domain=0:48,
    samples=100,
]
% Courbe Rumeur : croissance exponentielle rapide
\addplot[poporange, very thick, domain=0:24] {100*exp(0.35*x)}; 
\addlegendentry{Rumeur / Fake News (viralité non freinée)}
% Courbe Info vérifiée : démarrage lent (vérif), puis linéaire
\addplot[popblue, very thick, domain=0:48] {50*x + 100}; 
\addlegendentry{Information vérifiée (processus journalistique)}
% Seuil de "viralité massive"
\draw[popmuted, dashed] (0,10000) -- (48,10000) node[pos=0.5, above, font=\tiny, popmuted] {Seuil de viralité massive};
% Annotations
\node[poporange, font=\tiny, anchor=south west] at (axis cs:10, 3000) {Pic rumeur ~12h};
\node[popblue, font=\tiny, anchor=north west] at (axis cs:24, 1300) {Info vérifiée publiée};
\end{axis}
\end{tikzpicture}
\legende{Comparaison de la vitesse de diffusion : une rumeur (courbe orange, exponentielle) atteint des dizaines de milliers de personnes en quelques heures, tandis que l'information vérifiée (courbe bleue, linéaire après délai de vérification) met beaucoup plus de temps à toucher le même public. Échelle logarithmique sur l'axe vertical.}
\end{popfigure}

\begin{exemplebox}[Chiffres clés : la vitesse de la rumeur WhatsApp au Cameroun]
Au Cameroun, un message WhatsApp transféré dans 10 groupes de 200 membres chacun (2 000 personnes initiales) peut, si seulement 20\% le retransfèrent à 5 nouveaux groupes, atteindre \textbf{plus de 100 000 personnes en moins de 6 heures} (croissance géométrique). La correction (démenti officiel du MINAT, MINSANTE, MINDEF ou CNC) met souvent 24h à 48h à être diffusée sur les canaux officiels et ne suit pas les mêmes circuits de partage (moins d'émotion, moins de viralité, confiance moindre).
\end{exemplebox}

\courssub{Conséquences sociales de la désinformation non maîtrisée}

L'impact négatif des médias ne se limite pas à l'erreur factuelle ; il produit des effets réels, parfois irréversibles, sur le tissu social et la stabilité politique.

\begin{propriete}[Conséquences sociales majeures de la désinformation]
\begin{itemize}
    \item \textbf{Panique collective et comportements de ruée :} Ruées sur les stations-service (rumeur de pénurie de carburant), sur les banques/distributeurs (rumeur de faillite ou blocage des comptes), fuites de populations vers le village (rumeur d'attaque imminente).
    \item \textbf{Stigmatisation et discours de haine :} Désignation de boucs émissaires (ethnie, religion, origine géographique, corporation) accusés à tort de propager un virus, de commettre des crimes rituels, de voler des organes. Risque majeur de violences intercommunautaires (ex. rumeurs « kidnappeurs d'enfants » entraînant lynchages).
    \item \textbf{Polarisation de l'opinion (\emph{Polarisation affective}) :} La société se fracture en « camps » hermétiques au dialogue, chacun ne consommant que les médias qui confortent ses biais (\emph{bulles de filtre} / \emph{echo chambers}). Disparition du terrain d'entente.
    \item \textbf{Perte de confiance dans les institutions (Cynisme démocratique) :} Discrédit de la presse professionnelle (accusée de « complot », de « solde », de « presse aux ordres »), de la justice, de l'école, de l'armée. Conséquence : abstention électorale massive, désengagement civique, radicalisation.
    \item \textbf{Instabilité politique et menaces sur la paix :} Manipulation de l'opinion pour délégitimer un scrutin avant même le vote, provoquer une insurrection, ou justifier une répression préventive.
\end{itemize}
\end{propriete}

\begin{exemplebox}[Étude de cas : Couverture médiatique de la crise du NOSO (Nord-Ouest/Sud-Ouest)]
Depuis 2016, la crise anglophone a généré une « guerre de l'information » intense. Certains médias (locaux et internationaux) ont été accusés de \emph{framing} partial : soit en minimisant/occultant les exactions attribuées aux forces de défense, soit en occultant/minimisant celles des groupes armés séparatistes. Sur les réseaux sociaux, la diffusion de vidéos tronquées, de montages (deepfakes naissants) et de listes de « personnes à abattre » (enseignants, fonctionnaires, chefs traditionnels) sur WhatsApp a directement entraîné des assassinats ciblés et radicalisé les opinions des deux côtés, rendant le dialogue politique (Grand Dialogue National 2019) plus difficile.
\end{exemplebox}

\courssub{L'esprit critique comme remède : éduquer à l'information (EMI)}

Face à la pollution informationnelle, la régulation étatique (lois sur la cybercriminalité, la presse) est nécessaire mais insuffisante (juridiction territoriale, anonymat, volume). La réponse durable est l'\textbf{Éducation aux Médias et à l'Information (EMI)} pour former des citoyens \cle{résilients}, autonomes et responsables.

\begin{methode}[Vérifier une information en 4 étapes (Méthode S.A.D.R.)]
Appliquez systématiquement ce réflexe \textbf{avant} de croire, de commenter ou de partager :
\begin{enumerate}
    \item \textbf{S — Source :} Qui diffuse ? Est-ce un média reconnu (CRTV, AFP, \emph{Le Monde}, \emph{Jeune Afrique}, \emph{Data Cameroon}) ? Un site satirique (\emph{Le Gorafi}, \emph{Le Faso}) ? Un blog anonyme ? Un compte non vérifié (pas de badge bleu, peu d'abonnés, création récente) ? \rightarrow \emph{Action : Vérifier l'URL (orthographe exacte), la page « À propos », les mentions légales, le siège social.}
    \item \textbf{A — Auteur :} Qui a écrit/signé ? Un journaliste identifié (signature, bio, autres articles vérifiables) ? Un expert du domaine (titre, affiliation) ? Un témoin oculaire (géolocalisation cohérente) ? Ou un pseudonyme anonyme ? \rightarrow \emph{Action : Chercher le nom de l'auteur sur un moteur de recherche (Google Scholar, LinkedIn, site du média).}
    \item \textbf{D — Date et Contexte :} L'information est-elle actuelle ? Une vraie info de 2018 ressortie en 2025 devient une fake news par obsolescence (ex. nomination d'un ministre déjà remplacé). Le lieu correspond-il à l'image (météo, panneaux signalétiques, architecture, plaques d'immatriculation) ? \rightarrow \emph{Action : Vérifier la date de publication (pas seulement la date de partage), faire une recherche d'image inversée pour géolocaliser/dater.}
    \item \textbf{R — Recoupement (Cross-check) :} D'autres médias sérieux et indépendants la reprennent-ils ? Les sources primaires (communiqué officiel signé, procès-verbal, texte de loi, étude scientifique peer-reviewed, data.gov.cm) confirment-elles ? \rightarrow \emph{Action : Chercher les mots-clés sur Google News, AFP Factuel, \textbf{StopIntox (Cameroun)}, Data Cameroon, site officiel Gouvernement (spm.gov.cm, primature.cm).}
\end{enumerate}
\textbf{Règle d'or :} Si un seul doute persiste après S.A.D.R. : \textbf{NE PARTAGEZ PAS}. Signalez le contenu (fonction « Signaler » sur la plateforme).
\end{methode}

\begin{aretenir}[Outils de vérification accessibles au Cameroun]
\begin{itemize}
    \item \textbf{StopIntox (stopintox.cm) :} Plateforme camerounaise de \emph{fact-checking} (vérification des faits), membre de l'IFCN (International Fact-Checking Network). Publie des rapports réguliers.
    \item \textbf{Data Cameroon (datacameroon.com) :} Portail officiel des données ouvertes du gouvernement (budgets, statistiques, appels d'offres).
    \item \textbf{Google Reverse Image Search / TinEye / Yandex Images :} Pour tracer l'origine, la date et le lieu réel d'une photo ou vidéo.
    \item \textbf{InVID / WeVerify (extensions navigateur gratuites) :} Analyse vidéo avancée (keyframes, métadonnées EXIF, détection de montage).
    \item \textbf{Sites officiels :} \texttt{spm.gov.cm} (Services du Premier Ministre), \texttt{primature.cm}, \texttt{presidence.cm}, sites des ministères (minsante.gov.cm, minat.gov.cm, mindef.gov.cm) pour les communiqués authentiques.
\end{itemize}
\end{aretenir}

\begin{experience}[Atelier pratique : « Chasse aux infox » (à réaliser en classe ou en TD - 1h30)]
\textbf{Matériel :} Smartphones ou ordinateurs (par binôme), connexion internet, fiche d'analyse S.A.D.R. fournie par le professeur (papier ou numérique).
\textbf{Protocole :}
\begin{enumerate}
    \item Le professeur distribue 5 captures d'écran récentes (WhatsApp, Facebook, X) préparées à l'avance : 2 vraies infos vérifiées, 2 \cle{fake news} avérées (démenties par StopIntox ou source officielle), 1 article satirique (ex. page parodique camerounaise).
    \item Les élèves appliquent la méthode S.A.D.R. de manière autonome (15 min). Ils notent les preuves (URL, date, source primaire).
    \item Restitution collective : chaque binôme présente son verdict (Vrai / Faux / Satire / Non vérifiable) et \textbf{la preuve décisive} (lien source officiel, date de l'image originale, incohérence interne).
    \item Débat guidé : Pourquoi la fausse info semblait plausible ? (Mélange vrai/faux, émotion forte, autorité usurpée, contexte réel détourné, biais de confirmation).
\end{enumerate}
\textbf{Observation attendue :} Les fausses infos les plus crédibles mélangent habilement du vrai (nom réel d'un ministre, lieu réel, date récente, logo officiel volé) et du faux (citation inventée, chiffre falsifié, mesure inexistante).
\textbf{Interprétation pédagogique :} La vigilance porte sur le \textbf{détail vérifiable} (chiffre exact, citation sourcée, lien vers source primaire), pas sur l'enveloppe globale (mise en page, ton, logo).
\end{experience}

\begin{exercice}[Entraînement à l'esprit critique : Cas pratique WhatsApp]
\textbf{Consigne :} Vous recevez ce message transféré 5 fois sur WhatsApp (icône double flèche) : 
\begin{quote}
\textbf{URGENT : MINSANTE. À partir de demain, vaccination obligatoire porte-à-porte contre Ebola. Les agents porteront des blouses bleues. Refus = amende 50 000 FCFA. Partagez à tous vos groupes pour sauver des vies !}
\end{quote}
\begin{enumerate}
    \item Appliquez la méthode S.A.D.R. pour analyser ce message.
    \item Identifiez 3 indices de manipulation (langage, forme, contenu).
    \item Quelle attitude citoyenne adopter ? Justifiez en deux phrases.
\end{enumerate}
\end{exercice}

\begin{corrige}[Exercice n°1 : Analyse message WhatsApp « Ebola »]
\begin{enumerate}
    \item \textbf{S (Source) :} WhatsApp (messagerie chiffrée, anonyme, non traçable à l'origine). Pas de lien vers site officiel. \textbf{A (Auteur) :} Inconnu (transféré). \textbf{D (Date/Contexte) :} « À partir de demain » (flou, pas de date précise, pas de numéro de communiqué). Contexte : Pas d'épidémie d'Ebola déclarée au Cameroun à ce jour (vérification OMS/MINSANTE). \textbf{R (Recoupement) :} Aucun communiqué sur \texttt{minsante.gov.cm}, \texttt{spm.gov.cm}, CRTV, StopIntox. La vaccination porte-à-porte obligatoire avec amende n'existe pas dans le cadre légal camerounais (Loi sur la santé publique).
    \item \textbf{Indices de manipulation :} (1) \textbf{Urgence majuscule + injonction de partage} (pression psychologique, court-circuite la réflexion). (2) \textbf{Menace d'amende précise (50 000 FCFA)} pour faire peur et obtenir conformité. (3) \textbf{Absence de traçabilité officielle :} Pas de numéro de téléphone vert, pas de lien vers le site, pas de nom du Ministre, pas de référence de texte de loi.
    \item \textbf{Attitude citoyenne :} \textbf{Ne pas partager.} Signaler le message comme « Fake News / Spam » dans WhatsApp. Informer poliment l'expéditeur que c'est une fausse info en lui envoyant le lien du démenti de StopIntox ou du site MINSANTE. Consulter le site officiel du MINSANTE pour l'information sanitaire réelle.
\end{enumerate}
\end{corrige}

\begin{attention}[Écran de fumée : la satire prise au sérieux]
Des sites comme \emph{Le Gorafi} (France) ou des pages parodiques camerounaises (ex. \emph{Le Camerounais Satirique}, \emph{Actu Kmer} parodie) publient des articles humoristiques. Partagés hors contexte (capture d'écran sans l'URL, sans la mention « Satire »), ils sont crus par des milliers de personnes. \textbf{Règle absolue :} Vérifier toujours l'URL exacte et la rubrique « À propos / Mentions légales » du site source.
\end{attention}

\courssub{Synthèse : l'opinion publique à l'épreuve du numérique}

\begin{aretenir}[Bilan de la section : L'impact des médias sur l'opinion publique]
\begin{itemize}
    \item Les médias ont un \textbf{double pouvoir} : éclairer (information, éducation, sensibilisation, \emph{watchdog}) ou obscurcir (désinformation, propagande, \cle{fake news}, rumeur).
    \item L'influence passe par des \textbf{techniques professionnelles identifiables} (\emph{agenda-setting}, \emph{framing}, \emph{priming}, répétition, sensationnalisme, sémiologie visuelle) qu'il faut savoir décoder.
    \item La \textbf{vitesse de la rumeur} (exponentielle sur réseaux sociaux, amplifiée par l'algorithme) dépasse largement celle de la vérité (linéaire, coûteuse en vérification humaine).
    \item Les \textbf{conséquences sociales} sont réelles et graves : panique, haine, polarisation, défiance démocratique, violences physiques, instabilité.
    \item Le \textbf{rempart efficace} est l'esprit critique outillé : méthode \textbf{S.A.D.R.} (Source, Auteur, Date, Recoupement), usage réflexe des plateformes de \emph{fact-checking} (StopIntox), refus catégorique de partager sans vérifier.
\end{itemize}
\end{aretenir}

\coursec{La protection de l'opinion publique au Cameroun}

\courssub{Cadre juridique national : les textes fondateurs}

La protection de l'opinion publique repose sur l'équilibre entre \textbf{liberté d'expression} (constitutionnelle) et \textbf{responsabilité} (pénale, civile, administrative).

\begin{propriete}[Arsenal juridique camerounais de régulation médiatique]
\begin{itemize}
    \item \textbf{Constitution du 18 janvier 1996 (Préambule) :} Garantit « la liberté de communication, d'expression, de la presse » (art. 1er). Liberté soumise au respect de la loi, des droits d'autrui, de l'ordre public.
    \item \textbf{Loi n°90/052 du 19 décembre 1990 relative à la communication sociale (modifiée par Loi 96/04 et 2012/001) :} Texte fondateur. Définit le journaliste, le média, le droit à l'information, le secret des sources, le droit de réponse, la dépénalisation des délits de presse (amendes, pas de prison pour diffamation simple), création du CNC.
    \item \textbf{Loi n°2010/012 du 21 décembre 2010 relative à la cybercriminalité et la cyber-sécurité :} Sanctionne la propagation de fausses nouvelles (\cle{fake news}) via systèmes informatiques (Art. 78 : peine de 6 mois à 2 ans de prison et/ou amende 100 000 à 5 000 000 FCFA). Sanctionne aussi l'atteinte à l'intimité, la pédopornographie, le terrorisme numérique.
    \item \textbf{Loi n°2019/020 du 24 décembre 2019 (Code Pénal) :} Art. 113 (Diffamation), 114 (Injure), 115 (Outrage), 241 (Propagation de fausses nouvelles troublant l'ordre public). \textbf{Note :} La prison est rétablie pour certains délits de presse via le Code Pénal, créant une tension avec la Loi 90/052.
    \item \textbf{Loi n°2015/018 du 21 juillet 2015 relative à la lutte contre le terrorisme :} Sanctionne l'apologie du terrorisme et la diffusion de propagande terroriste via médias/réseaux sociaux.
    \item \textbf{Décret n°2012/001 du 10 janvier 2012 :} Organisation et fonctionnement du CNC.
\end{itemize}
\end{propriete}

\courssub{Institutions de régulation, de contrôle et de médiation}

\begin{center}
\begin{tabularx}{\linewidth}{|l|X|X|}\hline
\rowcolor{popblueL}
\textbf{Institution} & \textbf{Missions principales} & \textbf{Pouvoirs de sanction} \\ \hline
\textbf{CNC (Conseil National de la Communication)} & Régulation du secteur (attribution cartes de presse, agrément médias), veille déontologique, médiation (droit de réponse), avis sur projets de lois, observation élections. & Avertissement, blâme, suspension temporaire (max 3 mois), retrait carte de presse, suspension parution, amendes administratives. \\ \hline
\textbf{MINCOM (Ministère de la Communication)} & Tutelle administrative, élaboration politique sectorielle, subventions presse, gestion fonds d'aide. & Pouvoir de suspension d'urgence (art. 37 Loi 90/052) pour atteinte à l'ordre public / sécurité nationale (décision motivée, notifiée CNC). \\ \hline
\textbf{ANTIC (Agence Nationale des Technologies de l'Information et de la Communication)} & Sécurité des systèmes d'information, lutte cybercriminalité (CERT-CM), certification, sensibilisation. & Expertise technique, alerte, transmission au Procureur / Forces de l'ordre. Pas de pouvoir de sanction médiatique direct. \\ \hline
\textbf{Justice (Tribunaux, Parquets)} & Répression pénale (diffamation, fausses nouvelles, cybercriminalité, terrorisme), réparation civile (dommages-intérêts). & Peines de prison, amendes, dommages-intérêts, interdiction d'exercer, confiscation matériel. \\ \hline
\end{tabularx}
\end{center}

\begin{aretenir}[Le CNC : Organe indépendant ?]
Le CNC est une \textbf{institution administrative indépendante} (Loi 2012/001). Ses 11 membres (dont 4 élus par les journalistes, 3 par les éditeurs, 4 nommés par Président/PM/Présidents Assemblée/Sénat) prêtent serment. Il gère le Fonds d'Aide à la Presse Privée. Son avis est consultatif mais obligatoire pour certains actes. Il est le garant de l'éthique (Charte de déontologie du journaliste camerounais, 1997).
\end{aretenir}

\courssub{Mécanismes de protection concrets pour le citoyen}

\begin{methode}[Comment se protéger / faire valoir ses droits face à un média ?]
\begin{enumerate}
    \item \textbf{Droit de réponse (Loi 90/052, Art. 19-20) :} Toute personne physique/morale mise en cause a droit de réponse \textbf{gratuit}, dans les mêmes conditions (emplacement, durée, police) que l'article initial. Délai : 3 mois (presse écrite), 1 mois (audiovisuel/numérique). Refus = saisine CNC puis Tribunal.
    \item \textbf{Saisine du CNC (Plainte déontologique) :} Tout citoyen peut saisir le CNC pour manquement à la déontologie (information inexacte, partialité, atteinte vie privée, discrimination). Gratuit, non contentieux. CNC enquête, auditionne, rend avis/ sanction administrative.
    \item \textbf{Action en Justice (Voie pénale/civile) :} Citation directe ou plainte au Procureur (diffamation, injure, fausses nouvelles, cyberharcèlement). Demande de réparation (dommages-intérêts). Délai de prescription court (3 mois pour diffamation/presse).
    \item \textbf{Signalement plateforme / ANTIC :} Pour contenus illicites en ligne (haine, terrorisme, pédopornographie, fausses nouvelles virales) : signalement intégré (Facebook, X, TikTok, WhatsApp) + signalement ANTIC (\texttt{cert.cm} / \texttt{antic.cm}) + plainte Police/Gendarmerie.
\end{enumerate}
\end{methode}

\begin{exemplebox}[Cas pratique : Droit de réponse au Cameroun]
Un quotidien privé publie en Une : « Détournement à la Mairie de X : le Maire Y épinglé par le CONSUPE ». Le Maire Y dispose de 3 mois pour exercer son droit de réponse. Il rédige un texte (max 1/2 page), l'envoie en recommandé AR au Directeur de Publication. Le journal \textbf{doit} le publier dans le numéro suivant (même emplacement, même police) sans modification ni commentaire. Si refus : Maire Y saisit CNC (urgence) $\rightarrow$ CNC ordonne publication sous astreinte (amende/jour de retard) $\rightarrow$ Si persistance : Tribunal de Grande Instance (ordonnance de publication forcée + dommages-intérêts).
\end{exemplebox}

\courssub{Enjeux actuels de la régulation au Cameroun}

\begin{attention}[Défis non résolus pour le Bac]
\begin{itemize}
    \item \textbf{Conflit de lois :} Loi 90/052 (dépénalisation presse) vs Code Pénal 2019 (prison pour diffamation/fausses nouvelles). Insécurité juridique pour journalistes.
    \item \textbf{Régulation du numérique :} La Loi 90/052 est pré-internet. Le CNC peine à réguler les pure-players en ligne, les réseaux sociaux (transnationaux), les influenceurs. Projet de loi sur la régulation des réseaux sociaux en discussion.
    \item \textbf{Indépendance effective du CNC :} Nomination politique de 4/11 membres, budget dépendant de l'État, difficultés à sanctionner les médias pro-pouvoir.
    \item \textbf{Protection des sources :} Menacée par les réquisitions judiciaires (procureur) sur les serveurs/téléphones des journalistes (affaires récentes).
    \item \textbf{Éducation aux médias (EMI) :} Absente des programmes scolaires obligatoires (hors ECM optionnel), laissant les jeunes sans défense face aux algorithmes.
\end{itemize}
\end{attention}

\coursec{TD3 : La publicité dans les médias}

\courssub{La publicité : définition, fonctions et formes}

\begin{definition}[Publicité commerciale]
La \cle{publicité} est une communication de masse, payante, non personnelle, identifiant clairement son annonceur (marque, entreprise, institution), dont le but est d'influencer les attitudes et/ou les comportements d'une audience cible (consommateurs, électeurs, usagers) en faveur d'un produit, service, idée ou image.
\end{definition}

\begin{propriete}[Fonctions de la publicité]
\begin{enumerate}
    \item \textbf{Informationnelle :} Faire connaître l'existence, les caractéristiques, le prix, la disponibilité (nouveau produit, promotion).
    \item \textbf{Persuasive / Incitative :} Créer la préférence, le désir, l'intention d'achat (différenciation, image de marque, valeurs).
    \item \textbf{Mécanique / Rappel :} Maintenir la notoriété, fidéliser, déclencher l'achat d'impulsion (slogan, jingle, logo omniprésent).
    \item \textbf{Économique :} Financer les médias (modèle économique presse/TV gratuite), soutenir la concurrence, stimuler l'innovation.
    \item \textbf{Sociale / Culturelle :} Véhiculer des modèles de vie, des stéréotypes (genre, réussite, corps), créer des rites partagés (Noël, rentrée, CAN).
\end{enumerate}
\end{propriete}

\begin{center}
\begin{tabularx}{\linewidth}{|l|X|X|}\hline
\rowcolor{popblueL}
\textbf{Type de publicité} & \textbf{Caractéristiques} & \textbf{Exemples médias camerounais} \\ \hline
\textbf{Médias classiques (Offline)} & TV (spot 30s, parrainage), Radio (spot, chronique), Presse (encart, publireportage), Affichage (panneaux, bus), Cinéma. & Spot MTN/Orange sur CRTV/Equinoxe, Panneaux Brasseries/Guinness aux carrefours Yaoundé/Douala. \\ \hline
\textbf{Médias numériques (Online)} & Display (bannières), Search (Google Ads), Social Ads (Facebook/Insta/TikTok Ads), Influenceurs (posts sponsorisés, \#ad), Native Ads (contenu sponsorisé ressemblant à article), Emailing/SMS. & Campagnes Jumia/Konga (Black Friday), Influenceurs beauté (produits capillaires), SMS promotionnels banques (UBA, BICEC, Afriland). \\ \hline
\textbf{Publicité institutionnelle / d'intérêt général} & Annonceur = État, ONG, Collectivité. But = comportement civique/santé/éducation. Non commerciale. & Campagnes MINSANTE (vaccination), MINAT (recensement), ELECAM (inscription), ONG (scolarisation filles). \\ \hline
\textbf{Publireportage / Contenu de marque (\emph{Brand Content})} & Format éditorial (article, reportage, vidéo) payé par l'annonceur. Doit être signalé « Communication commerciale » / « Sponsorisé ». & Supplément spécial « 20 mai » dans \emph{Cameroon Tribune} payé par Mincom ; Reportage « Visite usine SOSUCAM » sur CRTV payé par l'entreprise. \\ \hline
\end{tabularx}
\end{center}

\courssub{Techniques publicitaires et influence sur le consommateur}

\begin{propriete}[Principaux leviers psychologiques utilisés]
\begin{itemize}
    \item \textbf{Appel à l'émotion :} Peur (assurance, sécurité), Joie/Bonheur (boissons, téléphonie), Nostalgie (produits traditionnels), Fierté/Identité (marques « Made in Cameroon », CAN), Sexe/Séduction (parfums, mode).
    \item \textbf{Preuve sociale / Témoignage :} « 9 femmes sur 10 recommandent », « Le choix des champions », influenceurs « comme vous ».
    \item \textbf{Autorité / Expertise :} Blouse blanche (dentifrice, médicament OTC), Ingénieur (ciment, téléphone), Chef traditionnel (produit terroir).
    \item \textbf{Rareté / Urgence :} « Offre limitée », « Stocks épuisés bientôt », « Promo 48h seulement ».
    \item \textbf{Répétition / Conditionnement :} Slogan + Jingle + Logo répétés jusqu'à l'automatisme (ex. « MTN \emph{Everywhere you go} », « Guinness \emph{Greatness} »).
    \item \textbf{Association d'idées (Transfert d'affect) :} Produit associé à valeurs positives : Nature (eau minérale), Famille (lait, huile), Modernité/Tech (smartphone), Tradition (bière, épices).
\end{itemize}
\end{propriete}

\begin{aretenir}[Publicité trompeuse et agressive (Interdites - Loi 90/052 Art. 34, Code Pénal, Loi Consommation)]
\begin{itemize}
    \item \textbf{Trompeuse :} Allégations fausses sur caractéristiques (composition, origine, quantité, prix, résultats), omissions essentielles, fausses promesses (guérison miracle, enrichissement rapide).
    \item \textbf{Agressive :} Harcèlement (appels/SMS répétés non sollicités), exploitation de la crédulité/faiblesse (personnes âgées, mineurs), pression physique/psychologique.
    \item \textbf{Comparative :} Autorisée si objective, vérifiable, porte sur caractéristiques essentielles, pas dévalorisante, pas confusion marque.
\end{itemize}
\end{aretenir}

\courssub{Régulation de la publicité au Cameroun}

\begin{center}
\begin{tabularx}{\linewidth}{|l|X|}\hline
\rowcolor{popblueL}
\textbf{Acteur} & \textbf{Rôle} \\ \hline
\textbf{CNC (Conseil National de la Communication)} & Avis sur éthique publicitaire (Charte déontologie), médiation plaintes publicités trompeuses/choquantes, surveillance contenu (protection mineurs, valeurs morales). \\ \hline
\textbf{MINCOM / MINCOMMERCE} & Tutelle, répression fraudes (contrôle prix affichés vs publicités), labels qualité. \\ \hline
\textbf{ARMP (Agence de Régulation des Marchés Publics)} & Contrôle publicité appels d'offres publics (transparence). \\ \hline
\textbf{ANOR (Agence des Normes et de la Qualité)} & Validation allégations techniques / normes (ex. « Norme ISO », « Certifié ANOR »). \\ \hline
\textbf{APC (Association de Protection des Consommateurs) / Ligue Camerounaise des Droits du Consommateur (LCDC)} & Plaintes, actions de groupe, sensibilisation, représentation au CNC. \\ \hline
\textbf{Autorégulation professionnelle (GAP - Groupement des Annonceurs du Cameroun, AACC - Assoc. Agences Conseil)} & Code de déontologie professionnel, pré-conseil, jury déontologique interne (rarement contraignant). \\ \hline
\end{tabularx}
\end{center}

\begin{savaistu}[La publicité pour l'alcool et le tabac au Cameroun]
Loi 90/052 (Art. 35) \& Décret d'application : Interdiction de la publicité pour le tabac (totale). Publicité pour l'alcool : autorisée mais encadrée (pas de mineurs, pas de lien performance sportive/conduite, horaires protégés TV/Radio : 22h-6h, messages sanitaires obligatoires « L'abus d'alcool est dangereux pour la santé »). En pratique : contournement par sponsoring sportif/culturel (brasseries = sponsors majeurs CAN, Coupe du Roi, festivals), placement de produit dans clips musicaux, influenceurs.
\end{savaistu}

\courssub{Exercice d'application : Analyse critique d'une campagne publicitaire}

\begin{exercice}[TD3 : Analyse d'une publicité télévisée / réseau social (Durée : 1h)]
\textbf{Consigne :} Choisissez une publicité récente (TV, Facebook, TikTok, panneau) d'une grande marque au Cameroun (Brasseries, MTN/Orange, Banque, Savonnerie, Télécom, Alimentaire). Répondez aux questions ci-dessous sur une fiche structurée.
\begin{enumerate}
    \item \textbf{Identification :} Marque, Produit, Support, Date/Contexte de diffusion, Cible apparente (âge, sexe, milieu, langue).
    \item \textbf{Analyse du message (Sémiologie) :} 
    \begin{itemize}
        \item Visuel : Couleurs, plans, personnages (rôle, tenue, attitude), décor, texte/image ratio.
        \item Sonore : Musique (style, tempo), Voix off (ton, langue), Slogan/Jingle.
        \item Narratif : Histoire racontée (problème $\rightarrow$ solution produit), valeurs véhiculées.
    \end{itemize}
    \item \textbf{Techniques d'influence identifiées :} (Cocher + justifier) Émotion / Preuve sociale / Autorité / Rareté / Répétition / Transfert d'affect / Humour / Choc.
    \item \textbf{Vérification éthique et légale :} 
    \begin{itemize}
        \item Y a-t-il allégation vérifiable (prix, composition, résultat) ? Est-elle sourcée ?
        \item Respect de la dignité humaine (pas de stéréotypes sexistes, tribaux, handicap) ?
        \item Protection des mineurs (horaire, contenu) ?
        \item Mention « Communication commerciale » / « Sponsorisé » si influenceur/native ads ?
        \item Message sanitaire obligatoire si alcool/tabac/médicament ?
    \end{itemize}
    \item \textbf{Esprit critique citoyen :} Cette pub vous semble-t-elle \textbf{responsable} ? Pourquoi ? Quel impact réel sur l'opinion/le comportement ? Proposez une amélioration.
\end{enumerate}
\textbf{Rendu :} Exposé oral 5 min / binôme + fiche écrite notée.
\end{exercice}

\begin{corrige}[Grille d'évaluation indicative TD3 (Professeur)]
\begin{center}
\begin{tabularx}{\linewidth}{|l|X|c|}\hline
\rowcolor{popblueL}
\textbf{Critères} & \textbf{Indicateurs de réussite} & \textbf{Points} \\ \hline
Identification complète & Marque, support, date, cible précise (pas « tout le monde ») & 2 \\ \hline
Analyse sémiologique & Vocabulaire précis (plan, cadrage, symboles, code couleur, langue), lien forme/sens & 4 \\ \hline
Techniques d'influence & 3 techniques min. nommées + justifiées par éléments concrets de la pub & 4 \\ \hline
Vérification éthique/légale & Référence aux textes (Loi 90/052, Code Conso, Décret alcool), détection stéréotypes, mentions obligatoires & 5 \\ \hline
Esprit critique / Proposition & Argumentation personnelle étayée, distinction séduction/manipulation, proposition constructive & 5 \\ \hline
\textbf{Total} & & \textbf{20} \\ \hline
\end{tabularx}
\end{center}
\end{corrige}

\begin{exemplebox}[Exemple d'analyse type : Campagne « MTN MoMo - \emph{Send Money Home} » (Fête du Travail / Fête Nationale)]
\begin{itemize}
    \item \textbf{Cible :} Diaspora urbaine, jeunes actifs, commerçants transfrontaliers (Nigeria, Gabon, Guinée Équatoriale).
    \item \textbf{Visuel :} Split screen : Fils à Douala (smartphone, sourire) $\leftrightarrow$ Mère au village (réception argent, achat riz/médicaments, joie). Couleurs MTN (Jaune/Noir) dominantes.
    \item \textbf{Sonore :} Pidgin English + Français + Langue locale (Bassa/Ewondo). Jingle entraînant. Slogan : « MoMo, c'est la vie facile ».
    \item \textbf{Techniques :} \textbf{Émotion} (lien familial, devoir filial), \textbf{Preuve sociale} (« Tout le monde utilise MoMo »), \textbf{Transfert d'affect} (Technologie moderne = soin traditionnels parents), \textbf{Répétition} (slogan/jingle).
    \item \textbf{Éthique :} Positif (inclusion financière, entraide familiale). Attention : frais de transaction non mentionnés dans le spot (info incomplète), promotion usage data/smartphone (exclut non-connectés).
    \item \textbf{Amélioration :} Afficher tarif indicatif « Ex: 10 000 FCFA envoyés = X FCFA reçus ». Version USSD (*126\#) mise en avant pour téléphones basiques.
\end{itemize}
\end{exemplebox}

\begin{aretenir}[Synthèse TD3 : Publicité et Citoyenneté]
\begin{itemize}
    \item La publicité est un \textbf{pouvoir d'influence massif} (financement médias, modelage désirs, normes sociales).
    \item Le citoyen averti \textbf{décrypte} : il distingue l'argument rationnel (prix, composant) de la séduction émotionnelle (image, rêve).
    \item Il connaît ses \textbf{droits} : signaler pub trompeuse/agressive (CNC, LCDC, MINCOMMERCE), exiger mention « Sponsorisé ».
    \item Il agit en \textbf{consommateur responsable} : compare, vérifie allégations (ANOR, labels), résiste à l'achat compulsif (rareté artificielle), privilégie l'info objective (tests comparatifs, avis vérifiés).
\end{itemize}
\end{aretenir}

\coursec{La protection de l'opinion publique au Cameroun}

Après avoir analysé \emph{comment} les médias influencent l'opinion publique (agenda-setting, priming, framing, spirale du silence), une question s'impose : \textbf{qui veille à ce que ce pouvoir d'influence ne devienne pas un pouvoir de manipulation ?} Dans une démocratie, l'opinion publique n'est pas une matière inerte que l'on modèle à volonté ; elle est le socle de la souveraineté nationale. La protéger, c'est protéger la paix sociale, la dignité des citoyens et le jeu démocratique lui-même. Le Cameroun, conscient de ces enjeux, a mis en place un arsenal juridique et institutionnel spécifique que nous allons décortiquer pas à pas.

\courssub{Pourquoi protéger l'opinion publique ? Les dangers de la manipulation}

L'opinion publique est vulnérable. Lorsque l'information est tronquée, faussée ou orientée dans un but partisan, commercial ou idéologique, les conséquences sont réelles et graves.

\begin{aretenir}[Les risques d'une opinion non protégée]
\begin{itemize}
    \item \textbf{Démantèlement du lien social :} la diffusion de rumeurs ethniques ou religieuses (ex. : messages de haine sur WhatsApp ou Facebook) peut déclencher des violences intercommunautaires.
    \item \textbf{Détournement du processus électoral :} la désinformation massive (\emph{fake news}) pendant une campagne fausse le choix du citoyen et délégitime les institutions issues du scrutin.
    \item \textbf{Atteinte à la dignité humaine :} calomnies, diffamation, exposition de la vie privée sans consentement détruisent des réputations et des vies.
    \item \textbf{Perte de confiance :} si les citoyens ne croient plus aucun média, c'est l'espace public lui-même qui se vide, laissant place au complotisme et au repli identitaire.
\end{itemize}
\end{aretenir}

\emph{Exemple concret :} en 2019-2020, la circulation de fausses informations sur des prétendus enlèvements d'enfants dans certaines localités du Cameroun a provoqué des lynchages d'innocents. La rumeur, amplifiée par les réseaux sociaux, a tué plus vite que la vérité n'a pu la rattraper. C'est pour éviter de tels drames que la protection juridique de l'opinion publique existe.

\courssub{Le cadre juridique : la Loi n° 90/052 du 19 décembre 1990 et ses amendements}

C'est la \textbf{loi fondamentale} régissant la communication sociale au Cameroun. Elle pose le principe de la liberté de la presse tout en fixant des garde-fous.

\begin{definition}[Loi n° 90/052 du 19 décembre 1990]
Texte fondateur portant régime de la communication sociale au Cameroun. Il libéralise la presse écrite et audiovisuelle, définit les conditions d'accès à la profession de journaliste, les droits et devoirs des organes de presse, et les infractions spécifiques (diffamation, injure, fausses nouvelles).
\end{definition}

\textbf{Principaux apports et évolutions :}
\begin{itemize}
    \item \textbf{Libéralisation :} fin du monopole d'État (CRTV) ; naissance de la presse privée écrite (1990) puis audiovisuelle (décret n° 2000/158 du 3 avril 2000).
    \item \textbf{Aménagement des délits de presse :} les réformes successives ont aménagé le régime des peines (amendes alourdies), tout en maintenant des peines privatives de liberté pour les infractions les plus graves (incitation à la haine, à la violence).
    \item \textbf{Adaptation au numérique :} la loi n° 2010/012 du 21 décembre 2010 relative à la cybersécurité et à la cybercriminalité étend la répression des infractions de presse aux médias en ligne et aux réseaux sociaux.
\end{itemize}

\begin{attention}[Liberté de la presse $\neq$ liberté de tout dire]
La liberté d'expression (art. 19 de la Déclaration universelle des droits de l'homme ; Préambule de la Constitution camerounaise) n'est pas absolue. Elle s'arrête là où commencent les droits d'autrui et l'ordre public. La loi camerounaise sanctionne notamment :
\begin{itemize}
    \item \textbf{La diffamation :} allégation ou imputation d'un fait précis qui porte atteinte à l'honneur (ex. : « M. X a détourné 50 millions » sans preuve).
    \item \textbf{L'injure :} expression outrageante, terme de mépris qui ne renferme l'imputation d'aucun fait précis (ex. : « M. X est un voleur »).
    \item \textbf{L'atteinte à la vie privée :} publication de photos, correspondances ou données de santé sans consentement.
    \item \textbf{L'incitation à la haine ou à la violence :} appel au meurtre, à la discrimination ethnique, religieuse ou régionale.
    \item \textbf{La propagation de fausses nouvelles :} information mensongère publiée de mauvaise foi et troublant la paix publique.
\end{itemize}
\textbf{Piège d'examen :} ne confondez pas \emph{critique légitime} (dénoncer une gestion opaque avec des preuves) et \emph{diffamation} (accuser sans preuve). La bonne foi et la base factuelle sont les meilleures défenses du journaliste poursuivi.
\end{attention}

\courssub{Les institutions de régulation : architecture institutionnelle}

La protection ne repose pas seulement sur des textes, mais sur des organes chargés de les faire appliquer.

\begin{popfigure}
\begin{tikzpicture}[
    node distance=1.2cm and 2cm,
    box/.style={draw=popblue, thick, rounded corners, fill=popblue!10, minimum width=3.5cm, minimum height=1.2cm, align=center, font=\small},
    inst/.style={draw=popdark, thick, rounded corners, fill=popcream, minimum width=3.5cm, minimum height=1cm, align=center, font=\small\bfseries},
    arrow/.style={-Stealth, thick, popdark},
    link/.style={-Stealth, dashed, poporange}
]
\node[inst, align=center] (mincom){MINCOM\\Ministère de la Communication\\ (Tutelle, politique sectorielle)};
\node[box, below=of mincom, align=center] (cnc){CNC\\Conseil National de la Communication\\ \textbf{Organe de régulation}\\(Éthique, déontologie, sanctions)};
\node[inst, below left=of cnc, xshift=-1.5cm, align=center] (justice){Justice\\Tribunaux\\ \textbf{Répression pénale}\\(Diffamation, haine, cybercriminalité)};
\node[inst, below right=of cnc, xshift=1.5cm, align=center] (autres){Autres acteurs\\ART (régulation des télécoms)\\OSC et observatoires des médias\\Plateformes de signalement};
\draw[arrow] (mincom) -- node[left, font=\tiny, pos=0.5] {Tutelle administrative} (cnc);
\draw[arrow] (cnc) -- node[left, font=\tiny, pos=0.5] {Saisine / Signalement} (justice);
\draw[link] (cnc) -- node[above, font=\tiny, pos=0.5] {Collab.} (autres);
\draw[link] (justice) -- node[below, font=\tiny, pos=0.5] {Application des peines} (autres);
\node[draw=popline, fill=white, rounded corners, font=\tiny, align=left, anchor=north east] at (current bounding box.north east) {
    \textbf{Légende :}\\
    Trait plein = hiérarchie / saisine directe\\
    Trait pointillé = coordination / collaboration
};
\end{tikzpicture}
\legende{Organigramme simplifié des institutions de protection de l'opinion publique au Cameroun. Le CNC est le pivot de la régulation déontologique ; la justice sanctionne les infractions pénales ; le MINCOM définit la politique publique de communication.}
\end{popfigure}

\courssub{Le Conseil National de la Communication (CNC) : le gardien de l'éthique}

Créé par la loi n° 90/052 de 1990, le CNC est une \textbf{autorité administrative de régulation}. Ses membres (magistrats, journalistes, personnalités qualifiées) sont nommés par décret pour un mandat renouvelable.

\begin{propriete}[Missions principales du CNC]
\begin{enumerate}
    \item \textbf{Veiller au respect de l'éthique et de la déontologie} dans la pratique du journalisme et de la communication.
    \item \textbf{Assurer l'équité d'accès aux médias publics :} temps de parole équitable pour les formations politiques, surtout en période électorale.
    \item \textbf{Sanctionner les dérives :} pouvoir de sanction administrative (avertissement, blâme, suspension d'émission ou de parution, interdiction temporaire d'exercer).
    \item \textbf{Donner des avis consultatifs} sur les projets de textes relatifs à la communication et sur l'attribution des fréquences.
    \item \textbf{Protéger le public :} traitement des plaintes des citoyens (droit de réponse, protection des mineurs).
\end{enumerate}
\end{propriete}

\begin{definition}[Déontologie journalistique]
Ensemble de règles morales et professionnelles qui s'imposent au journaliste dans l'exercice de sa fonction : \textbf{vérité} (vérification des faits), \textbf{indépendance} (refus des pressions politiques ou financières), \textbf{équité} (droit de réponse, contradictoire), \textbf{respect de la vie privée}, \textbf{protection des sources}, \textbf{refus de la corruption} (le fameux \emph{brown envelope}, « enveloppe brune » remise pour orienter un article). Au Cameroun, elle est codifiée dans la Charte des journalistes professionnels du Cameroun.
\end{definition}

\begin{exemplebox}[Exemples types de sanctions prononcées par le CNC]
Le CNC publie régulièrement des communiqués de sanctions. Voici des cas types, inspirés de ses rapports d'activité :
\begin{center}
\begin{tabularx}{\linewidth}{|l|X|X|}\hline
\rowcolor{popblueL}
\textbf{Organe / Journaliste} & \textbf{Manquement constaté} & \textbf{Sanction possible} \\ \hline
Radio privée (Yaoundé) & Diffusion d'allégations non vérifiées sur une affaire judiciaire en cours (violation de la présomption d'innocence) & Suspension temporaire de l'émission incriminée \\ \hline
Quotidien privé (Douala) & Publication de la photo d'un mineur victime d'agression sans floutage (atteinte à la vie privée, protection des mineurs) & Blâme public et mise en demeure \\ \hline
Journaliste (web TV) & Injures publiques envers une autorité administrative lors d'un direct sur Facebook & Interdiction temporaire d'exercer et saisine du procureur \\ \hline
Chaîne TV privée & Non-respect de l'équité du temps de parole pendant une campagne électorale & Mise en demeure de rééquilibrer, suspension temporaire du responsable de publication \\ \hline
\end{tabularx}
\end{center}
\emph{Ces exemples montrent que la régulation est effective, même si des critiques persistent sur sa sélectivité perçue.}
\end{exemplebox}

\courssub{La protection concrète du citoyen : droits et recours}

Le citoyen n'est pas démuni face aux médias. La loi lui donne des armes juridiques précises.

\begin{center}
\begin{tabularx}{\linewidth}{|l|X|X|}\hline
\rowcolor{popblueL}
\textbf{Droit du citoyen} & \textbf{Fondement} & \textbf{Recours / Mise en œuvre} \\ \hline
\textbf{Droit de réponse} & Loi n° 90/052 de 1990 & Toute personne mise en cause a droit à une réponse \textbf{gratuite}, dans des conditions de visibilité équivalentes (même page, même créneau horaire). En cas de refus : saisine du CNC puis du juge. \\ \hline
\textbf{Droit à l'image / vie privée} & Code civil ; loi n° 90/052 ; loi n° 2010/012 & Interdiction de publier l'image d'une personne dans un lieu privé sans son accord. Recours : référé (retrait, dommages-intérêts) et plainte pénale. \\ \hline
\textbf{Protection des mineurs} & Loi n° 90/052 ; réglementation des programmes audiovisuels & Interdiction de diffuser des images ou sons portant atteinte à la dignité des mineurs ; signalétique d'âge obligatoire à la télévision. Le CNC sanctionne les chaînes non conformes. \\ \hline
\textbf{Protection contre la haine et les fake news} & Loi n° 2010/012 (art. 78) ; Code pénal & Signalement sur les plateformes, plainte au commissariat ou à la gendarmerie, saisine du procureur. Peines : emprisonnement et amendes lourdes. \\ \hline
\textbf{Accès à l'information publique} & Textes sur l'accès aux documents administratifs & Droit de demander des documents administratifs ; en cas de refus, recours hiérarchique puis contentieux administratif. \\ \hline
\end{tabularx}
\end{center}

\begin{definition}[Droit de réponse]
Garantie légale permettant à toute personne physique ou morale, mise en cause nominalement ou identifiable dans un organe de presse, de faire publier ou diffuser sa réponse \textbf{gratuitement}, dans le même organe, avec une visibilité équivalente (même page, même créneau horaire). Le directeur de publication ne peut le refuser sans motif légitime, sous peine de sanctions.
\end{definition}

\courssub{La lutte contre les \emph{fake news} : un défi contemporain}

La loi n° 2010/012 du 21 décembre 2010 relative à la cybersécurité et à la cybercriminalité est l'outil répressif principal contre la désinformation en ligne.

\begin{savaistu}[La loi camerounaise de 2010 punit la diffusion de fausses informations par voie électronique]
L'\textbf{article 78} de la loi n° 2010/012 punit d'un emprisonnement de six mois à deux ans et d'une amende de 500 000 à 10 000 000 de francs CFA, ou de l'une de ces deux peines seulement, quiconque, par le biais d'un système informatique, conçoit, diffuse ou met en circulation une information fausse susceptible de porter atteinte à la paix publique ou à la dignité des personnes.
\textbf{Portée :} cela couvre les SMS, WhatsApp, Facebook, X (Twitter), les blogs et les sites web. Le simple « partage » (forward) peut constituer une \emph{diffusion}. La négligence caractérisée (partager sans vérifier une information sensationnelle) peut suffire à engager la responsabilité.
\end{savaistu}

\begin{methode}[Réagir face à une information suspecte : le réflexe citoyen]
\begin{enumerate}
    \item \textbf{Stop :} ne pas partager immédiatement (casser la viralité).
    \item \textbf{Source :} qui l'a écrite ? Site officiel ? Journaliste connu ? Compte anonyme ? Adresse web bizarre ?
    \item \textbf{Croiser :} l'information est-elle reprise par la CRTV, \emph{Cameroon Tribune} ou des médias de référence (RFI, AFP, BBC) ?
    \item \textbf{Contexte :} quelle est la date de l'article ? Le lieu ? S'agit-il d'une satire ou d'un compte parodique ?
    \item \textbf{Outils :} utiliser les plateformes de vérification (\emph{Africa Check}, \emph{Dubawa}, \emph{StopIntox} au Cameroun) et la recherche d'image inversée (Google Images).
    \item \textbf{Signaler :} sur la plateforme concernée (bouton « Signaler »), sur la plateforme officielle de signalement des contenus illicites, et au commissariat ou à la gendarmerie en cas de menace grave.
\end{enumerate}
\end{methode}

\courssub{La responsabilité du citoyen-consommateur : éducation aux médias et esprit critique}

La meilleure protection reste un citoyen \textbf{lucide}. L'État ne peut pas tout contrôler, et ne le devrait pas dans une démocratie. L'\textbf{Éducation aux Médias et à l'Information (EMI)} devient une compétence civique majeure.

\begin{aretenir}[Compétences EMI du citoyen averti]
\begin{itemize}
    \item \textbf{Identifier la nature du contenu :} information (faits vérifiés), opinion (éditorial, chronique), publicité (contenu sponsorisé), divertissement ou propagande.
    \item \textbf{Comprendre l'économie de l'attention :} les algorithmes des réseaux sociaux (Facebook, TikTok, YouTube) favorisent les émotions fortes (colère, peur, indignation) car elles retiennent l'utilisateur et rapportent des revenus publicitaires. \textbf{Mécanisme :} plus un contenu est clivant, plus il devient viral.
    \item \textbf{Détecter les biais cognitifs :} biais de confirmation (je ne lis que ce qui conforte mes idées), effet de simple exposition (plus je vois une information, plus je la crois vraie).
    \item \textbf{Produire de manière responsable :} ne pas relayer une information non vérifiée, respecter le droit d'auteur, ne pas harceler en ligne (le cyberharcèlement est puni par la loi n° 2010/012).
\end{itemize}
\end{aretenir}

\begin{exoresolu}[Analyse d'une situation : « Le buzz du faux communiqué »]
\textbf{Énoncé :} un élève de Première technique reçoit sur WhatsApp la capture d'écran d'un prétendu « communiqué du MINEDUB » annonçant l'annulation du Probatoire pour cause de fuites massives. Le document porte l'en-tête du ministère, un numéro de référence et une signature scannée. L'élève le partage sur trois groupes de classe (environ 150 personnes) en dix minutes. La panique s'installe. Le MINEDUB dément deux heures plus tard via sa page Facebook officielle.

\textbf{Questions :}
\begin{enumerate}
    \item Qualifiez l'infraction commise par l'auteur du faux communiqué et par l'élève qui l'a partagé.
    \item Quels recours pour les candidats trompés ? Quel rôle pour le CNC ?
    \item Proposez une démarche de vérification que l'élève aurait dû appliquer.
\end{enumerate}

\tcblower
\textbf{Solution détaillée :}
\begin{enumerate}
    \item \textbf{Auteur du faux :} faux et usage de faux en écriture publique (Code pénal) et diffusion de fausses informations par voie électronique troublant la paix publique (art. 78 de la loi n° 2010/012). Les peines peuvent se cumuler.
    \item \textbf{L'élève qui partage :} il participe à la diffusion (art. 78). Même s'il n'est pas l'auteur, le partage volontaire sans vérification caractérise l'élément matériel de l'infraction. Il risque une amende (500 000 à 10 000 000 FCFA) et/ou un emprisonnement (six mois à deux ans). Sa bonne foi pourrait être invoquée, mais sa négligence est caractérisée.
    \item \textbf{Recours des candidats :} signalement immédiat à la police judiciaire (brigade de cybercriminalité) ; possibilité de se constituer partie civile pour réparation du préjudice moral (stress, déplacements inutiles).
    \item \textbf{Rôle du CNC :} saisir le procureur, publier un démenti, et sanctionner tout organe de presse qui aurait relayé l'information sans vérification (manquement à la déontologie : vérité, recoupement des sources).
    \item \textbf{Démarche de vérification :}
    \begin{itemize}
        \item \textbf{Stop :} ne pas partager.
        \item \textbf{Investiguer la source :} consulter le site officiel du MINEDUB ou sa page Facebook certifiée ; vérifier le numéro de référence du communiqué.
        \item \textbf{Faire une recherche croisée :} taper « communiqué MINEDUB annulation Probatoire » sur un moteur de recherche ; vérifier si des médias crédibles (CRTV, \emph{Cameroon Tribune}) en parlent.
        \item \textbf{Tracer l'origine :} le document étant une image, effectuer une recherche d'image inversée, qui révèle souvent un ancien communiqué modifié ou un document d'une autre année.
    \end{itemize}
\end{enumerate}
\end{exoresolu}

\begin{exercice}[Entraînement : cas pratiques de protection de l'opinion]
\begin{enumerate}
    \item Un journaliste publie une enquête accusant un maire de détournement de fonds, basée sur des factures anonymes reçues par courriel. Le maire nie et porte plainte pour diffamation. Le journaliste invoque la « bonne foi » et la protection des sources. \textbf{Analysez} la recevabilité de la défense du journaliste au regard de la déontologie et de la loi n° 90/052.
    \item Une chaîne de télévision privée diffuse un reportage sur une rixe entre élèves dans un lycée. On y voit clairement le visage de mineurs non floutés, certains en pleurs, d'autres menottés. Le CNC est saisi par une association de parents. \textbf{Quelles sanctions le CNC peut-il prononcer ?} Sur quel fondement ?
    \item \textbf{Débat structuré :} « La loi sur la cybercriminalité de 2010 est-elle une protection nécessaire contre les \emph{fake news} ou un outil de musellement de la liberté d'expression sur les réseaux sociaux ? » \textbf{Argumentez} en quinze lignes maximum en mobilisant l'article 78 de la loi n° 2010/012, le Préambule de la Constitution et le rôle des algorithmes.
\end{enumerate}
\end{exercice}

\begin{corrige}[Exercice — Éléments de réponse]
\textbf{1. Affaire du journaliste et du maire :} la bonne foi suppose une enquête sérieuse, une base factuelle suffisante, l'absence d'animosité personnelle et l'offre d'un droit de réponse \emph{avant} publication. Des factures \emph{anonymes} et non authentifiées ne constituent pas une base factuelle solide. La protection des sources ne dispense pas de vérifier la \emph{véracité des faits rapportés}. Le journaliste risque une condamnation pour diffamation ; il aurait dû transmettre les documents à la justice plutôt que de publier.

\textbf{2. Reportage sur les mineurs :} violation des règles de protection des mineurs (atteinte à la dignité et à la vie privée des enfants, obligation de floutage). Sanctions possibles du CNC : blâme, suspension de l'émission ou du journaliste, mise en demeure de retirer ou de flouter les images. Une saisine du procureur est possible pour atteinte à l'intimité de la vie privée.

\textbf{3. Débat :} \emph{Thèse de la protection :} l'article 78 cible l'information \emph{fausse}, diffusée par voie électronique et troublant la paix publique ; il est nécessaire face à la viralité des rumeurs (lynchages de 2019-2020). \emph{Thèse du musellement :} la notion de « fausse information » peut être floue, avec un risque d'arbitraire et un effet d'autocensure (\emph{chilling effect}). \emph{Synthèse :} l'outil est nécessaire mais exige des garde-fous : définition précise du trouble à la paix, juge indépendant, proportionnalité des peines, distinction entre fait et opinion. L'éducation aux médias reste le rempart le plus durable.
\end{corrige}

\coursec{TD3 : La publicité dans les médias}

Après avoir examiné les mécanismes de protection de l'opinion publique au Cameroun, nous nous penchons maintenant sur l'un des acteurs majeurs qui façonnent cette opinion : la publicité. Omniprésente dans notre quotidien — à la radio, à la télévision, sur les panneaux d'affichage, dans la presse et sur nos smartphones — elle ne se contente pas de vendre des produits ; elle vend des modes de vie, des valeurs et des représentations. Ce travail dirigé vise à vous doter des outils critiques pour décrypter ses messages et en mesurer l'impact réel sur vos choix de consommateurs et de citoyens.

\courssub{Qu'est-ce que la publicité ? Définition et caractéristiques}

\begin{definition}[La publicité]
La \cle{publicité} est une communication de masse, \textbf{payante}, diffusée par un \cle{annonceur} (entreprise, association, institution) via des \cle{supports médias} (radio, télévision, presse, internet, affichage), dans le but d'\textbf{influencer les comportements} (achat, adhésion, changement d'attitude) d'un public \cle{cible}.
\end{definition}

\begin{aretenir}[Trois critères distinctifs]
\begin{itemize}
    \item \textbf{Payante :} l'annonceur achète l'espace ou le temps de diffusion (contrairement aux relations publiques ou au bouche-à-oreille).
    \item \textbf{Non personnelle :} le message s'adresse à une multitude, pas à un individu (contrairement à la vente directe).
    \item \textbf{Identifiée :} l'origine du message (l'annonceur) est connue ou identifiable.
\end{itemize}
\end{aretenir}

\begin{definition}[Le slogan]
Le \cle{slogan} (ou \emph{accroche}) est une phrase courte, percutante et facile à mémoriser, qui résume la promesse centrale de la marque ou du produit. Il repose souvent sur des figures de style (rime, allitération, antithèse). \textbf{Exemples :} « MTN, \emph{Everywhere you go} » ; « Orange, \emph{la vie change avec Orange} ».
\end{definition}

\courssub{Les supports publicitaires au Cameroun : un paysage médiatique pluriel}

Au Cameroun, le marché publicitaire s'appuie sur des supports traditionnels et numériques. Le choix du support dépend de la cible visée, du budget disponible et de l'objectif poursuivi.

\begin{center}
\begin{tabularx}{\linewidth}{|l|X|X|}
\hline
\rowcolor{popblueL}
\textbf{Support} & \textbf{Caractéristiques / Portée} & \textbf{Exemples camerounais} \\
\hline
\textbf{Télévision (CRTV, Canal 2 International, STV, Equinoxe TV)} & Fort impact audiovisuel, couverture nationale (CRTV), ciblage urbain (chaînes privées). Tranche de grande écoute (19 h – 22 h) très prisée. & Spots brasseries (33 Export, Castel), télécoms (MTN, Orange), banques (Afriland First Bank, BICEC). \\
\hline
\textbf{Radio (CRTV Radio, Balafon FM, Nostalgie, Sweet FM, radios communautaires)} & Proximité, langues locales, écoute mobile (taxi, marché, champ). Coût accessible pour les PME. & Spots en langues nationales (bassa, douala, ewondo, fulfuldé), annonces d'événements culturels. \\
\hline
\textbf{Presse écrite (Cameroon Tribune, Mutations, Le Jour, L'Œil du Sahel)} & Crédibilité, conservation du message, ciblage des décideurs. Diffusion papier en recul. & Encarts institutionnels (ministères), annonces légales, suppléments spéciaux (rentrée scolaire, fêtes). \\
\hline
\textbf{Affichage urbain (Douala, Yaoundé, Bafoussam, Garoua)} & Visibilité massive, répétition forcée (trajets quotidiens). Panneaux 4 m $\times$ 3 m, abribus, habillage de bus. & Grandes marques (cimenteries, savonneries, brasseries), campagnes électorales, événements. \\
\hline
\textbf{Numérique / Réseaux sociaux (Facebook, WhatsApp, Instagram, TikTok, YouTube)} & Ciblage précis (âge, lieu, centres d'intérêt), interactivité (clic, partage), mesure de performance en temps réel. & Marketing d'influence, publications sponsorisées, publicités avant les vidéos YouTube, messages WhatsApp Business. \\
\hline
\end{tabularx}
\end{center}

\begin{popfigure}
\begin{tikzpicture}[
    node distance=1.2cm and 1.1cm,
    box/.style={draw=popblue, thick, rounded corners, fill=popblue!10, minimum width=2.9cm, minimum height=1.3cm, align=center, font=\small},
    arrow/.style={-Stealth, thick, popdark},
    start/.style={draw=popgreen, thick, rounded corners, fill=popgreen!15, minimum width=2.9cm, minimum height=1.3cm, align=center, font=\small\bfseries},
    end/.style={draw=poporange, thick, rounded corners, fill=poporange!15, minimum width=2.9cm, minimum height=1.3cm, align=center, font=\small\bfseries},
]
\node[start, align=center] (cible){1. CIBLE\\Qui ? (âge, sexe,\\localisation,\\centres d'intérêt)};
\node[box, right=of cible, align=center] (message){2. MESSAGE\\Quoi ? (promesse,\\argument, preuve,\\ton, slogan)};
\node[box, right=of message, align=center] (support){3. SUPPORT\\Où ? (TV, radio, web,\\affichage, presse,\\influenceurs)};
\node[end, right=of support, align=center] (effet){4. EFFET RECHERCHÉ\\Résultat ? (notoriété,\\intention d'achat,\\fidélisation)};

\draw[arrow] (cible.east) -- node[above, font=\tiny, popdark] {détermine} (message.west);
\draw[arrow] (message.east) -- node[above, font=\tiny, popdark] {conditionne} (support.west);
\draw[arrow] (support.east) -- node[above, font=\tiny, popdark] {produit} (effet.west);

\draw[arrow, poppurple, dashed, bend left=45] (effet.north) to node[above, font=\tiny, poppurple] {Mesure des résultats et ajustement} (cible.north);
\end{tikzpicture}
\legende{Schéma : les quatre étapes logiques de la construction d'un message publicitaire efficace. La boucle en pointillés représente l'analyse des résultats, qui permet d'affiner la cible ou le message.}
\end{popfigure}

\begin{exemplebox}[Exemple concret : lancement d'une nouvelle boisson énergisante « Vital+ »]
\begin{itemize}
    \item \textbf{Cible :} jeunes urbains de 15 à 25 ans, élèves, étudiants, apprentis, amateurs de sport et de musique, à Douala et Yaoundé.
    \item \textbf{Message :} « Vital+ : l'énergie qui ne dort jamais. » (Promesse : endurance et dynamisme.) Ton : jeune, direct, familier.
    \item \textbf{Supports :} spots TV de 30 secondes pendant les matchs de football, vidéos courtes sur TikTok et Instagram (défi de danse), affichage près des lycées et universités, parrainage d'un tournoi de football de quartier.
    \item \textbf{Effet recherché :} faire connaître la marque auprès d'au moins 60 \% de la cible en trois mois et atteindre un volume de ventes régulier.
\end{itemize}
\end{exemplebox}

\courssub{Objectifs et techniques de persuasion : comment ça marche ?}

\begin{propriete}[Les trois objectifs fondamentaux de la publicité]
\begin{enumerate}
    \item \textbf{Faire connaître (notoriété) :} « J'existe, je suis là. » Objectif essentiel lors du lancement d'un produit.
    \item \textbf{Persuader (préférence, intention d'achat) :} « Je suis le meilleur choix pour toi. » Argumentation rationnelle (prix, qualité) ou émotionnelle (statut social, plaisir, peur).
    \item \textbf{Fidéliser (action répétée) :} « Reste avec moi. » Programmes de fidélité, communauté de marque, rappels constants.
\end{enumerate}
\end{propriete}

\begin{aretenir}[Principales techniques de persuasion (la boîte à outils du publicitaire)]
\begin{itemize}
    \item \textbf{Témoignage / parrainage par une célébrité :} une star, un expert ou « une personne comme vous » vante le produit. La crédibilité ou la sympathie liée à cette personnalité est transférée au produit. (Exemple : un footballeur célèbre vantant une marque de téléphonie.)
    \item \textbf{Slogan / jingle :} mémorisation par la répétition rythmée ou mélodique.
    \item \textbf{Images attractives :} beauté, humour, rêve, évasion. Le plaisir visuel détourne l'attention critique.
    \item \textbf{Promesses exagérées :} « le meilleur du monde », « zéro défaut ». Tolérées si elles sont purement subjectives, mais trompeuses si elles sont mesurables et fausses.
    \item \textbf{Preuve sociale :} « tout le monde l'achète », « numéro 1 des ventes ». Joue sur la conformité et la peur de rater quelque chose.
    \item \textbf{Urgence / pénurie :} « offre limitée », « stocks bientôt épuisés ». Provoque l'achat impulsif.
    \item \textbf{Musique / sonorité :} associe durablement une émotion (joie, nostalgie) à la marque.
\end{itemize}
\end{aretenir}

\begin{attention}[Publicité commerciale et information journalistique : ne les confondez pas !]
\begin{itemize}
    \item \textbf{Publicité :} espace \textbf{acheté}. But : \textbf{vendre / convaincre}. Contenu \textbf{contrôlé par l'annonceur}. Subjectivité assumée. Elle doit être \textbf{identifiée} comme telle (mentions « Publireportage », « Sponsorisé », « Annonce »).
    \item \textbf{Information journalistique :} espace \textbf{gagné} (intérêt public). But : \textbf{informer / éclairer}. Contenu \textbf{indépendant} (ligne éditoriale). Objectivité recherchée (vérification des faits, contradiction des sources).
    \item \textbf{Piège :} le \emph{publireportage} (article payant qui imite l'aspect d'un article d'information) ou le \emph{placement de produit} (produit visible dans une série ou une émission). \textbf{Posez-vous toujours la question :} « Est-ce une information ou une publicité ? »
\end{itemize}
\end{attention}

\courssub{Publicité commerciale et campagnes d'intérêt général}

\begin{center}
\begin{tabularx}{\linewidth}{|l|X|X|}
\hline
\rowcolor{popblueL}
\textbf{Critère} & \textbf{Publicité commerciale} & \textbf{Campagne d'intérêt général (service public / ONG)} \\
\hline
\textbf{Annonceur} & Entreprise privée (brasserie, banque, télécom) & Ministère (Santé publique, Éducation), ONG (Plan International, Croix-Rouge), organismes publics \\
\hline
\textbf{But ultime} & Profit, part de marché, ventes & Changement de comportement social, santé publique, citoyenneté \\
\hline
\textbf{Message} & « Achetez \emph{notre} produit » & « Adoptez \emph{ce} comportement » (vaccination, ceinture de sécurité, paiement des impôts, scolarisation des filles) \\
\hline
\textbf{Techniques} & Stars, slogans, émotions, humour, peur & Identiques : personnalités ambassadrices, slogans, émotions, humour \\
\hline
\textbf{Financement} & Budget marketing de l'entreprise & Budget de l'État, bailleurs de fonds, fonds propres des ONG \\
\hline
\textbf{Exemple camerounais} & Spot télévisé pour une bière ou un réseau mobile & « Moustiquaire imprégnée : toute la famille, toutes les nuits » (Ministère de la Santé publique / Programme national de lutte contre le paludisme) \\
\hline
\end{tabularx}
\end{center}

\begin{savaistu}[Le saviez-vous ?]
Les grandes campagnes de sensibilisation contre le \textbf{paludisme} (distribution de moustiquaires imprégnées) ou pour la \textbf{vaccination} (poliomyélite, rougeole, COVID-19) au Cameroun utilisent \textbf{exactement les mêmes techniques} que la publicité commerciale : jingles diffusés en boucle à la CRTV et sur les radios privées, spots télévisés avec des personnalités (artistes, footballeurs), affiches colorées dans les centres de santé et les marchés, messages SMS et WhatsApp de masse, influenceurs locaux mobilisés. La différence ? Le « produit » vendu est un \textbf{comportement de santé} gratuit, et le « profit » escompté est \textbf{collectif} : la baisse de la mortalité, notamment infantile.
\end{savaistu}

\courssub{Les dérives de la publicité : quand l'éthique est oubliée}

\begin{aretenir}[Trois dérives majeures à repérer]
\begin{enumerate}
    \item \textbf{Publicité mensongère ou trompeuse :} allégations fausses sur les propriétés d'un produit (exemple : « ce jus guérit le diabète »), prix cachés, photos excessivement retouchées, comparaisons biaisées. \textbf{Interdite par la loi.}
    \item \textbf{Ciblage des enfants :} exploitation de la naïveté et de l'incapacité critique des plus jeunes (jouets, sucreries, jeux vidéo). Pression exercée sur les parents par le harcèlement de l'enfant. \textbf{Encadrée :} pas d'appel direct à l'achat adressé aux enfants, protection des programmes jeunesse.
    \item \textbf{Promotion de produits nocifs :} alcool, tabac (cigarettes, chicha), paris sportifs et jeux d'argent, médicaments vendus sans ordonnance, produits de dépigmentation de la peau. \textbf{Risques :} dépendance, surendettement, cancers, troubles psychiques.
\end{enumerate}
\end{aretenir}

\begin{exemplebox}[Cas concret : les paris sportifs en ligne au Cameroun]
Depuis la fin des années 2010, les opérateurs de paris sportifs en ligne investissent massivement dans la publicité au Cameroun (maillots d'équipes locales, panneaux géants, spots télévisés, influenceurs sur les réseaux sociaux). \textbf{Dérive :} ciblage massif des jeunes (élèves, étudiants, jeunes chômeurs) avec la promesse de « gains faciles » et d'un « changement de vie ». Conséquences observées : addiction, détournement des frais de scolarité ou du loyer, endettement, dépression. \textbf{Réponse des pouvoirs publics :} encadrement progressif de la publicité pour les jeux d'argent (interdiction de cibler les mineurs, messages de mise en garde obligatoires du type « Jouer comporte des risques »).
\end{exemplebox}

\courssub{La régulation de la publicité au Cameroun : cadre légal et acteurs}

\begin{propriete}[Cadre juridique principal]
\begin{itemize}
    \item \textbf{Loi n° 90/052 du 19 décembre 1990} relative à la liberté de communication sociale (base : liberté d'expression, mais aussi responsabilité).
    \item \textbf{Loi n° 2011/012 du 6 mai 2011} relative à la protection du consommateur au Cameroun (information loyale, interdiction de la publicité mensongère, clauses abusives).
    \item \textbf{Textes sectoriels} (décrets et arrêtés ministériels) encadrant la publicité pour certains produits sensibles : médicaments, alcool, tabac, jeux d'argent.
\end{itemize}
\end{propriete}

\begin{definition}[Le Conseil National de la Communication (CNC)]
Organe de \textbf{régulation administrative} des médias au Cameroun. \textbf{Rôle clé en matière de publicité :}
\begin{itemize}
    \item veiller au respect de la déontologie (véracité, dignité humaine, non-discrimination, protection des mineurs) ;
    \item être saisi sur plainte (consommateur, association, concurrent) ou agir d'office ;
    \item pouvoirs : \textbf{mise en demeure}, \textbf{interdiction de diffusion}, \textbf{sanctions pécuniaires} (amendes), suspension ou retrait d'autorisation du support récalcitrant.
\end{itemize}
\end{definition}

\begin{exemplebox}[Exemple chiffré : les formats de spots publicitaires à la télévision camerounaise]
À la CRTV comme sur les chaînes privées, le format le plus courant reste le \textbf{spot de 30 secondes}.
\begin{itemize}
    \item \textbf{15 secondes :} rappel, teasing ou jingle seul (coût réduit, forte rotation).
    \item \textbf{30 secondes :} format standard « complet » (problème – solution – marque – slogan). \textbf{Format le plus vendu.}
    \item \textbf{45 à 60 secondes :} lancement majeur d'un nouveau produit, récit construit, publireportage court. Coût élevé, diffusion plus rare.
    \item \textbf{Tarification :} le prix d'un passage dépend de la chaîne, de l'horaire (les tranches de grande écoute coûtent plus cher) et du volume acheté ; les grands annonceurs négocient des \emph{forfaits} (par exemple un nombre de passages par mois).
\end{itemize}
\end{exemplebox}

\begin{popfigure}
\begin{tikzpicture}[scale=1]
% Diagramme circulaire manuel : TV 35%, Radio 20%, Affichage 15%, Numérique 12%, Presse 10%, Autres 8%
% Angles cumulés (départ 90°, sens horaire) : TV 126°, Radio 72°, Affichage 54°, Numérique 43.2°, Presse 36°, Autres 28.8°
\fill[popblue!80]   (0,0) -- (90:3)   arc (90:-36:3)   -- cycle;
\fill[poporange!80] (0,0) -- (-36:3)  arc (-36:-108:3) -- cycle;
\fill[popgreen!80]  (0,0) -- (-108:3) arc (-108:-162:3)-- cycle;
\fill[poppurple!80] (0,0) -- (-162:3) arc (-162:-205.2:3) -- cycle;
\fill[poppink!80]   (0,0) -- (-205.2:3) arc (-205.2:-241.2:3) -- cycle;
\fill[popgold!80]   (0,0) -- (-241.2:3) arc (-241.2:-270:3) -- cycle;
\draw[popdark, thick] (0,0) circle (3);
% Légende
\node[anchor=west, font=\small] at (3.6,2.2)  {\textcolor{popblue!80}{\rule{0.35cm}{0.35cm}}~TV (CRTV et privées) : 35 \%};
\node[anchor=west, font=\small] at (3.6,1.5)  {\textcolor{poporange!80}{\rule{0.35cm}{0.35cm}}~Radio : 20 \%};
\node[anchor=west, font=\small] at (3.6,0.8)  {\textcolor{popgreen!80}{\rule{0.35cm}{0.35cm}}~Affichage urbain : 15 \%};
\node[anchor=west, font=\small] at (3.6,0.1)  {\textcolor{poppurple!80}{\rule{0.35cm}{0.35cm}}~Numérique (réseaux sociaux, sites) : 12 \%};
\node[anchor=west, font=\small] at (3.6,-0.6) {\textcolor{poppink!80}{\rule{0.35cm}{0.35cm}}~Presse écrite : 10 \%};
\node[anchor=west, font=\small] at (3.6,-1.3) {\textcolor{popgold!80}{\rule{0.35cm}{0.35cm}}~Autres (cinéma, événementiel, SMS) : 8 \%};
\end{tikzpicture}
\legende{Diagramme circulaire : répartition estimée des dépenses publicitaires par support au Cameroun (ordres de grandeur indicatifs, années récentes). La télévision domine encore, mais la part du numérique progresse rapidement d'année en année.}
\end{popfigure}

\courssub{Travail dirigé : analyse guidée d'une publicité réelle (méthode en 5 questions)}

\begin{methode}[Comment analyser une publicité en 5 questions critiques (grille d'analyse)]
Appliquez systématiquement cette grille à \textbf{n'importe quelle publicité} (télévision, radio, affichage, Facebook, TikTok) :
\begin{enumerate}
    \item \textbf{QUI ? (Émetteur / annonceur)} Quelle entreprise, marque ou institution ? Quelle est sa réputation ? Quel intérêt financier ou idéologique poursuit-elle ?
    \item \textbf{QUOI ? (Contenu / promesse)} Quel produit, service ou comportement est promu ? Quelle est la \textbf{promesse centrale} (explicite ou implicite) ? Est-elle \textbf{vérifiable} ?
    \item \textbf{À QUI ? (Cible / récepteur)} Âge, sexe, lieu, niveau de vie, centres d'intérêt, vulnérabilités (enfants, personnes âgées, malades, personnes précaires) ? Pourquoi \emph{cette} cible ?
    \item \textbf{COMMENT ? (Forme / techniques)} Quelles techniques de persuasion sont utilisées (star, humour, peur, musique, slogan, images, preuve sociale, urgence) ? Quel \textbf{registre} (rationnel ou émotionnel) ? Quelle \textbf{langue} (français, anglais, camfranglais, langue locale) ?
    \item \textbf{AVEC QUEL EFFET ? (Impact visé / réel)} Effet \textbf{cognitif} (je sais), \textbf{affectif} (j'aime), \textbf{comportemental} (j'achète, j'adopte). La publicité est-elle \textbf{éthique} ? \textbf{Légale} ? Respecte-t-elle la dignité des personnes ? S'agit-il de \textbf{manipulation} ou d'\textbf{information loyale} ?
\end{enumerate}
\end{methode}

\begin{exoresolu}[Analyse guidée : spot TV « Super Malt – L'énergie du matin » (exemple type, 30 s)]
\textbf{Description :} lever de soleil sur Douala. Une jeune femme (25 ans, tenue professionnelle) se réveille fatiguée. Elle boit une canette de la boisson maltée. Son visage s'illumine. Elle court énergiquement vers son travail. Slogan : « Super Malt, \emph{l'énergie qui vous ressemble}. » Voix off : « Malt, vitamines, fer. Pour bien commencer la journée. » Logo et canette à l'écran, avec une mention de modération en petits caractères en bas de l'image.

\medskip
\textbf{Application de la grille :}
\begin{enumerate}
    \item \textbf{QUI ?} Une grande brasserie camerounaise, leader du marché des boissons. Intérêt : vendre du volume et fidéliser les jeunes actifs.
    \item \textbf{QUOI ?} Boisson maltée non alcoolisée. \textbf{Promesse :} énergie physique et mentale, réussite professionnelle, vitalité matinale. \textbf{Vérifiable ?} L'apport calorique (sucre) et les vitamines procurent une énergie réelle, mais \emph{pas} « magique » ; la teneur en sucre est élevée (attention au diabète et à l'obésité).
    \item \textbf{À QUI ?} Jeunes adultes urbains (20–35 ans), actifs, travailleurs, étudiants, aspirant à la réussite. \textbf{Vulnérabilité exploitée :} fatigue chronique, pression de la performance, recherche d'un « coup de boost ».
    \item \textbf{COMMENT ?} \textbf{Identification :} une héroïne « comme moi ». \textbf{Argument pseudo-scientifique :} « vitamines, fer ». \textbf{Émotion :} lever de soleil = espoir, course = dynamisme. \textbf{Slogan court.} \textbf{Langue :} français standard (cible lettrée). Mention légale discrète.
    \item \textbf{AVEC QUEL EFFET ?} \textbf{Visé :} achat quotidien, routine du « petit-déjeuner ». \textbf{Réel :} normalisation d'une consommation sucrée quotidienne ; association positive entre la marque et la réussite. \textbf{Éthique ?} Discutable : la boisson est présentée comme un produit « santé » alors qu'il s'agit d'une boisson sucrée. \textbf{Légal ?} Oui, si les mentions obligatoires sont présentes, mais la promesse nutritionnelle frôle la limite de l'allégation trompeuse.
\end{enumerate}
\end{exoresolu}

\begin{exercice}[À vous de jouer : analysez une publicité de votre environnement]
Choisissez \textbf{une} publicité actuelle que vous avez vue ou entendue cette semaine (spot TV sur CRTV ou STV, annonce radio sur Balafon FM ou Nostalgie, panneau d'affichage, publication sponsorisée Facebook/TikTok d'un influenceur camerounais, prospectus distribué au marché).
\begin{enumerate}
    \item Présentez-la brièvement (support, description visuelle ou sonore, slogan).
    \item Appliquez la \textbf{méthode des 5 questions} (Qui ? Quoi ? À qui ? Comment ? Avec quel effet ?).
    \item Concluez : \textbf{cette publicité est-elle loyale, éthique, efficace ?} Justifiez en trois arguments.
\end{enumerate}
\textit{Durée : 30 minutes en classe ou en devoir. Restitution orale ou écrite (une page maximum).}
\end{exercice}

\begin{corrige}[Exercice n°1 – Pistes de correction type (exemple : publicité « MTN Mobile Money – Envoie l'argent à maman au village »)]
\begin{enumerate}
    \item \textbf{Présentation :} spot TV de 30 secondes en tranche de grande écoute. Un jeune homme en ville (Douala) envoie 10 000 FCFA par Mobile Money à sa mère au village. La mère reçoit l'argent, sourit, achète des médicaments et de la nourriture. Slogan : « MTN MoMo, \emph{proche de vous, partout}. » Langues : français et langues nationales.
    \item \textbf{Analyse en 5 questions :}
        \begin{itemize}
            \item \textbf{QUI :} MTN Cameroon, leader des télécommunications. Intérêt : faire adopter le service Mobile Money, augmenter le volume des transactions, fidéliser les abonnés.
            \item \textbf{QUOI :} service de transfert d'argent instantané par téléphone. Promesse : rapidité, sécurité, proximité émotionnelle, utilité sociale (santé, alimentation). Vérifiable : oui, le service fonctionne techniquement.
            \item \textbf{À QUI :} jeunes travailleurs et commerçants urbains ayant de la famille au village. Cible large et transgénérationnelle. Vulnérabilité exploitée : le devoir filial et la culpabilité de la distance.
            \item \textbf{COMMENT :} \textbf{émotion forte} (lien mère-fils, soulagement), \textbf{preuve par l'usage} (démonstration en temps réel), \textbf{langues locales} (inclusion, confiance), \textbf{simplicité} (pas de jargon technique).
            \item \textbf{EFFET :} adoption ou réactivation du compte Mobile Money. \textbf{Éthique :} valorise l'entraide familiale et un usage réel du service. \textbf{Légal :} conforme (conditions générales et numéro court mentionnés).
        \end{itemize}
    \item \textbf{Conclusion :} \textbf{publicité loyale et éthique.} 1) Le service répond à un besoin réel (inclusion financière des populations). 2) La mise en scène est authentique (langues réelles, situations de la vie quotidienne). 3) Aucune promesse mensongère (on ne promet pas de « gagner » de l'argent, seulement de le transférer). \textbf{Efficace :} oui, grâce à une forte résonance culturelle.
\end{enumerate}
\end{corrige}

\begin{aretenir}[Bilan du TD3 : ce qu'il faut retenir pour l'examen et pour la vie]
\begin{itemize}
    \item La publicité est un \textbf{pouvoir d'influence massif} (économique, culturel, politique). La comprendre, c'est reprendre son \textbf{libre arbitre}.
    \item Au Cameroun, elle est \textbf{plurilingue, multi-supports et en mutation numérique rapide}.
    \item La \textbf{régulation existe} (CNC, loi sur la protection du consommateur, textes sectoriels), mais son \textbf{application} reste perfectible (moyens de contrôle limités, réactivité insuffisante face au numérique).
    \item \textbf{Outils du citoyen averti :} grille d'analyse (5 questions), esprit critique, signalement au CNC ou aux associations de consommateurs, éducation aux médias et à l'information.
    \item \textbf{Distinction vitale :} publicité commerciale (profit privé) $\neq$ campagne d'intérêt général (bien commun) $\neq$ information journalistique (devoir de vérité).
\end{itemize}
\end{aretenir}

\newpage

\coursec{Activité d'intégration}

\courssub{Objectifs de l'activité}
\begin{itemize}
    \item \textbf{Mobiliser} les connaissances sur les médias, l'opinion publique, la désinformation et le cadre juridique camerounais de protection de l'opinion.
    \item \textbf{Concevoir et administrer} un questionnaire d'enquête d'opinion auprès d'un échantillon de 30 personnes.
    \item \textbf{Traiter des données statistiques} brutes : calcul de fréquences, pourcentages, construction d'un diagramme en barres et d'un diagramme circulaire.
    \item \textbf{Analyser} un cas réel de \cle{fake news} ayant circulé au Cameroun (source, propagation, démenti).
    \item \textbf{Rédiger} un rapport d'enquête structuré, argumenté et conclu par trois conseils pratiques pour protéger l'opinion publique locale.
    \item \textbf{Communiquer} oralement les résultats et défendre sa démarche lors d'une restitution en classe.
\end{itemize}

\courssub{Situation}
Vous êtes élèves en classe de Première Technique au \textbf{Lycée Technique de Bafoussam}. Dans le cadre du module ECM « Les médias et leur impact sur l'opinion publique », votre professeur vous propose de réaliser une \textbf{enquête de terrain} pour mesurer la réalité médiatique de votre environnement immédiat (quartier, marché, établissement, famille).

\textbf{Contexte local :} La région de l'Ouest, et particulièrement Bafoussam, est un carrefour commercial et universitaire dense. La population y est très connectée (couverture 4G, fibre optique dans certains quartiers) mais reste fortement attachée aux médias traditionnels (radio communautaire \emph{Radio Medumba}, \emph{CRTV Radio}, \emph{CRTV Télévision}, chaînes privées comme \emph{Equinoxe TV} ou \emph{Canal 2 International}). La circulation de rumeurs sur les réseaux sociaux (WhatsApp, Facebook) y est rapide, notamment lors des périodes électorales ou de crises sanitaires.

\textbf{Votre mission :} En groupes de 4 à 5 élèves, vous devez interroger \textbf{exactement 30 personnes} (au total par groupe) à l'aide d'un questionnaire commun validé par le professeur. Le questionnaire porte sur trois dimensions :
\begin{enumerate}
    \item \textbf{Préférence médiatique :} « Quel est le média par lequel vous vous informez \emph{le plus souvent} ? » (Réponses proposées : Radio / Télévision / Réseaux sociaux / Presse écrite / Bouche-à-oreille).
    \item \textbf{Confiance accordée :} « Globalement, quel est votre niveau de confiance dans les informations diffusées par votre média principal ? » (Réponses : Élevée / Moyenne / Faible).
    \item \textbf{Exposition aux \cle{fake news} :} « Avez-vous déjà cru ou partagé une information qui s'est avérée fausse par la suite (fake news) ? » (Réponses : Oui / Non).
\end{enumerate}

\textbf{Données brutes simulées pour le corrigé (Groupe « Éthique Média ») :}
Le groupe a interrogé 30 personnes (15 femmes, 15 hommes, âgées de 18 à 60 ans, commerçants, élèves, fonctionnaires, artisans). Voici les réponses comptabilisées :

\begin{center}
\begin{tabularx}{\linewidth}{|l|X|}\hline
\textbf{Variable : Média principal} & \textbf{Effectifs} \\ \hline
Radio (CRTV, Medumba, FM locales) & 8 \\ \hline
Télévision (CRTV, Equinoxe, Canal 2) & 10 \\ \hline
Réseaux sociaux (Facebook, WhatsApp, TikTok) & 9 \\ \hline
Presse écrite (Cameroon Tribune, Mutations, Le Jour) & 2 \\ \hline
Bouche-à-oreille (famille, amis, marché) & 1 \\ \hline
\textbf{Total} & \textbf{30} \\ \hline
\end{tabularx}
\end{center}

\begin{center}
\begin{tabularx}{\linewidth}{|l|X|}\hline
\textbf{Variable : Niveau de confiance} & \textbf{Effectifs} \\ \hline
Élevée & 5 \\ \hline
Moyenne & 15 \\ \hline
Faible & 10 \\ \hline
\textbf{Total} & \textbf{30} \\ \hline
\end{tabularx}
\end{center}

\begin{center}
\begin{tabularx}{\linewidth}{|l|X|}\hline
\textbf{Variable : Fake news rencontrée} & \textbf{Effectifs} \\ \hline
Oui & 18 \\ \hline
Non & 12 \\ \hline
\textbf{Total} & \textbf{30} \\ \hline
\end{tabularx}
\end{center}

\courssub{Consignes}
\textbf{Phase 1 : Préparation et collecte (Séance 1)}
\begin{enumerate}
    \item Valider collectivement le questionnaire (ajout éventuel d'une variable « Âge » ou « Profession » pour croiser les données).
    \item Répartir les zones d'enquête (Marché A, Quartier B, Lycée, Arrêt de bus) pour diversifier l'échantillon.
    \item Effectuer l'enquête (en présentiel ou par téléphone) et centraliser les résultats dans un tableau de dépouillement unique par groupe.
\end{enumerate}

\textbf{Phase 2 : Traitement statistique et analyse (Séance 2)}
\begin{enumerate}
    \setcounter{enumi}{3}
    \item Calculer les \textbf{fréquences en pourcentage} pour chaque modalité des trois variables. Arrondir à une décimale.
    \item Construire un \textbf{diagramme en barres} représentant la répartition des médias principaux (axe des ordonnées : pourcentage ; axe des abscisses : médias). Titre, légende, unité obligatoires.
    \item Construire un \textbf{diagramme circulaire} (camembert) représentant la répartition du niveau de confiance. Pourcentages et étiquettes sur le graphique.
    \item \textbf{Identifier le média le plus influent} dans votre échantillon et justifier votre choix en croisant « Média principal » et « Confiance » (ex : la TV est première en audience mais la confiance y est-elle forte ?).
\end{enumerate}

\textbf{Phase 3 : Étude de cas \cle{fake news} (Séance 2-3)}
\begin{enumerate}
    \setcounter{enumi}{7}
    \item Choisir \textbf{une fake news réelle} ayant circulé au Cameroun ces 24 derniers mois (ex : recrutement fictif à la Fonction Publique, rumeur sur la santé d'une autorité, fausse annonce de gratuité scolaire, alerte enlèvement d'enfants).
    \item Remplir la fiche d'analyse ci-dessous :
    \begin{itemize}
        \item \textbf{Contenu :} Résumé du message faux.
        \item \textbf{Source initiale :} D'où vient le message ? (Capture d'écran WhatsApp, page Facebook anonyme, site web mimétique).
        \item \textbf{Vecteurs de propagation :} Groupes WhatsApp (combien ?), Partages Facebook, Relais par des influenceurs locaux, Bouche-à-oreille.
        \item \textbf{Démenti officiel :} Qui a démenti ? (Ministère concerné, Police, CRTV, Organe de régulation \cle{HAC}/\cle{CNC}, Plateforme de fact-checking comme \emph{DataCheck} ou \emph{StopIntox}). Date du démenti.
        \item \textbf{Impact observé :} Panique, pertes financières (frais de dossier arnaqués), troubles à l'ordre public, discrédit des institutions.
    \end{itemize}
\end{enumerate}

\textbf{Phase 4 : Rédaction et restitution (Séance 3)}
\begin{enumerate}
    \setcounter{enumi}{8}
    \item Rédiger un \textbf{rapport d'enquête} (2 à 3 pages manuscrites ou 800-1000 mots tapés) structuré ainsi :
    \begin{itemize}
        \item Introduction : Contexte, objectif, méthodologie (échantillon, dates).
        \item Résultats : Tableaux de fréquences, graphiques (insérés ou décrits), analyse croisée médium/confiance.
        \item Étude de cas : Fiche fake news analysée.
        \item \textbf{Recommandations :} \textbf{Trois conseils pratiques, numérotés, applicables immédiatement} par les habitants de votre quartier pour se protéger de la désinformation.
        \item Conclusion : Engagement citoyen personnel.
    \end{itemize}
    \item Présenter le rapport à la classe (5 min par groupe) et répondre aux questions (débat contradictoire).
\end{enumerate}

\courssub{Ressources à mobiliser}
\begin{definition}[Notions clés du programme ECM 1re Tech]
\textbf{Médias :} Ensemble des moyens de communication de masse (presse, radio, TV, internet). \textbf{Opinion publique :} Ensemble des attitudes, croyances, jugements partagés par un groupe sur une question d'intérêt général. \textbf{Influence médiatique :} Capacité des médias à fixer l'agenda (\cle{agenda-setting}), à cadrer les sujets (\cle{framing}) et à primer certains thèmes (\cle{priming}). \textbf{Fake news (Infox) :} Information fausse, délibérément fabriquée ou décontextualisée, diffusée dans un but de nuisance ou de profit. \textbf{Droit de réponse :} Droit pour toute personne mise en cause de faire publier sa réponse gratuitement (Loi 90/052 du 19 déc 1990, art. 11). \textbf{HAC (Haut Conseil de l'Audiovisuel) :} Organe de régulation du secteur audiovisuel (Loi 2019/018). \textbf{CNC (Conseil National de la Communication) :} Organe de régulation de la communication sociale (Loi 2019/019).
\end{definition}

\begin{methode}[Calcul de pourcentage et construction graphique]
\textbf{1. Fréquence en \% :} $f_i (\%) = \frac{n_i}{N} \times 100$ où $n_i$ = effectif de la modalité, $N$ = effectif total (30).\\
\textbf{2. Diagramme en barres :} Une barre par modalité qualitative. Hauteur = \%. Espace entre les barres.\\
\textbf{3. Diagramme circulaire :} Angle du secteur $= \frac{n_i}{N} \times 360^\circ$ ou $\% \times 3,6$.\\
\textbf{4. Analyse croisée :} Comparer le rang d'audience et le rang de confiance. Un média très suivi mais peu fiable signale une \cle{consommation contrainte} ou un défaut d'alternative.
\end{methode}

\begin{attention}[Pièges fréquents]
\begin{itemize}
    \item Confondre « média le plus cité » et « média le plus influent » (l'influence suppose un effet sur les opinions/attitudes, pas seulement l'audience).
    \item Oublier d'indiquer la source et la date sur les graphiques.
    \item Proposer des conseils vagues (« Faire attention ») au lieu d'actions concrètes (« Vérifier sur le site officiel du Ministère avant de partager »).
    \item Ne pas citer le cadre légal camerounais (HAC, CNC, Loi sur la cybercriminalité Loi 2010/012) dans les recommandations.
\end{itemize}
\end{attention}

\courssub{Démarche et corrigé commenté}
\begin{exoresolu}[Corrigé type pour le Groupe « Éthique Média »]
\textbf{Étape 1 : Tableaux de fréquences en pourcentages}

\textit{Variable 1 : Média principal (N=30)}
\begin{align*}
\text{Radio} &: \frac{8}{30} \times 100 = 26,7\% \\
\text{Télévision} &: \frac{10}{30} \times 100 = 33,3\% \\
\text{Réseaux sociaux} &: \frac{9}{30} \times 100 = 30,0\% \\
\text{Presse écrite} &: \frac{2}{30} \times 100 = 6,7\% \\
\text{Bouche-à-oreille} &: \frac{1}{30} \times 100 = 3,3\%
\end{align*}

\textit{Variable 2 : Niveau de confiance (N=30)}
\begin{align*}
\text{Élevée} &: \frac{5}{30} \times 100 = 16,7\% \\
\text{Moyenne} &: \frac{15}{30} \times 100 = 50,0\% \\
\text{Faible} &: \frac{10}{30} \times 100 = 33,3\%
\end{align*}

\textit{Variable 3 : Fake news rencontrée (N=30)}
\begin{align*}
\text{Oui} &: \frac{18}{30} \times 100 = 60,0\% \\
\text{Non} &: \frac{12}{30} \times 100 = 40,0\%
\end{align*}

\textbf{Étape 2 : Construction des graphiques (Description pour réalisation sur papier millimétré / Excel / Python)}

\begin{popfigure}
\begin{tikzpicture}
\begin{axis}[
    ybar,
    bar width=12pt,
    width=\linewidth,
    height=8cm,
    ymin=0, ymax=40,
    ylabel={Pourcentage (\%)},
    xlabel={Médias principaux},
    symbolic x coords={Radio, Télévision, Réseaux sociaux, Presse écrite, Bouche-à-oreille},
    xtick=data,
    xticklabel style={rotate=30, anchor=east, font=\small},
    nodes near coords={\pgfmathprintnumber\pgfplotspointmeta\%},
    nodes near coords align={vertical},
    title={Répartition des médias principaux (Enquête Bafoussam, 30 pers.)},
    title style={font=\bfseries\small},
    ymajorgrids=true,
    grid style=dashed,
]
\addplot[popblue, fill=popblueL] coordinates {
    (Radio, 26.7)
    (Télévision, 33.3)
    (Réseaux sociaux, 30.0)
    (Presse écrite, 6.7)
    (Bouche-à-oreille, 3.3)
};
\end{axis}
\end{tikzpicture}
\legende{Diagramme en barres : La Télévision arrive en tête (33,3\%), suivie de près par les Réseaux sociaux (30,0\%). La presse écrite et le bouche-à-oreille sont marginaux.}
\end{popfigure}

\begin{popfigure}
\begin{tikzpicture}
\def\angle{0}
\def\radius{3}
\foreach \percent/\name/\color in {16.7/Élevée/popgreen!70, 50.0/Moyenne/poporange!70, 33.3/Faible/popred!70} {
  \pgfmathsetmacro{\newangle}{\angle + \percent * 3.6}
  \draw[fill=\color] (0,0) -- ++(\angle:\radius) arc (\angle:\newangle:\radius) -- cycle;
  \pgfmathsetmacro{\midangle}{(\angle + \newangle)/2}
  \node at (\midangle:\radius*0.65) {\textbf{\name} \percent\%};
  \global\let\angle\newangle
}
\node[above] at (0,3.5) {\textbf{Niveau de confiance dans le média principal}};
\node[below] at (0,-3.5) {\textit{Total enquêtés : 30 | Source : Enquête classe 1re Tech}};
\end{tikzpicture}
\legende{Diagramme circulaire : La confiance « Moyenne » domine (50,0\%). Seuls 16,7\% accordent une confiance « Élevée », tandis qu'un tiers (33,3\%) se dit méfiant.}
\end{popfigure}

\textbf{Étape 3 : Identification du média le plus influent et analyse croisée}

\begin{center}
\begin{tabular}{|l|c|c|c|}\hline
\textbf{Média} & \textbf{Audience (\%)} & \textbf{Rang Audience} & \textbf{Confiance moyenne estimée*} \\ \hline
Télévision & 33,3 & 1 & Moyenne à Faible \\ \hline
Réseaux sociaux & 30,0 & 2 & Faible \\ \hline
Radio & 26,7 & 3 & Moyenne à Élevée \\ \hline
Presse écrite & 6,7 & 4 & Élevée \\ \hline
Bouche-à-oreille & 3,3 & 5 & Variable \\ \hline
\end{tabular}
\end{center}
\emph{*Estimation basée sur la répartition globale de la confiance (50\% Moyenne, 33\% Faible) et la littérature sur la crédibilité perçue : la Radio (proximité, langue locale) et la Presse écrite (professionnalisme) inspirent généralement plus de confiance que la TV (perçue comme partisane) et les Réseaux sociaux (non vérifiés).}

\textbf{Conclusion de l'analyse :} La \textbf{Télévision} est le média \textbf{le plus influent par l'audience} (premier rang, 33,3\%). Cependant, l'influence \emph{cognitive} (capacité à faire croire) est affaiblie par une confiance globalement moyenne/faible. Les \textbf{Réseaux sociaux} (2e audience, 30\%) constituent le \textbf{vecteur d'influence le plus risqué} car ils combinent forte audience et très faible confiance (propagation virale non vérifiée). La \textbf{Radio} (3e audience, 26,7\%) conserve un capital confiance supérieur, la rendant influente sur le long terme pour l'ancrage des valeurs.

\textbf{Étape 4 : Étude de cas \cle{fake news} réelle — « Recrutement fictif Fonction Publique 2024 »}

\begin{itemize}
    \item \textbf{Contenu :} Message viral sur WhatsApp/Facebook : « Le MINFOPRA lance un recrutement exceptionnel de 5000 agents (Catégories A, B, C, D). Inscription via lien \texttt{bit.ly/recrut2024cm} avant le 31 mars. Frais de dossier : 15 000 FCFA payables par Mobile Money. »
    \item \textbf{Source initiale :} Compte Facebook anonyme « Emploi Cameroun Officiel » (créé 3 jours avant), imitant le logo du MINFOPRA. Lien raccourci masquant un site \emph{phishing} hébergé à l'étranger.
    \item \textbf{Vecteurs de propagation :} Partagé dans plus de 50 groupes WhatsApp « Étudiants Bafoussam », « Chômeurs Diplômés Ouest », « Parents d'élèves ». Relais par 3 influenceurs locaux (pages > 10k abonnés) sans vérification. Bouche-à-oreille au marché (photocopies de l'annonce vendues 100 F).
    \item \textbf{Démenti officiel :} Communiqué \textbf{N°000001/MINFOPRA/SG/DRL du 12 mars 2024} : « Le Ministère dément formellement tout recrutement en cours. Attention aux arnaques. Seul le site \texttt{www.minfopra.gov.cm} et le journal \emph{Cameroon Tribune} font foi. » Reprise par \textbf{CRTV Info (13h/20h)}, \textbf{DataCheck Cameroun} (article de fact-checking), \textbf{HAC} (alerte aux médias audiovisuels).
    \item \textbf{Impact observé :} Au moins 120 jeunes de l'Ouest signalés victimes (perte 15 000 FCFA + données bancaires). Troubles devant la Délégation Régionale du MINFOPRA (manifestation de victimes). Discrédit temporaire des vrais concours officiels.
\end{itemize}

\textbf{Étape 5 : Extrait du Rapport d'enquête (Modèle de rédaction)}

\textbf{TITRE :} Enquête « Médias et Nous » — Analyse de l'écosystème informationnel à Bafoussam (Mai 2025)

\textbf{1. Introduction}
Dans le cadre du cours d'ECM, unité « Les médias et leur impact sur l'opinion publique », notre groupe « Éthique Média » a réalisé une enquête du 12 au 15 mai 2025 auprès de 30 résidents de l'arrondissement de Bafoussam 1er (Marché Central, Quartier Tamdja, Lycée Technique). Objectif : cartographier les pratiques informationnelles, mesurer la confiance et évaluer la vulnérabilité aux \cle{fake news}.

\textbf{2. Résultats et Analyse}
L'analyse du diagramme en barres (Fig. 1) révèle une \textbf{duopole informationnel} : Télévision (33,3\%) et Réseaux sociaux (30,0\%) captent 63,3\% de l'audience. La Radio (26,7\%) résiste grâce à l'accessibilité (piles, autoradio) et la langue locale (Medumba). La presse écrite (6,7\%) est en déclin structurel (coût, délai).
Le camembert (Fig. 2) montre une \textbf{crise de confiance} : 83,3\% des sondés ont une confiance « Moyenne » ou « Faible ». Paradoxe : on consomme massivement ce qu'on ne croit pas totalement. Cela s'explique par l'absence d'alternative locale fiable et la gratuité perçue des réseaux sociaux.
Le taux de 60\% de personnes ayant cru/partagé une \cle{fake news} confirme une \textbf{vulnérabilité critique}. L'étude de cas du « Recrutement MINFOPRA » illustre le mécanisme : \emph{urgence sociale (chômage) + autorité usurpée (MINFOPRA) + facilité technique (Mobile Money) = arnaque massive}.

\textbf{3. Trois conseils pratiques pour protéger l'opinion publique de Bafoussam}
\begin{enumerate}
    \item \textbf{Instaurer le « Réflexe Source Officielle » (RSO) :} Avant tout partage d'info administrative (concours, bourse, aide, alerte sécurité), \textbf{vérifier systématiquement sur le site \texttt{.gov.cm} de l'institution concernée ou sur la page Facebook \emph{vérifiée} (badge bleu) du Ministère}. Exemple : Pour un recrutement, aller sur \texttt{minfopra.gov.cm} \textbf{avant} de cliquer sur un lien WhatsApp.
    \item \textbf{Utiliser les canaux de signalement institutionnels :} En cas de doute sur une info circulant dans un groupe WhatsApp de quartier, \textbf{transférer le message au 1510 (Centre d'Appel Citoyen / Police) ou au numéro vert de l'ANTIC (Agence Nationale des Technologies de l'Information et de la Communication) : 8200}. Ne pas laisser l'infox circuler 24h sans alerter.
    \item \textbf{Soutenir et écouter la Radio Communautaire (Radio Medumba 94.0 FM) comme « Filtre de confiance » :} La radio locale, gérée par des professionnels connus du quartier, applique la déontologie (recoupement, droit de réponse). \textbf{S'engager à écouter le journal de 12h30 ou 19h} et à y signaler les rumeurs entendues au marché pour qu'elles soient traitées en direct (rubrique « Vrai ou Faux »). Cela renforce un média de proximité responsable, rempart contre la désinformation anonyme.
\end{enumerate}

\textbf{4. Conclusion}
Cette enquête nous a fait toucher du doigt la fragilité de notre opinion publique : sur-informée mais sous-vérifiée. En tant que futurs techniciens et citoyens, nous nous engageons à devenir des \cle{veilleurs numériques} dans nos familles et groupes WhatsApp, en appliquant la méthode « Source — Date — Auteur — Contexte » avant tout partage.
\tcblower
\textbf{Commentaire pédagogique du corrigé :}
\begin{itemize}
    \item Les calculs de pourcentages sont exacts (somme = 100\% ou 99,9\% par arrondi).
    \item Les graphiques respectent les normes : titres, unités, légende, source.
    \item L'analyse croisée (Audience vs Confiance) dépasse la simple lecture du classement : elle introduit la nuance « influence par contrainte » vs « influence par adhésion ».
    \item L'étude de cas est \textbf{ancrée, datée, sourcée} (MINFOPRA, DataCheck, HAC, ANTIC) et non inventée.
    \item Les 3 conseils sont \textbf{opérationnels} (numéros verts, sites .gov, média local nommé), \textbf{juridiquement fondés} (signalement police/ANTIC) et \textbf{socialement ancrés} (radio communautaire).
    \item La rédaction utilise le vocabulaire du programme : \cle{agenda-setting}, \cle{framing}, \cle{droit de réponse}, \cle{HAC}, \cle{CNC}, \cle{cybercriminalité}.
\end{itemize}
\end{exoresolu}

\courssub{Bilan de l'activité}
Cette activité d'intégration a permis de \textbf{valider l'acquisition des savoirs et savoir-faire} de l'unité « Les médias et leur impact sur l'opinion publique » dans une démarche \textbf{projet authentique} :

\begin{aretenir}[Compétences certifiées]
\begin{itemize}
    \item \textbf{Savoir (Connaissances) :} Les élèves distinguent clairement les types de médias, comprennent la formation de l'opinion (agenda-setting, spirale du silence), identifient les mécanismes de la désinformation (fabrication, viralité, biais cognitifs) et citent les acteurs de la régulation au Cameroun (HAC, CNC, MINPOSTEL, ANTIC, Police judiciaire).
    \item \textbf{Savoir-faire (Méthodes) :} Ils maîtrisent le cycle d'une enquête d'opinion : élaboration d'outil, échantillonnage raisonné, dépouillement, traitement statistique (fréquences, \%), représentation graphique normalisée (barres, camembert), analyse croisée de variables, étude de cas documentée, rédaction de rapport technique et argumenté.
    \item \textbf{Savoir-être (Attitudes) :} Ils développent l'\emph{esprit critique} (doute méthodique, vérification), la \emph{responsabilité numérique} (ne pas partager sans vérifier, signaler), le \emph{civisme actif} (connaissance des recours légaux, participation à l'assainissement de l'espace public) et la \emph{communication argumentée} (restitution orale, débat contradictoire).
\end{itemize}
\end{aretenir}

\begin{savaistu}[Ancrage camerounais et prolongements]
\begin{itemize}
    \item \textbf{Cadre légal :} La Loi n°2010/012 du 21 décembre 2010 relative à la cybercriminalité (art. 7 : « propagation de fausses nouvelles ») et la Loi n°2019/018/019 sur la communication sociale encadrent juridiquement les faits étudiés.
    \item \textbf{Acteurs réels :} \emph{DataCheck Cameroun}, \emph{StopIntox} (plateformes de fact-checking), \emph{ANTIC} (CERT-CM pour alertes cybersécurité), \emph{HAC} (sanctions aux médias non conformes).
    \item \textbf{Prolongement technique :} Les élèves des séries STT (Secrétariat-Bureautique), GEN (Gestion) ou TIC peuvent numériser le questionnaire (Google Forms / KoboToolbox), automatiser les graphiques (Excel / Python matplotlib) et publier le rapport sur le site de l'établissement ou le journal scolaire.
\end{itemize}
\end{savaistu}

\newpage

\coursec{Méthodes et exercices résolus}

\courssub{Savoir-faire 1 : Appliquer les notions de l'unité à des situations concrètes de la vie courante}

\begin{methode}[Analyser une situation concrète liée aux médias et à l'opinion publique]
Pour traiter une situation de vie (atelier, quartier, réseaux sociaux, entreprise) au Probatoire/Bac technique, suivre une démarche rigoureuse :
\begin{enumerate}
\item \textbf{Lire attentivement la situation} et souligner les faits : qui agit ? par quel média ? avec quel effet sur le public ?
\item \textbf{Identifier la notion du cours mobilisée} : médias (presse écrite, radio, télévision, médias numériques), opinion publique, manipulation, publicité, protection de l'opinion.
\item \textbf{Rappeler la définition exacte} de la notion (une phrase, apprise par cœur).
\item \textbf{Relier la définition aux faits} de la situation : citer précisément l'élément du texte qui illustre la notion.
\item \textbf{Dégager l'impact ou le problème} : influence positive (information, éducation) ou négative (rumeur, désinformation, manipulation).
\item \textbf{Conclure} par une phrase rédigée qui répond directement à la question posée.
\end{enumerate}
Réflexe : toujours passer du \cle{concret} (les faits) à la \cle{notion} (le cours), puis revenir au concret pour conclure.
\end{methode}

\begin{exoresolu}[Rumeur sur les réseaux sociaux dans un lycée technique]
Énoncé. Dans un lycée technique de Douala, une vidéo truquée circule sur WhatsApp : elle accuse le proviseur de détourner les fonds de la cantine. En quelques heures, des centaines d'élèves et de parents partagent la vidéo. Une enquête administrative démontre ensuite que l'accusation est fausse, mais la réputation du proviseur reste entachée.
\begin{enumerate}
\item Identifie le type de média concerné et la notion du cours illustrée par cette situation.
\item Montre, à partir de deux faits précis, l'impact de ce média sur l'opinion publique.
\end{enumerate}
\tcblower
Solution détaillée.
\begin{enumerate}
\item \textbf{Notion et média.} Le média concerné est un \cle{média numérique} (réseau social, messagerie WhatsApp), qui appartient aux médias de masse modernes. La notion du cours illustrée est \cle{l'impact des médias sur l'opinion publique}, ici sous sa forme négative : la \emph{désinformation} (diffusion d'une fausse information) et la \emph{manipulation de l'opinion}.
\item \textbf{Application aux faits.} L'opinion publique est l'ensemble des attitudes et des jugements partagés par un groupe de personnes sur une question d'intérêt général. Deux faits du texte montrent l'impact du média sur cette opinion :
\begin{itemize}
\item \og En quelques heures, des centaines d'élèves et de parents partagent la vidéo \fg{} : cela illustre la \textbf{rapidité et l'ampleur} de la diffusion propre aux médias numériques, qui façonnent très vite un jugement collectif.
\item \og La réputation du proviseur reste entachée \fg{} malgré l'enquête qui le disculpe : cela montre la \textbf{persistence} de l'opinion une fois formée ; la correction de l'information arrive trop tard et peine à effacer la première impression.
\end{itemize}
\end{enumerate}
\textbf{Conclusion.} Cette situation montre que les médias numériques, par leur vitesse de diffusion, peuvent imposer une fausse image à l'opinion publique ; d'où la nécessité de vérifier les informations avant de les partager et de protéger l'opinion contre la désinformation.
\end{exoresolu}

\courssub{Savoir-faire 2 : Rédiger et justifier une démarche, une réponse ou une conclusion}

\begin{methode}[Rédiger une réponse argumentée et justifiée en ECM]
Une réponse au Bac technique doit être structurée et justifiée :
\begin{enumerate}
\item \textbf{Annoncer clairement la réponse} dès la première phrase (ne pas tourner autour du pot).
\item \textbf{Justifier par un argument de cours} : définition, loi, principe (ex. liberté de la presse, code de déontologie journalistique).
\item \textbf{Justifier par un fait} tiré de la situation ou de la vie courante camerounaise.
\item \textbf{Enchaîner les idées} avec des connecteurs : d'abord, ensuite, en effet, cependant, par conséquent, enfin.
\item \textbf{Conclure} en reformulant la réponse et en ouvrant éventuellement sur une attitude citoyenne (esprit critique, vérification des sources).
\item \textbf{Relire} : orthographe, phrases complètes (sujet + verbe), vocabulaire du cours employé correctement.
\end{enumerate}
Formule à retenir : \cle{réponse} + \cle{argument de cours} + \cle{preuve factuelle} + \cle{conclusion}.
\end{methode}

\begin{exoresolu}[Faut-il vérifier une information avant de la partager ?]
Énoncé. Ton camarade de classe partage sur Facebook une annonce selon laquelle le gouvernement offre gratuitement des ordinateurs portables à tous les élèves qui s'inscrivent sur un site. Il t'encourage à faire de même. Rédige un texte de 8 à 12 lignes dans lequel tu réponds à la question : faut-il vérifier une information avant de la partager ? Justifie ta réponse.
\tcblower
Solution détaillée (modèle de rédaction).

\emph{Oui, il faut toujours vérifier une information avant de la partager.} D'abord, les médias numériques diffusent les messages instantanément et à un très large public : une fausse information partagée peut donc influencer l'opinion publique en quelques minutes, comme le montre la rapidité de propagation des rumeurs sur les réseaux sociaux. Ensuite, la désinformation est un danger réel : l'annonce des ordinateurs gratuits peut être une \emph{escroquerie} visant à collecter les données personnelles des élèves (nom, numéro de téléphone, identifiants). En effet, aucun communiqué officiel du MINESEC ou des médias publics (CRTV, Cameroon Tribune) ne confirme cette offre, ce qui est un indice de fausseté. Cependant, partager sans vérifier, c'est devenir complice de la manipulation de l'opinion et exposer ses proches à un préjudice. Par conséquent, la démarche citoyenne consiste à croiser les sources (sites officiels, médias reconnus), à vérifier l'auteur et la date du message, puis seulement à décider de partager ou non. Enfin, développer l'esprit critique face aux médias est un devoir du citoyen et du consommateur d'information. En conclusion, vérifier avant de partager protège à la fois soi-même, son entourage et l'opinion publique.

\emph{Justification de la démarche} : la réponse est annoncée dès la première phrase ; deux arguments de cours (vitesse de diffusion, désinformation) sont appuyés par des faits (absence de confirmation officielle, risque d'escroquerie) ; les connecteurs structurent le texte ; la conclusion répond à la question.
\end{exoresolu}

\courssub{Savoir-faire 3 : Distinguer les types de médias et leurs rôles}

\begin{methode}[Classer un média et dégager son rôle]
\begin{enumerate}
\item \textbf{Identifier le support} : papier, ondes hertziennes, écran, Internet.
\item \textbf{Classer} dans la bonne catégorie : presse écrite (journaux, magazines), radio, télévision, médias numériques (réseaux sociaux, blogs, sites d'information).
\item \textbf{Relever les caractéristiques propres} : la presse écrite permet l'analyse approfondie ; la radio touche les zones rurales et les personnes non alphabétisées ; la télévision combine image et son ; les médias numériques sont rapides, interactifs mais peu contrôlés.
\item \textbf{Attribuer les rôles} : informer, éduquer, divertir, former l'opinion, servir de contre-pouvoir (dénoncer les abus).
\item \textbf{Nuancer} : chaque média a des limites (coût, accès, fiabilité).
\end{enumerate}
\end{methode}

\begin{exoresolu}[Le bon média pour la bonne campagne]
Énoncé. Le ministère de la Santé publique veut lancer une campagne de vaccination contre le choléra dans les villages de l'Extrême-Nord du Cameroun, où une partie de la population ne sait ni lire ni écrire et où l'accès à Internet est faible.
\begin{enumerate}
\item Quel type de média conseilles-tu en priorité ? Justifie par deux caractéristiques.
\item Cite deux rôles que ce média jouera dans cette campagne.
\end{enumerate}
\tcblower
Solution détaillée.
\begin{enumerate}
\item \textbf{Média conseillé : la radio.} Justification par deux caractéristiques :
\begin{itemize}
\item La radio est un média \textbf{oral} : elle ne suppose pas de savoir lire, donc elle atteint les populations non alphabétisées, contrairement à la presse écrite.
\item La radio a une \textbf{large couverture géographique} et un coût d'accès faible (un simple poste à piles suffit) : elle touche les zones rurales où Internet et la télévision sont peu accessibles. De plus, elle peut émettre dans les \textbf{langues locales} (fulfuldé, etc.), ce qui renforce la compréhension du message.
\end{itemize}
\item \textbf{Deux rôles joués par la radio dans cette campagne :}
\begin{itemize}
\item \textbf{Informer} : annoncer les dates, les lieux et l'importance de la vaccination.
\item \textbf{Éduquer / former l'opinion} : expliquer les gestes d'hygiène, combattre les rumeurs sur les vaccins et convaincre la population de se faire vacciner.
\end{itemize}
\end{enumerate}
\textbf{Conclusion.} Le choix d'un média dépend du public visé et de ses conditions d'accès ; ici, la radio est le média le plus adapté pour informer et éduquer les populations rurales de l'Extrême-Nord.
\end{exoresolu}

\courssub{Savoir-faire 4 : Expliquer la formation de l'opinion publique et l'impact des médias}

\begin{methode}[Expliquer comment les médias influencent l'opinion publique]
\begin{enumerate}
\item \textbf{Définir l'opinion publique} : ensemble des jugements, attitudes et sentiments partagés par une population sur une question d'intérêt général.
\item \textbf{Décrire les mécanismes d'influence} : sélection des sujets (les médias imposent les thèmes de débat), répétition des messages, mise en avant de certaines images, cadrage des informations.
\item \textbf{Distinguer les effets positifs} (information, éducation, démocratie, contre-pouvoir) \textbf{des effets négatifs} (rumeurs, désinformation, propagande, manipulation, atteinte à la réputation).
\item \textbf{Illustrer chaque effet par un exemple} camerounais ou de la vie courante.
\item \textbf{Conclure} sur la nécessité de l'esprit critique et de la régulation.
\end{enumerate}
\end{methode}

\begin{exoresolu}[Les médias, miroir ou fabricant de l'opinion ?]
Énoncé. À l'approche d'une élection au Cameroun, les chaînes de télévision multiplient les débats politiques, les radios diffusent des spots des candidats et les réseaux sociaux relaient des sondages. Certains citoyens affirment : \og les médias ne reflètent pas l'opinion, ils la fabriquent \fg.
\begin{enumerate}
\item Définis l'opinion publique.
\item En t'appuyant sur deux mécanismes, montre comment les médias peuvent \og fabriquer \fg{} l'opinion.
\item Nuance ce jugement en citant un rôle positif des médias en période électorale.
\end{enumerate}
\tcblower
Solution détaillée.
\begin{enumerate}
\item \textbf{Définition.} L'opinion publique est l'ensemble des jugements, des attitudes et des sentiments qu'une population partage, à un moment donné, sur une question d'intérêt général (politique, santé, sécurité, etc.).
\item \textbf{Deux mécanismes de fabrication de l'opinion :}
\begin{itemize}
\item \textbf{La sélection et l'agenda} : en choisissant les sujets traités (débats, sondages), les médias imposent les thèmes dont parle la société ; ce qui n'est pas montré semble ne pas exister. Ici, les sondages répétés peuvent faire croire qu'un candidat est déjà vainqueur.
\item \textbf{La répétition et le cadrage} : la diffusion répétée des mêmes spots et des mêmes images finit par paraître vraie et oriente le jugement des électeurs ; présenter un candidat sous un jour favorable ou défavorable influence l'opinion sans que celle-ci s'en rende compte.
\end{itemize}
\item \textbf{Nuance.} Les médias jouent aussi un rôle positif : ils \textbf{informent} les électeurs sur les programmes, assurent le \textbf{débat démocratique} et exercent un \textbf{contre-pouvoir} en dénonçant les fraudes ou les abus. Ils ne fabriquent donc pas toujours l'opinion : ils peuvent l'éclairer.
\end{enumerate}
\textbf{Conclusion.} Les médias orientent fortement l'opinion publique par la sélection et la répétition des messages, mais ils restent indispensables à la vie démocratique ; le citoyen doit donc développer son esprit critique.
\end{exoresolu}

\courssub{Savoir-faire 5 : Mobiliser les mécanismes de protection de l'opinion publique au Cameroun}

\begin{methode}[Citer et expliquer les protections de l'opinion publique au Cameroun]
\begin{enumerate}
\item \textbf{Identifier le danger} contre lequel on veut protéger l'opinion : désinformation, diffamation, propagande, publicité mensongère, atteinte à la vie privée.
\item \textbf{Citer le cadre juridique} : la Constitution du 18 janvier 1996 (liberté d'expression et de la presse), la loi n° 90/052 du 19 décembre 1990 sur la communication sociale, le code pénal (répression de la diffamation et des fausses nouvelles), la loi n° 2010/012 du 21 décembre 2010 sur la cybersécurité et la cybercriminalité.
\item \textbf{Citer les institutions et acteurs} : le Conseil National de la Communication (CNC), les organes de régulation, la déontologie journalistique, les associations de consommateurs.
\item \textbf{Expliquer le rôle de chaque protection} en une phrase.
\item \textbf{Rappeler la part du citoyen} : vérification des sources, esprit critique, éducation aux médias.
\end{enumerate}
Réflexe : une protection = un texte ou une institution + sa fonction. Ne jamais citer un texte sans dire à quoi il sert.
\end{methode}

\begin{exoresolu}[Protéger l'opinion face aux fausses nouvelles]
Énoncé. Pendant une épidémie, des messages audio diffusés sur les réseaux sociaux affirment qu'un médicament traditionnel guérit la maladie et que les hôpitaux empoisonnent les malades. Plusieurs familles refusent alors d'emmener leurs proches à l'hôpital.
\begin{enumerate}
\item Quel danger pour l'opinion publique cette situation révèle-t-elle ?
\item Cite deux mécanismes de protection de l'opinion publique au Cameroun et explique comment chacun peut s'appliquer ici.
\end{enumerate}
\tcblower
Solution détaillée.
\begin{enumerate}
\item \textbf{Danger.} Il s'agit de la \cle{désinformation} (diffusion de fausses nouvelles) qui manipule l'opinion publique et met en danger la santé des populations, puisque des malades renoncent aux soins hospitaliers.
\item \textbf{Deux mécanismes de protection :}
\begin{itemize}
\item \textbf{Le cadre juridique} : le code pénal camerounais réprime la diffusion de fausses nouvelles, et la loi n° 2010/012 du 21 décembre 2010 sur la cybersécurité et la cybercriminalité sanctionne la propagation de fausses informations par voie électronique. Application : les auteurs des messages audio peuvent être recherchés et poursuivis, ce qui dissuade la désinformation.
\item \textbf{Le Conseil National de la Communication (CNC) et la déontologie} : le CNC veille au respect des règles de la communication sociale et peut mettre en garde les médias qui relaient de fausses informations ; les journalistes professionnels, tenus par la déontologie (vérification des sources, honnêteté), peuvent publier des démentis et des informations vérifiées du ministère de la Santé. Application : des communiqués officiels relayés par la CRTV et la presse sérieuse corrigent la rumeur et rétablissent la vérité dans l'opinion.
\end{itemize}
\end{enumerate}
\textbf{Conclusion.} La protection de l'opinion publique au Cameroun repose sur les textes de loi, les institutions de régulation et la déontologie, complétés par l'esprit critique des citoyens qui doivent vérifier les informations avant d'y croire et de les partager.
\end{exoresolu}

\courssub{Savoir-faire 6 : Analyser la publicité dans les médias (TD3)}

\begin{methode}[Analyser un message publicitaire]
\begin{enumerate}
\item \textbf{Définir la publicité} : technique de communication payante, diffusée par un média, visant à influencer le comportement du public en faveur d'un produit, d'un service ou d'une idée.
\item \textbf{Identifier les éléments du message} : l'annonceur (qui paie ?), la cible (à qui ?), le support (quel média ?), le slogan et les images.
\item \textbf{Repérer les techniques d'influence} : répétition, témoignage d'une vedette, promesses exagérées, jeux sur les émotions, prix \og exceptionnels \fg.
\item \textbf{Distinguer publicité honnête et publicité mensongère} : une publicité est mensongère lorsqu'elle attribue au produit des qualités qu'il n'a pas.
\item \textbf{Évaluer l'impact} sur l'opinion et le consommateur, puis proposer une attitude responsable (comparer, lire les étiquettes, se méfier des promesses miracles).
\end{enumerate}
\end{methode}

\begin{exoresolu}[TD3 : Une publicité pour une boisson \og miracle \fg]
Énoncé. Une entreprise diffuse à la télévision et sur Facebook une publicité pour une boisson énergisante. Une célèbre chanteuse camerounaise y déclare : \og Avec cette boisson, tu seras toujours le premier de la classe et jamais fatigué ! \fg{} Le spot est répété dix fois par jour. De nombreux élèves de Première technique achètent alors cette boisson.
\begin{enumerate}
\item Identifie l'annonceur, la cible et les supports de cette publicité.
\item Releve deux techniques d'influence utilisées.
\item Cette publicité est-elle honnête ? Justifie et indique l'attitude responsable du consommateur.
\end{enumerate}
\tcblower
Solution détaillée.
\begin{enumerate}
\item \textbf{Éléments du message.} L'\textbf{annonceur} est l'entreprise qui fabrique la boisson énergisante (c'est elle qui paie la diffusion). La \textbf{cible} est la jeunesse, en particulier les élèves. Les \textbf{supports} sont la télévision et un média numérique (Facebook).
\item \textbf{Deux techniques d'influence :}
\begin{itemize}
\item \textbf{Le témoignage d'une vedette} : la chanteuse célèbre sert de modèle ; les jeunes achètent par identification et admiration, non par jugement sur le produit.
\item \textbf{La répétition} (dix passages par jour) : le message s'imprime dans la mémoire et finit par paraître vrai ; s'y ajoute la \textbf{promesse exagérée} (\og toujours le premier de la classe \fg) qui flatte les désirs des élèves.
\end{itemize}
\item \textbf{Honnêteté de la publicité.} Cette publicité est \textbf{mensongère}, car aucune boisson ne peut garantir d'être premier de la classe ni supprimer la fatigue : elle attribue au produit des qualités qu'il ne possède pas. Une telle publicité trompe le consommateur et manipule l'opinion ; elle est contraire à la protection du consommateur. L'\textbf{attitude responsable} consiste à ne pas croire les promesses miracles, à lire la composition du produit, à comparer les informations et à se rappeler que la réussite scolaire dépend du travail, non d'une boisson.
\end{enumerate}
\textbf{Conclusion.} La publicité est un moyen légitime d'information commerciale, mais elle devient condamnable lorsqu'elle est mensongère ; le citoyen-consommateur doit exercer son esprit critique face aux techniques d'influence des médias.
\end{exoresolu}

\begin{attention}[Les 5 erreurs qui coûtent des points au Bac]
\begin{enumerate}
\item \textbf{Confondre média et support technique} : dire \og WhatsApp est une opinion \fg{} ou mélanger presse écrite et réseaux sociaux. Apprends la typologie exacte (presse écrite, radio, télévision, médias numériques) et cite le bon média pour la bonne situation.
\item \textbf{Donner une définition vague de l'opinion publique} : écrire \og c'est ce que les gens pensent \fg{} est insuffisant. La définition exigée mentionne les jugements et attitudes \emph{partagés} par une population sur une question \emph{d'intérêt général}.
\item \textbf{Citer un texte de loi sans expliquer son rôle} : écrire \og la loi de 1990 \fg{} seul ne rapporte rien. Chaque protection citée (Constitution de 1996, loi n° 90/052 du 19 décembre 1990, loi de 2010 sur la cybercriminalité, CNC) doit être suivie de sa fonction concrète.
\item \textbf{Répondre sans justifier} : au Bac technique, une affirmation sans argument de cours ni fait précis perd la moitié des points. Applique toujours la formule : réponse + argument + preuve + conclusion.
\item \textbf{Traiter la publicité comme toujours négative ou toujours positive} : le correcteur attend une analyse nuancée (information commerciale légitime \emph{versus} publicité mensongère) et l'identification précise des techniques d'influence (répétition, vedette, promesse exagérée), pas un simple jugement moral.
\end{enumerate}
\end{attention}

\newpage

\coursec{Exercices}

\courssub{Je vérifie mes connaissances}

\begin{exercice}[QCM : notions et institutions]
Pour chaque question, une seule réponse est exacte. Recopie la lettre correspondante.
\begin{enumerate}
\item Un média est :
\begin{enumerate}
\item[a)] une opinion partagée par la majorité des citoyens ;
\item[b)] un moyen de diffusion de l'information vers un large public ;
\item[c)] un organe de contrôle de la presse.
\end{enumerate}
\item L'opinion publique est :
\begin{enumerate}
\item[a)] l'avis du gouvernement sur les questions de société ;
\item[b)] l'ensemble des idées, attitudes et jugements partagés par une population sur un sujet donné ;
\item[c)] le contenu des journaux télévisés.
\end{enumerate}
\item Une \emph{fake news} est :
\begin{enumerate}
\item[a)] une information fausse créée ou relayée pour tromper le public ;
\item[b)] une information vraie mais désagréable ;
\item[c)] une publicité mensongère uniquement.
\end{enumerate}
\item Au Cameroun, l'organe de régulation de la communication sociale est :
\begin{enumerate}
\item[a)] le MINCOM ;
\item[b)] le Conseil National de la Communication (CNC) ;
\item[c)] l'Assemblée nationale.
\end{enumerate}
\item Le droit de réponse permet à toute personne mise en cause par un média de :
\begin{enumerate}
\item[a)] faire fermer le média ;
\item[b)] exiger la publication de sa réaction dans le même média ;
\item[c)] refuser toute interview.
\end{enumerate}
\end{enumerate}
\end{exercice}

\begin{exercice}[Vrai ou faux ? Justifie ta réponse]
Réponds par Vrai ou Faux, puis justifie en deux ou trois phrases.
\begin{enumerate}
\item Les réseaux sociaux (WhatsApp, Facebook, TikTok) sont des médias au même titre que la radio et la télévision.
\item L'opinion publique se forme uniquement à partir des informations diffusées par les médias.
\item La liberté de la presse est garantie par la Constitution camerounaise du 18 janvier 1996.
\item Relayer une rumeur sur WhatsApp sans la vérifier n'engage aucune responsabilité.
\end{enumerate}
\end{exercice}

\begin{exercice}[Je définis]
Définis clairement, en deux ou trois phrases chacune, les notions suivantes : média ; opinion publique ; presse écrite ; média audiovisuel ; rumeur ; droit de réponse.
\end{exercice}

\begin{exercice}[Je classe et je calcule]
\begin{enumerate}
\item Classe les supports suivants dans la catégorie qui convient (presse écrite, média audiovisuel, média numérique) : \emph{Cameroon Tribune} ; la CRTV ; un site d'information en ligne ; la radio Magic FM ; une page Facebook d'actualités ; un journal scolaire.
\item Dans un quartier de Yaoundé, une enquête porte sur 500 personnes : 320 s'informent principalement par les réseaux sociaux, 100 par la radio, 60 par la télévision et 20 par la presse écrite. Calcule le pourcentage de personnes correspondant à chaque canal d'information.
\end{enumerate}
\end{exercice}

\courssub{Je m'entraîne}

\begin{exercice}[Rumeur sur WhatsApp]
Dans un quartier de Douala, un message vocal diffusé sur WhatsApp affirme que l'eau du robinet est empoisonnée. En quelques heures, les habitants stockent de l'eau et certains commerces ferment.
\begin{enumerate}
\item Explique le mécanisme de propagation de cette rumeur (rapidité, absence de vérification, effet de peur).
\item Identifie deux risques concrets pour la vie du quartier.
\item Propose deux réflexes de vérification qu'un citoyen responsable aurait dû adopter avant de partager le message.
\end{enumerate}
\end{exercice}

\begin{exercice}[Analyse d'un spot publicitaire]
Un spot télévisé de 30 secondes vante une boisson énergisante : un jeune footballeur célèbre boit le produit avant de marquer le but de la victoire ; des jeunes dansent ; le slogan affirme : « Bois X, deviens un champion ! »
\begin{enumerate}
\item Dégage la cible de cette publicité.
\item Relève deux techniques de persuasion utilisées (témoignage de star, promesse exagérée, musique et émotion...).
\item Cette publicité respecte-t-elle la déontologie publicitaire ? Justifie ta réponse.
\end{enumerate}
\end{exercice}

\begin{exercice}[Deux sources contradictoires]
À propos de l'ouverture d'un nouveau marché à Bafoussam, un journal publie un article signé, daté, citant le maire et des chiffres officiels ; un post Facebook anonyme affirme le contraire sans aucune preuve.
\begin{enumerate}
\item Indique quelle source est la plus fiable en justifiant avec au moins trois critères (auteur identifiable, sources citées, recoupement possible).
\item Rédige la démarche de vérification complète qu'un élève de Première technique devrait suivre avant de croire ou de partager une information.
\end{enumerate}
\end{exercice}

\begin{exercice}[La protection de l'opinion publique]
\begin{enumerate}
\item Cite deux institutions ou textes qui protègent l'opinion publique au Cameroun et précise le rôle de chacun.
\item Explique en quoi la déontologie du journaliste (vérification des sources, respect de la vie privée, droit de réponse) protège les citoyens.
\item Donne un exemple concret de sanction possible contre un média qui diffuse de fausses informations.
\end{enumerate}
\end{exercice}

\courssub{Je me prépare au Bac}

\begin{exercice}[Exercice documentaire type Probatoire]
\textbf{Document 1} (extrait d'un article de presse) : « Lors de l'inauguration du lycée technique de Nkongsamba, le ministre a annoncé la création de 200 places d'internat. L'information a été confirmée par le préfet et publiée au bulletin officiel de la région. »

\textbf{Document 2} (post Facebook) : « Scandale ! Le lycée technique de Nkongsamba n'ouvrira jamais, les fonds ont disparu. Partagez massivement ! » (aucun auteur identifié, aucune source citée).
\begin{enumerate}
\item Présente chaque document (nature, source, degré de fiabilité).
\item À l'aide d'au moins trois critères de vérification, indique quel document donne l'information fiable et justifie ta réponse.
\item Explique pourquoi le Document 2 peut être qualifié de \emph{fake news} et quels dangers sa diffusion représente pour l'opinion publique.
\item Décris en quatre étapes la démarche de vérification d'une information reçue sur les réseaux sociaux.
\item Cite l'institution camerounaise compétente pour sanctionner les dérives des médias et précise son rôle.
\end{enumerate}
\end{exercice}

\begin{exercice}[Étude de cas : la rumeur du marché]
À Douala, un message WhatsApp accuse les commerçants d'un marché de vendre de la viande avariée. En deux jours, la rumeur se propage, les clients désertent le marché et plusieurs boutiques ferment. L'enquête sanitaire révèle ensuite que l'information était fausse.
\begin{enumerate}
\item Définis les notions de rumeur et d'opinion publique, puis montre comment la rumeur a influencé l'opinion publique dans ce cas.
\item Analyse le mécanisme de propagation de la rumeur (vitesse des réseaux sociaux, émotion, absence de vérification).
\item Identifie les responsabilités engagées : celle de l'auteur du message, celle des personnes qui l'ont relayé, celle des consommateurs.
\item Présente deux mesures de protection applicables au Cameroun (loi sur la communication sociale, rôle du CNC, sanctions pénales contre la diffusion de fausses nouvelles).
\item Propose trois actions concrètes pour éviter qu'une telle situation ne se reproduise dans ton quartier.
\end{enumerate}
\end{exercice}

\begin{exercice}[Question de synthèse type Baccalauréat]
Sujet : « Les médias sont-ils un danger ou une chance pour la démocratie camerounaise ? »

Rédige un développement structuré d'environ deux pages comportant :
\begin{enumerate}
\item une introduction avec une mise en contexte, une définition des termes clés (médias, opinion publique, démocratie) et une problématique claire ;
\item un développement en deux parties : d'abord les dangers des médias (désinformation, manipulation de l'opinion, rumeurs), ensuite leurs apports (information des citoyens, éducation civique, contre-pouvoir), chaque argument étant illustré d'un exemple camerounais ;
\item une conclusion équilibrée qui répond à la problématique et souligne le rôle de la régulation (CNC, lois) et de l'éducation aux médias.
\end{enumerate}
\end{exercice}

\newpage

\coursec{Corrigés des exercices}

\coursec{Leçon 03 : Les médias et leur impact sur l'opinion publique}
\courssub{Corrigés des exercices d'application et de synthèse}

\begin{corrige}[Exercice 1 — QCM : notions et institutions]
\begin{enumerate}
\item \textbf{Réponse b)} : un \cle{média} est un moyen de diffusion de l'information vers un large public (radio, télévision, presse écrite, médias numériques). La réponse a) définit l'\cle{opinion publique} et la réponse c) définit un organe de régulation comme le \cle{CNC}.
\item \textbf{Réponse b)} : l'\cle{opinion publique} est l'ensemble des idées, attitudes et jugements partagés par une population sur un sujet donné. Elle n'est ni l'avis du gouvernement ni le simple contenu des journaux.
\item \textbf{Réponse a)} : une \cle{fake news} est une information fausse créée ou relayée dans le but de tromper le public. Une information vraie mais désagréable reste une information ; la publicité mensongère n'est qu'un cas particulier.
\item \textbf{Réponse b)} : le \cle{Conseil National de la Communication (CNC)} est l'organe de régulation de la communication sociale au Cameroun. Le \cle{MINCOM} est le ministère de tutelle et l'Assemblée nationale vote les lois.
\item \textbf{Réponse b)} : le \cle{droit de réponse} permet à toute personne mise en cause par un média d'exiger la publication de sa réaction dans ce même média. Il ne donne pas le pouvoir de fermer le média.
\end{enumerate}
\textbf{Points de vigilance} : ne pas confondre média (support de diffusion) et opinion publique (contenu des idées partagées) ; ne pas confondre l'organe de régulation (CNC) avec le ministère de tutelle (MINCOM).
\end{corrige}

\begin{corrige}[Exercice 2 — Vrai ou faux ? Justifie ta réponse]
\begin{enumerate}
\item \textbf{Vrai.} Les réseaux sociaux (WhatsApp, Facebook, TikTok) sont des \cle{médias numériques} : ils diffusent de l'information vers un large public, comme la radio et la télévision. Ils appartiennent à la même famille de moyens de communication de masse, même si leur fonctionnement est interactif et décentralisé.
\item \textbf{Faux.} L'\cle{opinion publique} se forme aussi à partir de l'éducation reçue, des discussions en famille et au quartier, de la religion, des expériences personnelles et du travail. Les médias l'influencent fortement, mais ils ne sont pas sa seule source.
\item \textbf{Vrai.} Le préambule de la \cle{Constitution du 18 janvier 1996} garantit la liberté d'expression, la liberté de la presse et le droit à l'information. Cette liberté s'exerce toutefois dans le respect des lois en vigueur.
\item \textbf{Faux.} Relayer une \cle{rumeur}, même sans l'avoir créée, engage la responsabilité de celui qui la diffuse. La loi camerounaise sur la communication sociale (loi n° 90/052 du 19 décembre 1990) et le code pénal sanctionnent la diffusion de fausses nouvelles ; chaque citoyen doit vérifier une information avant de la partager.
\end{enumerate}
\textbf{Points de vigilance} : une justification courte mais précise est exigée (deux ou trois phrases) ; citer un texte (Constitution de 1996, loi sur la communication sociale) renforce toujours la réponse.
\end{corrige}

\begin{corrige}[Exercice 3 — Je définis]
\begin{itemize}
\item \textbf{\cle{Média}} : moyen technique de diffusion de l'information vers un large public (radio, télévision, presse écrite, médias numériques).
\item \textbf{\cle{Opinion publique}} : ensemble des idées, des attitudes et des jugements partagés par une population sur un sujet donné à un moment donné.
\item \textbf{\cle{Presse écrite}} : ensemble des médias qui diffusent l'information sous forme de textes imprimés ou publiés : journaux, magazines, revues (exemple : \emph{Cameroon Tribune}, un journal scolaire).
\item \textbf{\cle{Média audiovisuel}} : média qui transmet l'information par le son et/ou l'image : la radio (son seul) et la télévision (son et image), comme la CRTV ou Magic FM.
\item \textbf{\cle{Rumeur}} : information non vérifiée, souvent fausse, qui circule de bouche à oreille ou sur les réseaux sociaux et dont la source est inconnue ou incertaine.
\item \textbf{\cle{Droit de réponse}} : droit reconnu à toute personne mise en cause par un média d'exiger la publication ou la diffusion de sa réaction dans ce même média, afin de rétablir la vérité.
\end{itemize}
\textbf{Points de vigilance} : chaque définition doit contenir l'idée essentielle (pour le média : diffusion + large public ; pour la rumeur : non vérifiée + circulation) ; un exemple concret est un plus.
\end{corrige}

\begin{corrige}[Exercice 4 — Je classe et je calcule]
\begin{enumerate}
\item Classement des supports :
\begin{center}
\begin{tabularx}{\linewidth}{|l|X|}\hline
\rowcolor{popblueL} \textbf{Catégorie} & \textbf{Supports} \\ \hline
Presse écrite & \emph{Cameroon Tribune} ; un journal scolaire \\ \hline
Média audiovisuel & la CRTV (télévision) ; la radio Magic FM \\ \hline
Média numérique & un site d'information en ligne ; une page Facebook d'actualités \\ \hline
\end{tabularx}
\end{center}
\item Calcul des pourcentages (total : $320 + 100 + 60 + 20 = 500$ personnes) :
\begin{align*}
\text{Réseaux sociaux} &: \frac{320}{500}\times 100 = 64\ \%\\
\text{Radio} &: \frac{100}{500}\times 100 = 20\ \%\\
\text{Télévision} &: \frac{60}{500}\times 100 = 12\ \%\\
\text{Presse écrite} &: \frac{20}{500}\times 100 = 4\ \%
\end{align*}
\textbf{Vérification} : $64 + 20 + 12 + 4 = 100\ \%$. Le calcul est exact.
\end{enumerate}
\textbf{Points de vigilance} : toujours vérifier que la somme des pourcentages vaut 100 \% ; présenter la formule (part/total $\times$ 100) avant le résultat ; conclure par une phrase (ici, les réseaux sociaux dominent largement l'accès à l'information dans ce quartier).
\end{corrige}

\begin{corrige}[Exercice 5 — Rumeur sur WhatsApp]
\begin{enumerate}
\item \textbf{Mécanisme de propagation} : le message vocal est transféré en quelques secondes de groupe en groupe, ce qui explique la rapidité de la diffusion. Aucun relais ne vérifie l'information auprès d'une source officielle (CAMWATER, mairie, services sanitaires). Enfin, la peur pour la santé pousse chacun à prévenir ses proches, ce qui accélère encore la propagation : c'est l'effet boule de neige de la \cle{rumeur}.
\item \textbf{Deux risques concrets} : la panique et les comportements irrationnels (stockage excessif, achats précipités) ; la perturbation de la vie économique du quartier (fermeture des commerces, pertes de revenus). On peut ajouter la défiance envers les services publics de l'eau.
\item \textbf{Deux réflexes de vérification} : chercher une confirmation auprès d'une source officielle (communiqué de la CAMWATER, de la mairie ou des autorités sanitaires) ; recouper l'information avec au moins deux médias sérieux (radio, télévision, presse) avant tout partage. On peut aussi vérifier l'auteur et la date du message.
\end{enumerate}
\textbf{Points de vigilance} : distinguer clairement les trois éléments du mécanisme (vitesse, absence de vérification, émotion) ; les réflexes proposés doivent être concrets et réalisables par un citoyen ordinaire.
\end{corrige}

\begin{corrige}[Exercice 6 — Analyse d'un spot publicitaire]
\begin{enumerate}
\item \textbf{Cible} : les jeunes, notamment les adolescents et les jeunes adultes amateurs de football et de musique, qui s'identifient au footballeur célèbre et rêvent de réussite sportive.
\item \textbf{Techniques de persuasion} : le \cle{témoignage d'une star} (le footballeur célèbre sert de modèle et transfère son prestige sur le produit) ; la \cle{promesse exagérée} (« Bois X, deviens un champion ! » attribue au produit un pouvoir miraculeux) ; on peut ajouter la musique, la danse et l'émotion, qui créent une ambiance positive associée à la boisson.
\item \textbf{Respect de la déontologie} : cette publicité est \textbf{discutable}. Elle ne respecte pas pleinement la \cle{déontologie} publicitaire, car le slogan fait une promesse mensongère : aucune boisson ne peut transformer un consommateur en champion. Une publicité ne doit pas induire le public en erreur sur les qualités réelles du produit, surtout lorsqu'elle vise des jeunes, public vulnérable. La présence d'une star et de la musique est licite, mais l'exagération trompeuse ne l'est pas.
\end{enumerate}
\textbf{Points de vigilance} : la question 3 exige une justification argumentée, pas un simple oui ou non ; il faut mobiliser le principe déontologique d'honnêteté (pas de publicité mensongère) et mentionner la vulnérabilité du public ciblé.
\end{corrige}

\begin{corrige}[Exercice 7 — Deux sources contradictoires]
\begin{enumerate}
\item \textbf{Source la plus fiable : l'article de journal.} Justification par trois critères :
\begin{itemize}
\item \textbf{Auteur identifiable} : l'article est signé, on peut donc identifier et interroger le journaliste ; le post Facebook est anonyme.
\item \textbf{Sources citées} : l'article cite le maire et des chiffres officiels, vérifiables ; le post ne cite aucune source ni aucune preuve.
\item \textbf{Recoupement possible} : les informations de l'article (date, déclarations, chiffres) peuvent être confrontées à d'autres médias et aux documents officiels ; rien ne permet de vérifier le post anonyme.
\end{itemize}
\item \textbf{Démarche de vérification complète} :
\begin{enumerate}
\item \textbf{Identifier la source} : qui publie ? L'auteur est-il nommé, connu, compétent ?
\item \textbf{Vérifier la date et le contexte} : l'information est-elle récente et à jour ?
\item \textbf{Recouper} : comparer avec au moins deux autres sources sérieuses et indépendantes (médias reconnus, communiqués officiels).
\item \textbf{Contrôler les preuves} : chiffres, citations, documents officiels sont-ils fournis et vérifiables ?
\item \textbf{Décider} : ne croire et ne partager que si les étapes précédentes confirment l'information ; dans le doute, s'abstenir de relayer.
\end{enumerate}
\end{enumerate}
\textbf{Points de vigilance} : les trois critères doivent être nommés explicitement et appliqués aux deux sources ; la démarche doit être présentée sous forme d'étapes ordonnées, comme exigé par le savoir-faire « rédiger et justifier une démarche ».
\end{corrige}

\begin{corrige}[Exercice 8 — La protection de l'opinion publique]
\begin{enumerate}
\item \textbf{Deux institutions ou textes protecteurs} :
\begin{itemize}
\item \textbf{La Constitution du 18 janvier 1996} : son préambule garantit la liberté d'expression, la liberté de la presse et le droit à l'information, fondements d'une opinion publique libre et éclairée.
\item \textbf{Le Conseil National de la Communication (CNC)} : organe de régulation qui veille au respect de la déontologie, à l'équité de traitement de l'information et peut adresser des mises en garde ou proposer des sanctions contre les médias fautifs.
\item On peut ajouter la \textbf{loi n° 90/052 du 19 décembre 1990} sur la communication sociale, qui encadre l'exercice de la profession de journaliste, et le \textbf{code pénal}, qui punit la diffusion de fausses nouvelles.
\end{itemize}
\item \textbf{La déontologie protège les citoyens} : la vérification des sources empêche la diffusion d'informations fausses qui tromperaient l'opinion ; le respect de la vie privée protège l'honneur et la dignité des personnes ; le \cle{droit de réponse} permet à toute personne mise en cause de rétablir la vérité. Ensemble, ces règles garantissent une information honnête, condition d'une opinion publique libre.
\item \textbf{Exemple de sanction} : le CNC peut prononcer une mise en garde, un blâme ou proposer la suspension temporaire d'un programme ou d'un organe de presse ; devant les tribunaux, la diffusion de fausses nouvelles peut être punie d'amendes et de peines de prison prévues par le code pénal.
\end{enumerate}
\textbf{Points de vigilance} : citer précisément les textes (Constitution de 1996, loi de 1990) ; ne pas attribuer au CNC un pouvoir de censure préalable : il intervient après diffusion, dans le respect de la liberté de la presse.
\end{corrige}

\begin{corrige}[Exercice 9 — Exercice documentaire type Probatoire]
\begin{enumerate}
\item \textbf{Présentation des documents} :
\begin{itemize}
\item \textbf{Document 1} : article de presse relatant une déclaration officielle ; ses sources sont identifiables (le ministre, le préfet) et l'information est confirmée par une publication au bulletin officiel de la région. Degré de fiabilité : \textbf{élevé}.
\item \textbf{Document 2} : post Facebook anonyme, sans source citée, sans preuve, avec un appel au partage massif et un ton accusateur. Degré de fiabilité : \textbf{très faible}.
\end{itemize}
\item \textbf{Information fiable : le Document 1.} Trois critères : \textbf{auteur identifiable} (journal et sources nommées contre anonymat) ; \textbf{sources citées et vérifiables} (ministre, préfet, bulletin officiel contre aucune source) ; \textbf{recoupement possible} (confirmation officielle publiée, contre une affirmation isolée et invérifiable).
\item \textbf{Le Document 2 est une \cle{fake news}} car il diffuse une information fausse (le lycée est inauguré et doté de moyens annoncés), sans auteur ni preuve, avec l'intention de tromper (« Partagez massivement ! »). Dangers : il sème la défiance envers les institutions, manipule l'opinion publique sur la base d'un mensonge et peut décourager les inscriptions au lycée, nuisant ainsi à l'intérêt général.
\item \textbf{Démarche de vérification en quatre étapes} :
\begin{enumerate}
\item Identifier l'auteur et la source du message ;
\item Vérifier la date et le contexte de l'information ;
\item Recouper avec au moins deux sources sérieuses (médias reconnus, communiqués officiels) ;
\item Ne partager que si l'information est confirmée ; dans le doute, s'abstenir.
\end{enumerate}
\item \textbf{Institution compétente} : le \textbf{Conseil National de la Communication (CNC)}. Son rôle : réguler la communication sociale, veiller au respect de la déontologie et de l'éthique, et sanctionner les dérives (mises en garde, blâmes, propositions de suspension). Les tribunaux peuvent en outre punir pénalement la diffusion de fausses nouvelles.
\end{enumerate}
\textbf{Barème indicatif (sur 20)} : présentation des documents 4 pts ; critères de fiabilité 4 pts ; analyse de la fake news et dangers 4 pts ; démarche en quatre étapes 4 pts ; institution et rôle 4 pts.

\textbf{Points de vigilance} : présenter un document, c'est toujours indiquer sa nature, sa source et sa fiabilité ; les critères doivent être explicitement appliqués aux deux documents ; la démarche doit comporter exactement quatre étapes numérotées ; ne pas oublier que la sanction pénale relève des tribunaux, pas du CNC.
\end{corrige}

\begin{corrige}[Exercice 10 — Étude de cas : la rumeur du marché]
\begin{enumerate}
\item \textbf{Définitions et lien avec le cas} : une \cle{rumeur} est une information non vérifiée, souvent fausse, qui circule sans source certaine ; l'\cle{opinion publique} est l'ensemble des idées et jugements partagés par une population sur un sujet. Ici, la rumeur sur la viande avariée a fait naître dans la population la conviction que le marché était dangereux : l'opinion publique s'est formée sur une information fausse, ce qui a provoqué la désertion du marché.
\item \textbf{Mécanisme de propagation} : WhatsApp permet un transfert instantané de groupe en groupe (vitesse) ; la peur pour la santé suscite une émotion forte qui pousse à partager « pour protéger » ses proches ; personne ne vérifie auprès des services sanitaires avant de relayer. La rumeur s'auto-amplifie : plus elle circule, plus elle paraît vraie.
\item \textbf{Responsabilités engagées} : l'\textbf{auteur} du message porte la responsabilité principale (diffusion de fausses nouvelles, punie par le code pénal) ; les \textbf{relais} sont co-responsables, car partager sans vérifier contribue à la désinformation ; les \textbf{consommateurs} ont réagi sans esprit critique, transformant une rumeur en crise économique pour des commerçants innocents.
\item \textbf{Mesures de protection au Cameroun} : la \textbf{loi sur la communication sociale} (loi n° 90/052 du 19 décembre 1990) encadre la diffusion de l'information ; le \textbf{CNC} veille au respect de la déontologie et peut sanctionner les médias fautifs ; le \textbf{code pénal} punit la diffusion de fausses nouvelles (amendes, peines de prison).
\item \textbf{Trois actions concrètes} : vérifier toute alerte sanitaire auprès des autorités compétentes (services d'hygiène, mairie) avant d'y croire ; sensibiliser le quartier à l'éducation aux médias (causeries, affiches sur la vérification des informations) ; signaler les messages mensongers aux administrateurs des groupes et aux autorités au lieu de les relayer.
\end{enumerate}
\textbf{Barème indicatif (sur 20)} : définitions et application au cas 4 pts ; mécanisme de propagation 4 pts ; responsabilités 4 pts ; mesures de protection 4 pts ; actions concrètes 4 pts.

\textbf{Points de vigilance} : les définitions doivent être données avant l'application au cas ; bien distinguer les trois niveaux de responsabilité ; les mesures doivent être précises (nom de la loi, rôle du CNC) et non vagues ; les actions proposées doivent être réalisables localement.
\end{corrige}

\begin{corrige}[Exercice 11 — Question de synthèse type Probatoire]
\textbf{Corrigé rédigé (modèle de copie).}

\emph{Introduction.} Au Cameroun comme partout dans le monde, les médias occupent une place croissante dans la vie quotidienne : radio, télévision, presse écrite et réseaux sociaux informent chaque jour des millions de citoyens. Un \cle{média} est un moyen de diffusion de l'information vers un large public ; l'\cle{opinion publique} désigne l'ensemble des idées et jugements partagés par une population ; la \cle{démocratie} est un régime où le pouvoir appartient au peuple, qui l'exerce notamment par des choix éclairés. Or une information de qualité est la condition de ces choix. Dès lors, on peut se demander : les médias sont-ils un danger ou une chance pour la démocratie camerounaise ? Nous verrons d'abord qu'ils peuvent menacer la vie démocratique, puis qu'ils constituent aussi un instrument précieux au service des citoyens.

\emph{Première partie : les dangers des médias.} Les médias peuvent d'abord véhiculer la \cle{désinformation} : les \emph{fake news} circulent rapidement, surtout sur les réseaux sociaux, et trompent les citoyens. Ainsi, une rumeur WhatsApp accusant faussement des commerçants d'un marché de Douala a provoqué la fermeture de boutiques avant que l'enquête sanitaire ne rétablisse la vérité. Ensuite, les médias peuvent servir la \cle{manipulation de l'opinion} : un média aux mains d'intérêts particuliers peut orienter les esprits, imposer un point de vue unique et fausser le débat public. Enfin, la course à l'audience favorise parfois le sensationnalisme au détriment de la vérification, ce qui fragilise la confiance des citoyens dans toutes les informations, même vraies.

\emph{Deuxième partie : les apports des médias.} Les médias sont d'abord une source d'\cle{information} indispensable : grâce à la radio et à la télévision, comme la CRTV ou les radios privées, le citoyen camerounais connaît les décisions publiques, les campagnes de vaccination ou les dates des élections. Ils assurent ensuite une \cle{éducation civique} : les émissions de débat et les journaux expliquent les droits et devoirs du citoyen. Enfin, ils jouent un rôle de \cle{contre-pouvoir} : en enquêtant et en dénonçant les abus (mauvaise gestion, corruption), la presse oblige les responsables à rendre compte, ce qui est au cœur de la démocratie. La liberté de la presse, garantie par le préambule de la Constitution du 18 janvier 1996, protège précisément cette mission.

\emph{Conclusion.} Les médias sont donc à la fois un danger potentiel — désinformation, manipulation, rumeurs — et une chance réelle — information, éducation, contre-pouvoir — pour la démocratie camerounaise. Tout dépend de l'usage qui en est fait. C'est pourquoi la régulation par le \cle{CNC} et les lois (loi de 1990 sur la communication sociale, code pénal contre les fausses nouvelles), jointe à l'\cle{éducation aux médias} des citoyens dès l'école, permet de limiter les dangers et de faire des médias un véritable instrument de démocratie.

\textbf{Barème indicatif (sur 20)} : introduction (mise en contexte, définitions, problématique, annonce du plan) 4 pts ; première partie (deux à trois dangers illustrés) 4 pts ; deuxième partie (deux à trois apports illustrés) 4 pts ; conclusion équilibrée avec ouverture sur la régulation 4 pts ; cohérence, orthographe et exemples camerounais 4 pts.

\textbf{Points de vigilance} : la problématique doit être formulée par une phrase interrogative ; chaque argument doit être illustré d'un exemple camerounais ; éviter la conclusion unilatérale (« les médias sont un danger ») : le sujet appelle une réponse nuancée ; citer au moins un texte (Constitution de 1996, loi de 1990) et une institution (CNC) ; respecter le plan dialectique en deux parties annoncé dans l'introduction.
\end{corrige}

\newpage

\coursec{Fiche bilan}

\begin{popfigure}
\begin{tikzpicture}[
  mindmap,
  concept color=popblue,
  level 1/.append style={level distance=4.5cm, sibling angle=60},
  level 2/.append style={level distance=3cm, sibling angle=45},
  every node/.style={font=\small, text=popdark},
  branch1/.style={concept color=poporange, grow=30},
  branch2/.style={concept color=popgreen, grow=90},
  branch3/.style={concept color=poppurple, grow=150},
  branch4/.style={concept color=poppink, grow=210},
  branch5/.style={concept color=popgold, grow=270},
  branch6/.style={concept color=popblue, grow=330},
  text width=2.8cm, align=center
]
\node[concept, text width=3.2cm] {Médias \& Opinion Publique}
  child[branch1] { node[concept, align=center]{1. Médias\\Définitions \& Typologie} }
  child[branch2] { node[concept, align=center]{2. Opinion Publique\\Formation \& Caractères} }
  child[branch3] { node[concept, align=center]{3. Impact Médias\\Fonctions \& Effets} }
  child[branch4] { node[concept, align=center]{4. Protection Cameroun\\Cadre Légal \& CNC} }
  child[branch5] { node[concept, align=center]{5. Publicité (TD3)\\Rôle, Régulation, Éthique} }
  child[branch6] { node[concept, align=center]{6. Enjeux Citoyens\\Esprit Critique \& Fake News} };
\end{tikzpicture}
\legende{Carte mentale de synthèse : Les médias et leur impact sur l'opinion publique}
\end{popfigure}

\begin{bilan}
\courssub{Définitions Clés}
\begin{itemize}
  \item \cle{Médias} : Ensemble des moyens techniques de communication de masse (presse, radio, TV, internet) permettant de diffuser une information vers un large public.
  \item \cle{Opinion publique} : Ensemble des attitudes, croyances et jugements partagés par une collectivité sur une question d'intérêt général à un moment donné.
  \item \cle{Publicité} : Communication commerciale payante, non personnelle, visant à promouvoir un produit, service ou idée auprès d'une cible.
  \item \cle{Agenda-setting} : Capacité des médias à imposer les thèmes dont le public débat (fixation de l'agenda).
  \item \cle{CNC} : Conseil National de la Communication, organe de régulation du secteur de la communication sociale au Cameroun.
\end{itemize}

\courssub{Repères Légaux \& Institutionnels (Cameroun)}
\begin{center}
\begin{tabularx}{\linewidth}{|l|X|}\hline
\rowcolor{popblueL} \textbf{Texte / Organe} & \textbf{Rôle Principal} \\ \hline
Constitution (Préambule) & Garantit la liberté de presse, d'expression et de communication. \\ \hline
Loi n°90/052 (19 déc. 1990) & Régime de la communication de masse : libertés, obligations, sanctions. \\ \hline
Loi n°2016/015 (14 déc. 2016) & Création du CNC, missions, composition, fonctionnement. \\ \hline
CNC & Régulation, déontologie, attribution des cartes de presse, médiation, sanctions administratives. \\ \hline
Code Pénal (Art. 113, 305...) & Répression : diffamation, injure, fausses nouvelles, outrage aux corps constitués. \\ \hline
\end{tabularx}
\end{center}

\courssub{Méthodes Express}
\begin{methode}[Analyser un message médiatique]
  \begin{enumerate}
    \item \textbf{Source} : Qui émet ? (Média, journaliste, influenceur, institution).
    \item \textbf{Angle / Cadrage (Framing)} : Quel aspect est mis en avant ? Quel vocabulaire ?
    \item \textbf{Cible} : À qui s'adresse le message ? (Langage, support, horaire).
    \item \textbf{Intention} : Informer, persuader, divertir, manipuler, vendre ?
    \item \textbf{Vérification} : Croiser les sources, vérifier la date, l'auteur, les images (outils fact-checking).
  \end{enumerate}
\end{methode}

\begin{methode}[Analyser une publicité (TD3)]
  \begin{enumerate}
    \item Identifier le \textbf{produit/marque} et la \textbf{cible} (âge, sexe, CSP, centres d'intérêt).
    \item Repérer la \textbf{promesse} (bénéfice promis) et la \textbf{preuve} (argument rationnel/émotionnel).
    \item Analyser la \textbf{création} : slogan, visuel, musique, storytelling, humour, peur, témoignage.
    \item Vérifier la \textbf{légalité/éthique} : respect Code de la publicité (CNCC), non-discrimination, véracité, protection mineurs.
  \end{enumerate}
\end{methode}
\end{bilan}

\begin{aretenir}[Checklist avant le Bac]
  $\square$ Définir les médias et citer 3 types (presse, audiovisuel, numérique) avec exemples locaux.\par
  $\square$ Définir l'opinion publique et expliquer ses 3 caractéristiques (collectivité, sujet d'intérêt général, variabilité).\par
  $\square$ Citer les 4 fonctions classiques des médias (Laswell) : surveillance, corrélation, transmission culturelle, divertissement.\par
  $\square$ Expliquer l'\emph{agenda-setting} et donner un exemple concret (ex: traitement d'une grève, d'une élection).\par
  $\square$ Citer la base légale de la liberté de presse au Cameroun (Préambule Constitution, Loi 90/052).\par
  $\square$ Lister 3 missions du CNC (régulation, déontologie, carte de presse, médiation, sanction).\par
  $\square$ Différencier \emph{information}, \emph{communication}, \emph{propagande} et \emph{publicité}.\par
  $\square$ Analyser un spot publicitaire : cible, promesse, preuve, technique de persuasion, conformité éthique.\par
  $\square$ Identifier 3 infractions de presse sanctionnées par le Code Pénal (diffamation, fausses nouvelles, injure...).\par
  $\square$ Appliquer la démarche de vérification d'une info (source, date, auteur, recoupement) face à une \emph{fake news}.\par
  $\square$ Argumenter sur le rôle du citoyen : esprit critique, consommation responsable, signalement (CNC, plateformes).
\end{aretenir}

\findecours

\end{document}
